% Œil et la vision — SVTEEHB 3ème — Cours signé OnBuch+
% Programme officiel MINESEC. Compiler avec : tectonic N10-il-vision.tex
\newcommand{\DOCMATIERE}{SVTEEHB}
\newcommand{\DOCNIVEAU}{3ème}
\newcommand{\DOCMODULE}{SVTEEHB – classe de 3ème}
\newcommand{\DOCLECON}{N10}
\newcommand{\DOCTITRE}{Œil et la vision}
% OnBuch+ — préambule « cours signés » ENRICHI (compilé avec tectonic / XeLaTeX)
% Système visuel « Pop » d'OnBuch+ : fond crème (jamais blanc), bordures encre
% épaisses, ombres « dures » décalées, titres Archivo Black en capitales.
% Version MATHÉMATIQUES : graphes (pgfplots/tikz), boîtes pédagogiques
% (définition, propriété, méthode, attention, le savais-tu, exercice résolu,
% exercice, TP, bilan), page de garde et signature OnBuch+ ancrée en bas.
%
% Macros attendues dans main.tex AVANT \input{preamble} :
%   \DOCMATIERE  \DOCNIVEAU  \DOCMODULE  \DOCLECON (numéro)  \DOCTITRE
\documentclass[11pt]{article}

\usepackage{fontspec}
\usepackage[french]{babel}
\usepackage[a4paper,margin=2cm,top=3.4cm,bottom=2.8cm,headheight=1.4cm,footskip=1.3cm]{geometry}
\usepackage{amsmath,amssymb,amsfonts}
\usepackage{mathtools} % \xmapsto, \coloneqq... souvent utilisés par les modèles
\usepackage{cancel} % simplifications barrées (bilans de forces)
% \abs{z} (module, valeur absolue) et \norm{u} (norme), utilisés par les modèles
\providecommand{\abs}{}\let\abs\relax\DeclarePairedDelimiter{\abs}{\lvert}{\rvert}
\providecommand{\norm}{}\let\norm\relax\DeclarePairedDelimiter{\norm}{\lVert}{\rVert}
\usepackage[table]{xcolor} % [table] : fournit aussi \rowcolors (alternance de lignes)
\usepackage{pagecolor}
\usepackage{tikz}
\usetikzlibrary{arrows.meta,positioning,calc,shapes.geometric,shapes.misc,shapes.arrows,decorations.pathmorphing,decorations.pathreplacing,decorations.markings,patterns,shadows,mindmap,trees,fit,backgrounds,matrix,intersections,angles,quotes}
\usepackage{pgfplots}
\usepackage[version=4]{mhchem} % équations nucléaires \ce{^{235}_{92}U}
\usepackage[european,siunitx]{circuitikz} % schémas électriques
\pgfplotsset{compat=1.18}
\usepgfplotslibrary{fillbetween,groupplots} % aires entre courbes (calcul intégral)
\usepackage{enumitem}
\usepackage{fancyhdr}
% [most] charge toutes les bibliothèques tcolorbox (listings, minted, raster,
% documentation...) et gonfle inutilement la mémoire TeX sur un document long
% (~300 boîtes) : on ne charge que ce qui est réellement utilisé.
\usepackage[skins,breakable]{tcolorbox}
\usepackage{graphicx}
\usepackage{colortbl}
\usepackage{booktabs}
\usepackage{tabularx}
\usepackage{multirow}
\usepackage{diagbox} % case d'en-tête diagonale des tableaux à double entrée
% Colonnes extensibles centrée / gauche / droite (C, L, R), souvent utilisées par les modèles
\newcolumntype{C}{>{\centering\arraybackslash}X}
\newcolumntype{L}{>{\raggedright\arraybackslash}X}
\newcolumntype{R}{>{\raggedleft\arraybackslash}X}
\usepackage{array}
\usepackage{multicol}
\usepackage{siunitx}
% parse-numbers=false : accepte aussi \SI{2,5 \times 10^{-3}}{...}
\sisetup{locale=FR,output-decimal-marker={,},group-minimum-digits=4,parse-numbers=false}
\DeclareSIUnit\ohm{\ensuremath{\symup{\Omega}}} % U+2126 absent de la police
\DeclareSIUnit\division{div} % réglages d'oscilloscope (ms/div, V/div)
\DeclareSIUnit\tr{tr} % tours (vitesse de rotation, ex. \SI{1500}{\tr\per\minute})
\DeclareSIUnit\molar{M} % concentration molaire (ex. \SI{3}{\molar}, \SI{1}{\micro\molar})
\DeclareSIUnit\an{an} % année (le modèle confond parfois avec \year, un compteur TeX) — ex. \SI{5}{\kilo\meter\per\an}
\DeclareSIUnit\FCFA{FCFA} % franc CFA (ex. \SI{500}{\FCFA\per\kilo\gram})
\DeclareSIUnit\degreeBrix{\textdegree Brix} % teneur en sucre d'un sirop/jus (réfractométrie)
\DeclareSIUnit\month{mois} % le modèle utilise parfois \month, un compteur TeX, comme unité de durée
\DeclareSIUnit\microliter{\micro\liter} % le modèle écrit parfois \microliter en un seul token
\DeclareSIUnit\million{millions} % dénombrement (ex. \SI{15}{\million\per\mL} spermatozoïdes)
\usepackage{qrcode}
% \degree : commande fréquemment utilisée par les modèles pour un angle ou une
% température en dehors de \SI{}{} (non fournie par les paquets chargés).
\providecommand{\degree}{^{\circ}}
% \qed (fin de démonstration), utilisé par les modèles sans amsthm
\providecommand{\qed}{\hfill$\square$}
% \barre : double barre des valeurs interdites dans un tableau de variations (tkz-tab)
\providecommand{\barre}{\ensuremath{\|}}
% environnement proof (amsthm non chargé) utilisé par les modèles
\ifdefined\proof\else\newenvironment{proof}[1][Démonstration]{\par\noindent\textit{#1.}\ }{\qed\par}\fi
\usepackage{lastpage}
\usepackage{hyperref}
\hypersetup{hidelinks}

% ============================================================
% Palette « Pop » OnBuch+ — mêmes hex que la classe P (plus_ui.dart)
% ============================================================
\definecolor{poporange}{HTML}{F59321}
\definecolor{popdark}{HTML}{E07A0C}
\definecolor{popcream}{HTML}{FFF6E8}
\definecolor{popink}{HTML}{1C1714}
\definecolor{poppurple}{HTML}{7A5AE0}
\definecolor{popgreen}{HTML}{2E9E6A}
\definecolor{poppink}{HTML}{E0457B}
\definecolor{popblue}{HTML}{2F7DE1}
\definecolor{popgold}{HTML}{C98A12}
\definecolor{popmuted}{HTML}{8A7F76}
\definecolor{popline}{HTML}{EADFD2}
\definecolor{popgray}{HTML}{8A7F76}
\colorlet{popgrey}{popgray}
% Alias tolérants (le modèle nomme parfois les couleurs différemment)
\colorlet{popwhite}{white}
\colorlet{popblack}{popink}
\colorlet{popred}{poppink}
\colorlet{popyellow}{popgold}
% Teintes claires (fonds de tableaux/encadrés) — le modèle nomme parfois
% les couleurs avec un suffixe L (ex. popblueL) sans les avoir définies.
\colorlet{poporangeL}{poporange!20}
\colorlet{popdarkL}{popdark!20}
\colorlet{poppurpleL}{poppurple!20}
\colorlet{popgreenL}{popgreen!20}
\colorlet{poppinkL}{poppink!20}
\colorlet{popblueL}{popblue!20}
\colorlet{popgoldL}{popgold!20}
\colorlet{popredL}{popred!20}
\colorlet{popgrayL}{popgray!20}
% Teintes claires pour fonds de tableaux / zones
\colorlet{popblueL}{popblue!10!white}
\colorlet{poporangeL}{poporange!14!white}
\colorlet{popgreenL}{popgreen!12!white}
\colorlet{poppurpleL}{poppurple!12!white}
\colorlet{poppinkL}{poppink!12!white}
\colorlet{popgoldL}{popgold!14!white}

\pagecolor{popcream}

% ============================================================
% Typographie
% ============================================================
\defaultfontfeatures{Ligatures=TeX}
\setmainfont{PlusJakartaSans-Medium.ttf}[Path=../../../fonts/,BoldFont=PlusJakartaSans-Bold.ttf]
% Maths en sans-serif (Fira Math) avec chiffres et lettres droites de Plus
% Jakarta, pour que les formules se fondent dans le texte.
\usepackage[math-style=ISO,bold-style=ISO]{unicode-math}
\setmathfont{FiraMath-Regular.otf}[Path=../../../fonts/]
\setmathfont{PlusJakartaSans-Medium.ttf}[Path=../../../fonts/,range={up/num,up/{Latin,latin},\mathup}]
\usepackage{icomma} % virgule décimale sans espace en mode math
% Caractères mathématiques Unicode absents des polices de texte (Plus Jakarta,
% Archivo Black) : rendus par leur équivalent mathématique, en texte comme en
% formule — sinon ils s'affichent en carré vide (ex. « Arithmétique dans ℤ »).
\usepackage{newunicodechar}
\newunicodechar{ℝ}{\ensuremath{\mathbb{R}}}
\newunicodechar{ℤ}{\ensuremath{\mathbb{Z}}}
\newunicodechar{ℂ}{\ensuremath{\mathbb{C}}}
\newunicodechar{ℕ}{\ensuremath{\mathbb{N}}}
\newunicodechar{ℚ}{\ensuremath{\mathbb{Q}}}
\newunicodechar{𝔻}{\ensuremath{\mathbb{D}}}
\newunicodechar{α}{\ensuremath{\alpha}}
\newunicodechar{β}{\ensuremath{\beta}}
\newunicodechar{γ}{\ensuremath{\gamma}}
\newunicodechar{δ}{\ensuremath{\delta}}
\newunicodechar{ε}{\ensuremath{\varepsilon}}
\newunicodechar{θ}{\ensuremath{\theta}}
\newunicodechar{λ}{\ensuremath{\lambda}}
\newunicodechar{μ}{\ensuremath{\mu}}
\newunicodechar{σ}{\ensuremath{\sigma}}
\newunicodechar{τ}{\ensuremath{\tau}}
\newunicodechar{φ}{\ensuremath{\varphi}}
\newunicodechar{ψ}{\ensuremath{\psi}}
\newunicodechar{ω}{\ensuremath{\omega}}
\newunicodechar{Σ}{\ensuremath{\Sigma}}
% majuscules produites par \MakeUppercase dans les titres (α-aminés -> Α)
\newunicodechar{Α}{\ensuremath{\alpha}}
\newunicodechar{Β}{\ensuremath{\beta}}
\newunicodechar{Γ}{\ensuremath{\Gamma}}
\newunicodechar{Δ}{\ensuremath{\Delta}}
\newunicodechar{Ε}{\ensuremath{\varepsilon}}
\newunicodechar{Θ}{\ensuremath{\Theta}}
\newunicodechar{Λ}{\ensuremath{\Lambda}}
\newunicodechar{Μ}{\ensuremath{\mu}}
\newunicodechar{Τ}{\ensuremath{\tau}}
\newunicodechar{Φ}{\ensuremath{\Phi}}
\newunicodechar{Ψ}{\ensuremath{\Psi}}
\newunicodechar{Ω}{\ensuremath{\Omega}}
\newunicodechar{↦}{\ensuremath{\mapsto}}
\newunicodechar{∈}{\ensuremath{\in}}
\newunicodechar{∉}{\ensuremath{\notin}}
\newunicodechar{∧}{\ensuremath{\wedge}}
\newunicodechar{⇔}{\ensuremath{\Leftrightarrow}}
\newunicodechar{⟺}{\ensuremath{\Longleftrightarrow}}
\newunicodechar{⇒}{\ensuremath{\Rightarrow}}
\newunicodechar{⟹}{\ensuremath{\Longrightarrow}}
\newunicodechar{∘}{\ensuremath{\circ}}
\newunicodechar{∪}{\ensuremath{\cup}}
\newunicodechar{∩}{\ensuremath{\cap}}
\newunicodechar{≡}{\ensuremath{\equiv}}
\newunicodechar{⊂}{\ensuremath{\subset}}
\newunicodechar{⊥}{\ensuremath{\perp}}
\newunicodechar{∃}{\ensuremath{\exists}}
\newunicodechar{∀}{\ensuremath{\forall}}
\newunicodechar{∛}{\ensuremath{\sqrt[3]{\,}}}
\newunicodechar{∜}{\ensuremath{\sqrt[4]{\,}}}
\newunicodechar{⃗}{\ensuremath{{}^{\to}}}
\newunicodechar{ⁿ}{\ifmmode^{n}\else\textsuperscript{\ensuremath{n}}\fi}
\newunicodechar{ˣ}{\ifmmode^{x}\else\textsuperscript{\ensuremath{x}}\fi}
\newunicodechar{ᵇ}{\ifmmode^{b}\else\textsuperscript{\ensuremath{b}}\fi}
\newunicodechar{ʳ}{\ifmmode^{r}\else\textsuperscript{\ensuremath{r}}\fi}
\newunicodechar{ᵘ}{\ifmmode^{u}\else\textsuperscript{\ensuremath{u}}\fi}
\newunicodechar{ʸ}{\ifmmode^{y}\else\textsuperscript{\ensuremath{y}}\fi}
\newunicodechar{ᵃ}{\ifmmode^{a}\else\textsuperscript{\ensuremath{a}}\fi}
\newunicodechar{ᵉ}{\ifmmode^{e}\else\textsuperscript{\ensuremath{e}}\fi}
\newunicodechar{ᵏ}{\ifmmode^{k}\else\textsuperscript{\ensuremath{k}}\fi}
\newunicodechar{ᵖ}{\ifmmode^{p}\else\textsuperscript{\ensuremath{p}}\fi}
\newunicodechar{ᴺ}{\ifmmode^{N}\else\textsuperscript{\ensuremath{N}}\fi}
\newunicodechar{ᶜ}{\ifmmode^{c}\else\textsuperscript{\ensuremath{c}}\fi}
\newunicodechar{ʰ}{\ifmmode^{h}\else\textsuperscript{\ensuremath{h}}\fi}
\newunicodechar{ᵠ}{\ifmmode^{q}\else\textsuperscript{\ensuremath{q}}\fi}
\newunicodechar{ₙ}{\ifmmode_{n}\else\textsubscript{\ensuremath{n}}\fi}
\newunicodechar{ₐ}{\ifmmode_{a}\else\textsubscript{\ensuremath{a}}\fi}
\newunicodechar{ᵢ}{\ifmmode_{i}\else\textsubscript{\ensuremath{i}}\fi}
\newunicodechar{ᵣ}{\ifmmode_{r}\else\textsubscript{\ensuremath{r}}\fi}
\newunicodechar{ₖ}{\ifmmode_{k}\else\textsubscript{\ensuremath{k}}\fi}
\newunicodechar{ᵦ}{\ifmmode_{\beta}\else\textsubscript{\ensuremath{\beta}}\fi}
\newunicodechar{ₓ}{\ifmmode_{x}\else\textsubscript{\ensuremath{x}}\fi}
\newfontfamily\popdisplay{ArchivoBlack-Regular.ttf}[Path=../../../fonts/]
\newfontfamily\poplogo{SpaceGrotesk-Bold.ttf}[Path=../../../fonts/]
\newfontfamily\popsemi{PlusJakartaSans-SemiBold.ttf}[Path=../../../fonts/]
\newfontfamily\popxbold{PlusJakartaSans-ExtraBold.ttf}[Path=../../../fonts/]
\newfontfamily\popxboldls{PlusJakartaSans-ExtraBold.ttf}[Path=../../../fonts/,LetterSpace=3.0]

\setlist[enumerate]{leftmargin=1.7em,itemsep=5pt,topsep=4pt}
\setlist[itemize]{leftmargin=1.7em,itemsep=4pt,topsep=4pt,label={\color{poporange}\textbullet}}
\setlength{\parindent}{0pt}
\setlength{\parskip}{5pt}
\renewcommand{\arraystretch}{1.25}

% pgfplots : style Pop par défaut
\pgfplotsset{
  every axis/.append style={axis line style={popink,line width=1pt},tick style={popink},
    grid style={popline,line width=0.5pt},label style={font=\small},tick label style={font=\footnotesize}},
  popcurve/.style={poporange,line width=1.8pt,no markers,smooth},
}
% Même style disponible dans un \draw TikZ simple (hors pgfplots)
\tikzset{popcurve/.style={poporange,line width=1.8pt}}
\tikzset{popgrid/.style={popline,very thin}}
\colorlet{popcurve}{poporange} % le nom du style est parfois utilisé comme couleur

% ============================================================
% Pill / squiggle
% ============================================================
\newcommand{\poppill}[3]{%
  \tikz[baseline=-0.55ex]\node[draw=popink,line width=1.2pt,rounded corners=2.6pt,%
    fill=#2,inner xsep=6.5pt,inner ysep=3pt]{%
    \popxboldls\fontsize{7}{7}\selectfont\color{#3}\MakeUppercase{#1}};%
}
\newcommand{\popsquiggle}[1]{%
  \tikz[baseline=-0.4ex]{
    \draw[#1,line width=2.6pt,line cap=round]
      plot [smooth] coordinates {(0,0.09) (0.34,0.01) (0.68,0.21) (1.02,0.01) (1.36,0.10)};
  }%
}
% Mot-clé surligné (marqueur orange sous le texte)
\newcommand{\cle}[1]{{\popxbold\color{popdark}#1}}

% ============================================================
% En-tête / pied
% ============================================================
\pagestyle{fancy}
\fancyhf{}
\renewcommand{\headrulewidth}{0pt}
\renewcommand{\footrulewidth}{0pt}
\lhead{\raisebox{-0.15ex}{\poplogo\fontsize{13}{13}\selectfont\color{popink}OnBuch\color{poporange}+}}
\rhead{\poppill{\DOCMATIERE\ \textbullet\ \DOCNIVEAU\ \textbullet\ Leçon \DOCLECON}{white}{popink}}
\lfoot{\popxbold\fontsize{7.6}{7.6}\selectfont\color{popmuted}\MakeUppercase{Cours signé OnBuch+}\ \textbullet\ {\popsemi\DOCTITRE}}
\rfoot{\tikz[baseline=-0.6ex]\node[rounded corners=8.5pt,fill=popink,minimum height=17pt,minimum width=17pt,inner xsep=5pt,inner sep=0]{\popxbold\fontsize{8}{8}\selectfont\color{popcream}\thepage\,/\,\pageref*{LastPage}};}

% ============================================================
% Titres
% ============================================================
\newcommand{\chaptitre}[2]{%
  \begin{tcolorbox}[enhanced,colback=poporange,colframe=popink,boxrule=2.6pt,arc=11pt,
      drop shadow=popink,left=15pt,right=15pt,top=12pt,bottom=11pt,before skip=3pt,after skip=18pt]
    \poppill{#2}{popink}{popcream}\par\vspace{9pt}
    {\raggedright\hyphenpenalty=10000\exhyphenpenalty=10000\popdisplay\fontsize{22}{24}\selectfont\color{popcream}\MakeUppercase{#1}\par}
  \end{tcolorbox}
}
% Section numérotée : pastille orange avec le numéro + titre Archivo
\newcounter{popsec}
\newcommand{\coursec}[1]{%
  \stepcounter{popsec}\setcounter{popsub}{0}%
  \par\vspace{16pt}\noindent
  \begin{minipage}{\linewidth}
  \tikz[baseline=-0.7ex]\node[draw=popink,line width=1.6pt,fill=poporange,rounded corners=4pt,minimum size=22pt,inner sep=0]{\popdisplay\fontsize{12}{12}\selectfont\color{popcream}\thepopsec};%
  \hspace{8pt}{\popdisplay\fontsize{15}{17}\selectfont\color{popink}\MakeUppercase{#1}}\par
  \vspace{4pt}\noindent{\color{poporange}\rule{36pt}{3.6pt}}
  \end{minipage}\par\nopagebreak\vspace{8pt}
}
\newcounter{popsub}
\newcommand{\courssub}[1]{%
  \stepcounter{popsub}%
  \par\vspace{9pt}\noindent{\popxbold\fontsize{11.5}{13}\selectfont\color{popink}{\color{poporange}\thepopsec.\thepopsub}\hspace{6pt}#1}\par\nopagebreak\vspace{4pt}
}

% ============================================================
% Boîtes pédagogiques — même recette : fond blanc, bordure épaisse colorée,
% ombre dure encre, badge plein. Titre optionnel [#1].
% ============================================================
\tcbset{popbase/.style={breakable,enhanced,colback=white,boxrule=2.3pt,arc=9pt,
  shadow={3pt}{-3pt}{0pt}{popink},left=14pt,right=14pt,top=11pt,bottom=11pt,before skip=11pt,after skip=15pt,
  fontupper=\popsemi\fontsize{10.6}{14.5}\selectfont\color{popink},
  fontlower=\popsemi\fontsize{10.6}{14.5}\selectfont\color{popink}}}
% Coche dessinée (le glyphe ✓ n'existe pas dans les polices du document)
\newcommand{\popcheck}[1]{\tikz[baseline=-0.5ex]\draw[#1,line width=1.6pt,line cap=round,line join=round](-0.13,0.0)--(-0.04,-0.09)--(0.13,0.1);}
\providecommand{\coche}{\popcheck{popgreen}}
\providecommand{\resultat}[1]{\par\smallskip\noindent\fcolorbox{popgreen}{popgreenL}{\parbox{\dimexpr\linewidth-2\fboxsep-2\fboxrule\relax}{\textbf{Résultat.} #1}}\par\smallskip}
\newcommand{\popboxhead}[3]{\poppill{#1}{#2}{white}\ifx\relax#3\relax\else\hspace{6pt}{\popxbold\fontsize{10.6}{12}\selectfont\color{popink}#3}\fi\par\vspace{7pt}}

\newtcolorbox{aretenir}[1][]{popbase,colframe=popgreen,before upper={\popboxhead{À retenir}{popgreen}{#1}}}
\newtcolorbox{exemplebox}[1][]{popbase,colframe=poporange,before upper={\popboxhead{Exemple}{poporange}{#1}}}
\newtcolorbox{definition}[1][]{popbase,colframe=popblue,before upper={\popboxhead{Définition}{popblue}{#1}}}
\newtcolorbox{propriete}[1][]{popbase,colframe=poppurple,before upper={\popboxhead{Propriété}{poppurple}{#1}}}
\newtcolorbox{methode}[1][]{popbase,colframe=popgold,before upper={\popboxhead{Méthode}{popgold}{#1}}}
\newtcolorbox{attention}[1][]{popbase,colframe=poppink,before upper={\popboxhead{Attention}{poppink}{#1}}}
\newtcolorbox{experience}[1][]{popbase,colframe=popblue,colback=popblueL,before upper={\popboxhead{Expérience}{popblue}{#1}}}
% « Le savais-tu ? » : carte violette pleine (contexte, culture, Cameroun)
\newtcolorbox{savaistu}[1][]{popbase,colback=poppurple,colframe=popink,
  fontupper=\popsemi\fontsize{10.4}{14.2}\selectfont\color{white},
  before upper={\poppill{Le savais-tu\,?}{white}{poppurple}\ifx\relax#1\relax\else\hspace{6pt}{\popxbold\color{white}#1}\fi\par\vspace{7pt}}}
% Exercice résolu : énoncé en haut, \tcblower puis la solution sur fond crème
\newcounter{popexo}\newcounter{popexr}
\newtcolorbox{exoresolu}[1][]{popbase,colframe=poporange,colbacklower=popcream,
  segmentation style={popink,line width=1.2pt,dashed},
  before upper={\stepcounter{popexr}\popboxhead{Exercice résolu \thepopexr}{poporange}{#1}},
  before lower={\poppill{Solution}{popink}{popcream}\par\vspace{6pt}}}
% Exercice (non résolu en place) — la correction est dans la partie Corrigés
\newtcolorbox{exercice}[1][]{popbase,colframe=popink,
  before upper={\stepcounter{popexo}\popboxhead{Exercice \thepopexo}{popink}{#1}}}
% Corrigé
\newtcolorbox{corrige}[1][]{popbase,colframe=popgreen,colback=popgreenL,
  before upper={\popboxhead{Corrigé}{popgreen}{#1}}}
% Objectifs / prérequis / bilan
\newtcolorbox{objectifs}{popbase,colframe=popink,colback=poporangeL,
  before upper={\popboxhead{Objectifs de la leçon}{popink}{}}}
\newtcolorbox{prerequis}{popbase,colframe=popink,colback=popblueL,
  before upper={\popboxhead{Prérequis}{popblue}{}}}
\newtcolorbox{bilan}{popbase,colframe=popink,colback=white,boxrule=2.8pt,
  before upper={\popboxhead{Fiche bilan}{poporange}{L'essentiel à connaître pour l'examen}}}
% Figure encadrée avec légende
\newtcolorbox{popfigure}[1][]{enhanced,breakable=false,colback=white,colframe=popline,boxrule=1.4pt,arc=8pt,
  left=10pt,right=10pt,top=10pt,bottom=8pt,before skip=10pt,after skip=12pt,halign=center,
  lower separated=false,halign lower=center,
  fontlower=\popsemi\fontsize{9}{11.5}\selectfont\color{popmuted},
  #1}
\newcommand{\legende}[1]{\par\vspace{6pt}{\popsemi\fontsize{9}{11.5}\selectfont\color{popmuted}#1}}

% ============================================================
% Signature OnBuch+
% ============================================================
\newcommand{\poplogoinline}[1]{{\poplogo\fontsize{#1}{#1}\selectfont\color{popink}OnBuch\color{poporange}+}}
\newcommand{\signatureonbuch}{%
  \begin{tcolorbox}[enhanced,colback=popink,colframe=popink,boxrule=2.2pt,arc=10pt,
      drop shadow=poporange,left=14pt,right=14pt,top=10pt,bottom=10pt,before skip=0pt,after skip=0pt]
    \tikz[baseline=-0.7ex]\node[circle,fill=popgreen,draw=popcream,line width=1.3pt,minimum size=22pt,inner sep=0]{\popcheck{white}};%
    \hspace{9pt}\begin{minipage}[c]{0.62\linewidth}
      {\popxbold\fontsize{12}{13}\selectfont\color{popcream}Signé\ \textbullet\ {\poplogo OnBuch\color{poporange}+}}\par\vspace{2pt}
      {\raggedright\popsemi\fontsize{8}{10.5}\selectfont\color{popline}\MakeUppercase{Équipe pédagogique OnBuch+}\\\MakeUppercase{Conforme au programme officiel MINESEC}\par}
    \end{minipage}\hfill
    {\poplogo\fontsize{22}{22}\selectfont\color{popcream}OnBuch\color{poporange}+}
  \end{tcolorbox}%
}

% ============================================================
% Page de garde
% #1 titre · #2 matière · #3 niveau · #4 module · #5 n° leçon · #6 durée ·
% #7 description / objectifs (texte ou itemize)
% ============================================================
\newcommand{\pagegarde}[7]{%
  \thispagestyle{empty}
  \newgeometry{a4paper,margin=1.7cm,top=1.6cm,bottom=1.4cm}%
  \begin{tikzpicture}[remember picture,overlay]
    \fill[poporange!18] ([xshift=-3.2cm,yshift=-3.6cm]current page.north east) circle (4.6cm);
    \fill[poppurple!14] ([xshift=2.4cm,yshift=7.6cm]current page.south west) circle (3.8cm);
  \end{tikzpicture}%
  {\poplogo\fontsize{30}{30}\selectfont\color{popink}OnBuch\color{poporange}+}\hfill
  \raisebox{0.6ex}{\poppill{Cours signé \textbullet\ Premium}{popink}{popcream}}\par
  \vspace{1pt}\popsquiggle{poppurple}\par
  \vspace{8pt}
  \begin{tcolorbox}[enhanced,colback=poporange,colframe=popink,boxrule=2.8pt,arc=13pt,
      drop shadow=popink,left=18pt,right=18pt,top=12pt,bottom=13pt,after skip=10pt]
    \poppill{#2 \textbullet\ #3}{popink}{popcream}\hspace{5pt}\poppill{#4}{white}{popink}\par\vspace{8pt}
    {\popdisplay\fontsize{14}{14}\selectfont\color{popink}LEÇON #5}\par\vspace{4pt}
    {\raggedright\hyphenpenalty=10000\exhyphenpenalty=10000\popdisplay\fontsize{25}{27}\selectfont\color{popcream}\MakeUppercase{#1}\par}
    \vspace{8pt}
    {\popxbold\fontsize{9.5}{9.5}\selectfont\color{popink}\MakeUppercase{Durée conseillée : #6}}
  \end{tcolorbox}
  \begin{tcolorbox}[enhanced,colback=white,colframe=popink,boxrule=2.2pt,arc=10pt,
      drop shadow=popink,left=16pt,right=16pt,top=10pt,bottom=10pt,after skip=8pt]
    \poppill{Dans cette leçon}{poppurple}{white}\par\vspace{5pt}
    \popsemi\fontsize{9.3}{12.6}\selectfont\color{popink}#7
  \end{tcolorbox}
  \noindent\begin{minipage}[t]{0.487\linewidth}
    \begin{tcolorbox}[enhanced,colback=popgreen,colframe=popink,boxrule=2.2pt,arc=10pt,
        drop shadow=popink,left=11pt,right=11pt,top=10pt,bottom=10pt,height=3.1cm,valign=center]
      \tcbox[colback=white,colframe=popink,boxrule=1.3pt,arc=5pt,boxsep=0pt,nobeforeafter,
        left=3pt,right=3pt,top=3pt,bottom=3pt,valign=center]{\qrcode[height=1.9cm]{https://play.google.com/store/apps/details?id=com.luvvix.onbuch&hl=fr}}%
      \hspace{8pt}\begin{minipage}[c]{0.5\linewidth}
        {\popdisplay\fontsize{11}{12.5}\selectfont\color{white}\MakeUppercase{Télécharge OnBuch}}\par\vspace{4pt}
        {\raggedright\popsemi\fontsize{8.4}{11}\selectfont\color{white}Gratuit \textbullet\ Google Play\\Tuteur IA, annales, concours, résultats\par}
      \end{minipage}
    \end{tcolorbox}
  \end{minipage}\hfill
  \begin{minipage}[t]{0.487\linewidth}
    \begin{tcolorbox}[enhanced,colback=poppurple,colframe=popink,boxrule=2.2pt,arc=10pt,
        drop shadow=popink,left=11pt,right=11pt,top=10pt,bottom=10pt,height=3.1cm,valign=center]
      \tikz[baseline=-0.7ex]\node[circle,fill=white,draw=popink,line width=1.3pt,minimum size=1.6cm,inner sep=0]{\popdisplay\fontsize{24}{24}\selectfont\color{poppurple}+};%
      \hspace{8pt}\begin{minipage}[c]{0.6\linewidth}
        {\popdisplay\fontsize{11}{12.5}\selectfont\color{white}\MakeUppercase{Passe à OnBuch+}}\par\vspace{4pt}
        {\raggedright\popsemi\fontsize{8.4}{11}\selectfont\color{white}Tous les cours signés, Tuteur IA illimité, fiches PDF — onglet OnBuch+ de l'app\par}
      \end{minipage}
    \end{tcolorbox}
  \end{minipage}
  \vspace{8pt}
  \popcontenu
  \vfill
  \signatureonbuch
  \restoregeometry
  \newpage
}

% Rangée « Ce cours contient » (page de garde)
\newcommand{\poptile}[3]{% #1 couleur #2 gros texte #3 légende
  \begin{tcolorbox}[enhanced,colback=#1,colframe=popink,boxrule=2pt,arc=9pt,shadow={2.5pt}{-2.5pt}{0pt}{popink},
      left=6pt,right=6pt,top=9pt,bottom=9pt,halign=center,valign=center,height=1.75cm,width=\linewidth,nobeforeafter]
    {\popdisplay\fontsize{12.5}{13}\selectfont\color{white}\MakeUppercase{#2}}\par\vspace{3pt}
    {\centering\popsemi\fontsize{7.8}{9.5}\selectfont\color{white}#3\par}
  \end{tcolorbox}}
\newcommand{\popcontenu}{%
  \par\noindent{\popxboldls\fontsize{7.5}{7.5}\selectfont\color{popmuted}CE COURS SIGNÉ CONTIENT}\par\vspace{7pt}
  \noindent\begin{minipage}[t]{0.235\linewidth}\poptile{popblue}{Cours}{complet, illustré, pas à pas}\end{minipage}\hfill
  \begin{minipage}[t]{0.235\linewidth}\poptile{poporange}{Méthodes}{et exercices résolus}\end{minipage}\hfill
  \begin{minipage}[t]{0.235\linewidth}\poptile{poppink}{Exercices}{type examen + corrigés}\end{minipage}\hfill
  \begin{minipage}[t]{0.235\linewidth}\poptile{popgreen}{Fiche bilan}{l'essentiel à retenir}\end{minipage}\par}

% Sommaire
\newtcolorbox{sommaire}{enhanced,colback=white,colframe=popink,boxrule=2.2pt,arc=10pt,
  drop shadow=popink,left=18pt,right=18pt,top=13pt,bottom=13pt,before skip=6pt,after skip=20pt,
  before upper={\poppill{Sommaire}{poppurple}{white}\par\vspace{9pt}\popsemi\fontsize{10.6}{14.5}\selectfont\color{popink}}}

% Signature de fin de cours — poussée tout en bas de la dernière page
\newcommand{\findecours}{%
  \par\vfill\nopagebreak
  \begin{center}\popsquiggle{poporange}\end{center}\vspace{4pt}
  \signatureonbuch
}
\AtBeginDocument{\def\llbracket{\mathopen{[\mkern-3.5mu[}}\def\rrbracket{\mathclose{]\mkern-3.5mu]}}} % intervalles d'entiers (glyphe ⟦ absent de la police)
% Glyphes absents de Fira Math : redéfinis à partir de glyphes disponibles
\AtBeginDocument{%
  \def\setminus{\mathbin{\backslash}}\let\smallsetminus\setminus
  \def\vdots{\mathord{\vbox{\baselineskip3.2pt\lineskiplimit0pt\kern4pt\hbox{.}\hbox{.}\hbox{.}}}}%
  \def\ddots{\mathinner{\mkern1mu\raise6.5pt\hbox{.}\mkern2mu\raise3.6pt\hbox{.}\mkern2mu\raise0.7pt\hbox{.}\mkern1mu}}%
  \def\triangle{\mathord{\scalebox{0.9}{$\symup{\Delta}$}}}%
  \def\nabla{\mathord{\raisebox{\depth}{\scalebox{1}[-1]{$\symup{\Delta}$}}}}%
  \def\checkmark{\popcheck{popdark}}%
  \def\star{\mathbin{\ast}}\def\bigstar{\mathord{\ast}}%
  \def\diamond{\mathbin{\scalebox{0.6}{$\Diamond$}}}%
  \def\bigcup{\mathop{\mathchoice{\raisebox{-0.3em}{\scalebox{1.7}{$\cup$}}}{\scalebox{1.25}{$\cup$}}{\cup}{\cup}}\displaylimits}%
  \def\bigcap{\mathop{\mathchoice{\raisebox{-0.3em}{\scalebox{1.7}{$\cap$}}}{\scalebox{1.25}{$\cap$}}{\cap}{\cap}}\displaylimits}%
  \def\lesseqqgtr{\mathrel{\lessgtr}}\def\bigoplus{\mathop{\oplus}\displaylimits}%
}
\providecommand{\ssi}{si et seulement si} % abréviation fréquente
\AtBeginDocument{\providecommand{\wideparen}{\overparen}} % arc de cercle

%% FIGFIX-BEGIN v3 — correctifs globaux des figures (docs/onbuch-generation/tools/figfix)
\usepackage{adjustbox}\usepackage{etoolbox}
% toute figure TikZ/pgfplots plus large que la ligne est réduite à la largeur disponible
\BeforeBeginEnvironment{tikzpicture}{\begin{adjustbox}{max width=\linewidth}}
\AfterEndEnvironment{tikzpicture}{\end{adjustbox}}
% légendes pgfplots lisibles par-dessus les courbes
\pgfplotsset{every axis/.append style={legend style={fill=white,fill opacity=0.92,text opacity=1,draw=black!15,inner sep=3pt}}}
% étiquettes d'arêtes (syntaxe edge["..."]) avec fond blanc
\tikzset{every edge quotes/.append style={fill=white,inner sep=1.2pt}}
% bulles de cartes mentales : texte à la ligne dans la bulle, bulle un peu plus grande
\tikzset{every concept/.append style={text width=2.0cm,align=center,minimum size=2.3cm,inner sep=2pt}}
%% FIGFIX-END
\begin{document}

\pagegarde{Œil et la vision}{SVTEEHB}{3ème}{SVTEEHB – classe de 3ème}{N10}{3 séances (fiche de progression harmonisée)}{Cette leçon étudie l'anatomie de l'œil et le mécanisme de la vision, puis les principales anomalies et maladies de la vue (myopie, hypermétropie, presbytie, glaucome, cataracte) ainsi que les règles d'hygiène visuelle. Elle s'appuie sur l'observation de schémas, la dissection d'œil de bœuf et l'exploitation de documents pour aboutir à des applications concrètes de la vie courante.
\vspace{4pt}
\begin{multicols}{2}\small
\begin{itemize}
\item À la fin de cette leçon, je sais identifier et nommer sur un schéma les principales structures de l'œil (cornée, iris, cristallin, rétine, nerf optique...)
\item À la fin de cette leçon, je sais décrire le trajet de la lumière dans l'œil jusqu'à la formation de l'image sur la rétine
\item À la fin de cette leçon, je sais expliquer le mécanisme de l'accommodation à partir de faits expérimentaux
\item À la fin de cette leçon, je sais distinguer et schématiser la myopie, l'hypermétropie et la presbytie, et indiquer leur correction
\item À la fin de cette leçon, je sais décrire le glaucome et la cataracte et leurs conséquences sur la vision
\item À la fin de cette leçon, je sais énoncer et justifier les règles d'hygiène de la vision
\end{itemize}\end{multicols}}

\begin{sommaire}
\begin{multicols}{2}
\begin{enumerate}
\item Anatomie de l'œil
\item Formation de l'image et accommodation
\item Anomalies de la vision : myopie, hypermétropie, presbytie
\item Maladies de la vision : cataracte et glaucome
\item Hygiène de la vision
\item Synthèse et activités d'intégration
\item Activité d'intégration
\item Méthodes et exercices résolus
\item Exercices
\item Corrigés des exercices
\item Fiche bilan
\end{enumerate}
\end{multicols}
\end{sommaire}

\begin{objectifs}
\begin{itemize}
\item À la fin de cette leçon, je sais identifier et nommer sur un schéma les principales structures de l'œil (cornée, iris, cristallin, rétine, nerf optique...)
\item À la fin de cette leçon, je sais décrire le trajet de la lumière dans l'œil jusqu'à la formation de l'image sur la rétine
\item À la fin de cette leçon, je sais expliquer le mécanisme de l'accommodation à partir de faits expérimentaux
\item À la fin de cette leçon, je sais distinguer et schématiser la myopie, l'hypermétropie et la presbytie, et indiquer leur correction
\item À la fin de cette leçon, je sais décrire le glaucome et la cataracte et leurs conséquences sur la vision
\item À la fin de cette leçon, je sais énoncer et justifier les règles d'hygiène de la vision
\item À la fin de cette leçon, je sais interpréter des documents (schémas, clichés, résultats d'examen ophtalmologique) pour diagnostiquer un trouble de la vision
\item À la fin de cette leçon, je sais rédiger et justifier une démarche, une réponse ou une conclusion dans un problème de synthèse
\end{itemize}
\end{objectifs}

\begin{prerequis}
\begin{itemize}
\item Notion de récepteur sensoriel et de message nerveux (4ème : sensibilité et système nerveux)
\item Propagation rectiligne de la lumière et formation des images (physique, 4ème)
\item Rôle du cerveau dans la réception et l'interprétation des messages nerveux
\item Notion de lentille convergente et de foyer (complément de physique utile)
\item Conduite d'une observation et exploitation d'un schéma annoté
\end{itemize}
\end{prerequis}

\newpage

\definecolor{choroid}{RGB}{40,20,10}
\definecolor{sclera}{RGB}{230,230,230}
\definecolor{retina}{RGB}{200,150,140}
\definecolor{cornea}{RGB}{200,230,255}
\definecolor{iriscol}{RGB}{80,50,30}
\definecolor{lenscol}{RGB}{220,210,190}
\definecolor{vitreous}{RGB}{235,245,255}
\definecolor{aqueous}{RGB}{220,240,255}
\definecolor{nerve}{RGB}{255,200,180}

\coursec{Anatomie de l'œil}

À Yaoundé, un élève de 3ème remarque qu'il ne distingue plus clairement le tableau depuis le fond de la classe, alors qu'il lit sans peine son cahier. Sa grand-mère, elle, écarte le journal à bout de bras pour lire, et son oncle chauffeur de taxi-moto se plaint de maux de tête après avoir longtemps roulé sous le soleil de poussière. Comment l'œil nous permet-il de voir net à différentes distances ? Pourquoi ces troubles apparaissent-ils et comment peut-on les corriger ou les prévenir ? Cette leçon répond à ces questions en explorant l'organe de la vision.

\courssub{Aspect externe et protections de l'œil}

\begin{experience}[Observation de l'aspect externe de l'œil]
\textbf{Matériel :} Miroir, lampe de poche, photographie d'œil humain (ou observation mutuelle entre binômes).
\textbf{Protocole :}
\begin{enumerate}
\item Observer son propre œil dans un miroir sous un éclairage modéré. Noter les structures visibles : paupières, cils, conjonctive (blanc de l'œil), iris (partie colorée), pupille (orifice central noir).
\item Éclairer brièvement l'œil avec la lampe de poche (de côté, jamais en face) et observer la réaction de la pupille.
\item Identifier le rôle protecteur de chaque structure.
\end{enumerate}
\textbf{Observations :} Les paupières se ferment par réflexe (clignement) si un objet approche trop vite ; les cils retiennent la poussière ; la conjonctive est une muqueuse transparente et vascularisée qui recouvre la sclérotique et tapisse l'intérieur des paupières ; l'iris présente une coloration variable (marron, noir, bleu, vert) ; la pupille apparaît comme un trou noir au centre de l'iris ; sous l'effet de la lumière, le diamètre de la pupille diminue (myosis), dans l'obscurité il augmente (mydriase).
\textbf{Interprétation :} L'œil est un organe fragile, protégé par des annexes externes. La variation du diamètre pupillaire régule la quantité de lumière entrant dans l'œil (réflexe pupillaire).
\textbf{Conclusion :} L'aspect externe révèle les premières structures de protection et de régulation lumineuse.
\end{experience}

\begin{definition}[L'œil]
L'œil est l'\cle{organe récepteur sensoriel} de la vision. Il transforme l'\cle{excitation lumineuse} (onde électromagnétique visible, \SIrange{400}{700}{\nano\meter}) en \cle{message nerveux} (influx nerveux) transmis au cerveau par le nerf optique pour y être interprété (sensation visuelle).
\end{definition}

\begin{propriete}[Annexes de protection de l'œil]
\begin{itemize}
\item \textbf{Paupières :} replis cutanéo-muqueux mobiles ; protègent par le clignement (réflexe de défense) et répartissent le film lacrymal.
\item \textbf{Cils :} poils implantés sur le bord libre des paupières ; filtrent les particules (poussière, insectes) et déclenchent le réflexe de clignement au contact.
\item \textbf{Conjonctive :} muqueuse transparente, vascularisée, continue avec la cornée (limbe) et la muqueuse palpebrale ; sécrète du mucus (composante du film lacrymal) et assure une défense immunitaire locale.
\item \textbf{Glandes lacrymales :} situées en haut et en externe de l'orbite ; sécrètent les \cle{larmes} (eau, sels, lysozyme bactéricide, lipides) qui lubrifient, nourrissent la cornée (avasculaire), nettoient et désinfectent la surface oculaire. Les larmes s'écoulent vers les points lacrymaux (angle interne) puis le canal lacrymo-nasal (d'où l'écoulement nasal en pleurant).
\end{itemize}
\end{propriete}

\begin{attention}[Iris vs Pupille : ne pas confondre !]
C'est l'erreur n°1 en observation et en schéma.
\begin{itemize}
\item L'\cle{iris} est une \textbf{structure} musculaire et pigmentée (le « couleur des yeux »), en forme de diaphragme annulaire.
\item La \cle{pupille} est l'\textbf{orifice central} de l'iris (le « trou noir »). Ce n'est \emph{pas} une structure matérielle, mais une \emph{ouverture} dont le diamètre varie.
\end{itemize}
\textbf{Astuce :} L'iris \emph{est} coloré ; la pupille \emph{laisse passer} la lumière (donc elle paraît noire car la lumière n'en ressort pas).
\end{attention}

\courssub{Anatomie interne : coupe sagittale du globe oculaire}

Pour aller plus loin que l'aspect externe, il faut « ouvrir » l'œil. En classe, nous réalisons la dissection d'un œil de bœuf (proche de l'œil humain par sa taille et son organisation) ou nous étudions un schéma de coupe sagittale médiane.

\begin{experience}[Dissection d'un œil de bœuf (TP)]
\textbf{Matériel :} Œil de bœuf frais (boucherie), plateau de dissection, scalpel, ciseaux, pince, loupe, gants, papier journal, texte imprimé (petits caractères).
\textbf{Protocole :}
\begin{enumerate}
\item Observer l'œil entier : identifier le nerf optique (arrière), les muscles oculomoteurs (insertions latérales), la cornée (antérieure, transparente), la sclérotique (blanc opaque).
\item Inciser la sclérotique au niveau de l'équateur (milieu) avec le scalpel, puis couper tout le tour avec les ciseaux pour ouvrir le globe en deux hémisphères (antérieur et postérieur).
\item Dans l'hémisphère antérieur : observer l'iris, la pupille, extraire le cristallin (lentille transparente, élastique). Le poser sur le texte imprimé : il grossit les caractères (rôle de lentille convergente).
\item Dans l'hémisphère postérieur : observer la rétine (fine membrane translucide, rosée), la choroïde (couche foncée, très vascularisée, pigmentée en noir), la sclérotique (blanche, fibreuse, dure). Noter l'adhérence rétine/choroïde sauf au niveau de la tête du nerf optique (tache aveugle).
\item Observer l'humeur vitrée (gel transparent) remplissant la cavité postérieure.
\end{enumerate}
\textbf{Observations clés :} Le cristallin est une lentille biconvexe, dure chez le bœuf (plus souple chez l'humain jeune). La choroïde est noire (mélanine) : elle absorbe la lumière parasite (évite les reflets internes). La rétine se décolle facilement de la choroïde (risque de décollement de rétine). Le nerf optique émerge en un point dépourvu de photorécepteurs.
\textbf{Interprétation :} L'œil est une chambre noire sphérique à parois rigides (sclérotique), dont la paroi interne (rétine) reçoit l'image formée par un système optique (cornée + cristallin). Les milieux traversés sont transparents.
\end{experience}

\begin{exemplebox}[Le cristallin de bœuf, une loupe naturelle]
En posant le cristallin extrait sur une ligne de texte (police 8 pt), on observe un grossissement net des caractères. Cela prouve sa géométrie \cle{convergente} (biconvexe) et sa transparence. Chez l'humain jeune, il est élastique et change de forme (accommodation) ; chez le bœuf adulte comme chez le personne âgée, il durcit (sclérose du cristallin), base de la \cle{presbytie}.
\end{exemplebox}

\begin{popfigure}
\begin{tikzpicture}[scale=1.1, every node/.style={font=\small}]
% Couleurs

\definecolor{choroid}{RGB}{40,20,10}







% Sclérotique (extérieur)
\draw[fill=sclera, draw=black, thick] (-3.5,0) .. controls (-3.5,2.2) and (3.5,2.2) .. (3.5,0) .. controls (3.5,-2.2) and (-3.5,-2.2) .. cycle;
% Choroïde
\draw[fill=choroid, draw=black] (-3.3,0) .. controls (-3.3,2.0) and (3.3,2.0) .. (3.3,0) .. controls (3.3,-2.0) and (-3.3,-2.0) .. cycle;
% Rétine
\draw[fill=retina, draw=black] (-3.1,0) .. controls (-3.1,1.85) and (3.1,1.85) .. (3.1,0) .. controls (3.1,-1.85) and (-3.1,-1.85) .. cycle;
% Cornée (bombement antérieur)
\draw[fill=cornea, draw=black, thick] (-3.5,0) .. controls (-3.5,0.6) and (-1.5,0.8) .. (0,0.8) .. controls (1.5,0.8) and (3.5,0.6) .. (3.5,0);
% Limbe (jonction cornée/sclérotique)
\draw[dashed, thin] (-3.5,0) -- (-3.5,0.1);
% Iris et pupille
\draw[fill=iriscol, draw=black, thick] (-1.2,0.3) -- (-1.2,-0.3) .. controls (-1.0,-0.5) and (1.0,-0.5) .. (1.2,-0.3) -- (1.2,0.3) .. controls (1.0,0.5) and (-1.0,0.5) .. cycle;
\draw[fill=black] (-0.6,0) ellipse (0.3 and 0.3); % Pupille (myosis modéré)
% Cristallin
\draw[fill=lenscol, draw=black, thick] (-0.8,0.6) .. controls (-0.8,1.0) and (0.8,1.0) .. (0.8,0.6) .. controls (0.8,0.2) and (-0.8,0.2) .. cycle;
% Zonule (fibres de Zinn) - suggérées
\foreach \x in {-0.7,-0.3,0.3,0.7} {
 \draw[thin, gray] (\x,0.6) -- (\x*1.5,1.2);
 \draw[thin, gray] (\x,-0.6) -- (\x*1.5,-1.2); % pour postérieur mais invisible ici
}
% Humour aqueuse (antérieur cristallin)
\draw[fill=aqueous, draw=none] (-1.2,0.3) .. controls (-1.2,0.5) and (1.2,0.5) .. (1.2,0.3) -- (1.2,-0.3) .. controls (1.2,-0.5) and (-1.2,-0.5) .. (-1.2,-0.3) -- cycle;
% Correction : zone antérieure à l'iris et postérieure à la cornée
% On simplifie : zone entre cornée et iris/cristallin
% Humour vitrée (postérieur cristallin)
\draw[fill=vitreous, draw=none] (-3.1,0) .. controls (-3.1,1.8) and (0.8,1.5) .. (0.8,0.6) .. controls (0.8,0.2) and (-3.1,-1.8) .. (-3.1,-1.85) .. controls (-3.1,-1.8) and (0.8,-1.5) .. (0.8,-0.6) .. controls (0.8,-0.2) and (-3.1,0) .. cycle;
% Note: pgfplots/tikz fill between would be better but manual path is safer for compatibility.
% Nerf optique
\draw[fill=nerve, draw=black, thick] (-3.5,-0.5) -- (-5,-0.4) -- (-5,0.4) -- (-3.5,0.5) -- cycle;
% Tache aveugle (disque optique) - zone sans rétine/choroïde
\draw[fill=choroid, draw=black] (-3.3,-0.3) .. controls (-3.3,-0.1) and (-3.0,-0.1) .. (-3.0,-0.3) .. controls (-3.0,-0.5) and (-3.3,-0.5) .. cycle; % trou dans choroide
\draw[fill=retina, draw=black] (-3.1,-0.25) .. controls (-3.1,-0.1) and (-2.9,-0.1) .. (-2.9,-0.25) .. controls (-2.9,-0.4) and (-3.1,-0.4) .. cycle; % trou dans retine
% Muscles oculomoteurs (insertions)
\draw[thick, red!70!black] (-3.5,1.5) -- (-4.5,2.2) node[above left, black] {Muscle droit sup.};
\draw[thick, red!70!black] (-3.5,-1.5) -- (-4.5,-2.2) node[below left, black] {Muscle droit inf.};
\draw[thick, red!70!black] (3.5,1.5) -- (4.5,2.2) node[above right, black] {Muscle droit lat.};
\draw[thick, red!70!black] (3.5,-1.5) -- (4.5,-2.2) node[below right, black] {Muscle droit méd.};
% Flèches lumière
\draw[->, very thick, yellow!80!orange] (0,2.5) -- (0,0.9) node[midway, right] {Lumière};
\draw[->, very thick, yellow!80!orange] (0,0.9) -- (0,0.2) node[midway, right] {Cornée};
\draw[->, very thick, yellow!80!orange] (0,0.2) -- (0,-0.1) node[midway, right] {Pupille};
\draw[->, very thick, yellow!80!orange] (0,-0.1) -- (0,-0.5) node[midway, right] {Cristallin};
\draw[->, very thick, yellow!80!orange] (0,-0.5) -- (0,-1.5) node[midway, right] {Humeur vitrée};
\draw[->, very thick, yellow!80!orange] (0,-1.5) -- (0,-1.85) node[midway, right] {Rétine};
% Légendes par repères
\node[align=left, anchor=west] at (4.0, 2.0) {\textbf{1} Cornée};
\node[align=left, anchor=west] at (4.0, 1.5) {\textbf{2} Humour aqueuse};
\node[align=left, anchor=west] at (4.0, 0.5) {\textbf{3} Iris};
\node[align=left, anchor=west] at (4.0, 0.0) {\textbf{4} Pupille};
\node[align=left, anchor=west] at (4.0, -0.5) {\textbf{5} Cristallin};
\node[align=left, anchor=west] at (4.0, -1.2) {\textbf{6} Humour vitrée};
\node[align=left, anchor=west] at (4.0, -1.8) {\textbf{7} Rétine};
\node[align=left, anchor=west] at (4.0, -2.3) {\textbf{8} Choroïde};
\node[align=left, anchor=west] at (4.0, -2.8) {\textbf{9} Sclérotique};
\node[align=left, anchor=west] at (-6.0, 0.0) {\textbf{10} Nerf optique};
\node[align=left, anchor=west] at (-6.0, -1.5) {\textbf{11} Tache aveugle};
\end{tikzpicture}
\legende{Coupe sagittale médiane de l'œil humain (schéma simplifié). Trajet de la lumière (flèches orange) de la cornée vers la rétine. Les trois tuniques (sclérotique, choroïde, rétine) et les milieux transparents sont identifiés.}
\end{popfigure}

\begin{methode}[Lire et légender un schéma anatomique de l'œil (ordre de la lumière)]
Pour ne rien oublier et respecter le sens physiologique, suis le parcours de la lumière :
\begin{enumerate}
\item \textbf{Milieux transparents (dioptries) :} Cornée $\to$ Humour aqueuse $\to$ Cristallin $\to$ Humour vitrée. (Ce sont les \cle{dioptries} : surfaces séparant deux milieux transparents d'indices différents).
\item \textbf{Diaphragme :} Iris (muscles sphincter et dilatateur) $\to$ Pupille (orifice).
\item \textbf{Tuniques (de l'extérieur vers l'intérieur) :} Sclérotique (blanche, fibreuse, protection/maintien) $\to$ Choroïde (noire, vascularisée, nutrition + absorption lumière) $\to$ Rétine (translucide, neurosensorielle, photoréception).
\item \textbf{Sortie du message :} Rétine $\to$ Nerf optique (tache aveugle) $\to$ Chiasma optique $\to$ Cortex visuel occipital.
\end{enumerate}
\textbf{Astuce examen :} Sur un schéma à légender, commence toujours par la cornée (antérieur) et finis par le nerf optique (postérieur).
\end{methode}

\begin{propriete}[Les trois tuniques du globe oculaire]
\begin{center}
\begin{tabularx}{\linewidth}{|l|X|X|X|}\hline
\rowcolor{popblueL}
\textbf{Tunique} & \textbf{Structure / Situation} & \textbf{Composition histologique} & \textbf{Rôle principal} \\ \hline
\cle{Sclérotique} (tunique fibreuse) & Postérieure (5/6) + Cornée (antérieure 1/6) & Tissu conjonctif dense, fibres de collagène (blanc opaque) ; Cornée : collagène ordonné + kératocytes (transparence) & \textbf{Protection} mécanique, maintien de la forme sphérique, point d'ancrage des muscles oculomoteurs. Cornée = première lentille convergente (puissance $\approx$ 43 D). \\ \hline
\cle{Choroïde} (tunique vasculaire) & Moyenne, entre sclérotique et rétine & Très vascularisée (capillaires), riche en \textbf{mélanine} (pigment noir) & \textbf{Nutrition} de la rétine externe (photorécepteurs) et de la sclérotique ; \textbf{absorption de la lumière} parasite (évite la diffusion interne, « chambre noire »). \\ \hline
\cle{Rétine} (tunique interne / neurosensorielle) & Interne, au contact de l'humeur vitrée & \textbf{Épithélium pigmentaire} (extérieur) + \textbf{Couche nerveuse} (intérieure : photorécepteurs, bipolaires, ganglionnaires) & \textbf{Réception lumineuse} (photoréception) et début du traitement de l'image (circuits rétiniens). Se décolle facilement de la choroïde (risque pathologique). \\ \hline
\end{tabularx}
\end{center}
\end{propriete}

\begin{propriete}[Les milieux transparents (dioptries) traversés par la lumière]
\begin{itemize}
\item \textbf{Cornée :} Fenêtre antérieure, avasculaire (nutrition par larmes et humour aqueuse), indice $n \approx 1,376$. Puissance $\approx$ \SI{43}{\per\meter}. Fixe (ne s'accommode pas).
\item \textbf{Humeur aqueuse :} Liquide clair (eau $99\%$, sels, glucose, acide ascorbique, protéines), sécrétée par les procès ciliaires (corps ciliaire), circule chambre postérieure $\to$ pupille $\to$ chambre antérieure $\to$ réseau trabéculaire $\to$ veine sclérale (canal de Schlemm). Indice $n \approx 1,336$. Maintient la pression intraoculaire (PIO $\approx$ \SI{10}{--}\SI{21}{mmHg}) et nourrit cornée/cristallin.
\item \textbf{Cristallin :} Lentille biconvexe, transparente, élastique (jeune), avasculaire. Indice $n \approx 1,406$ (noyau) à $1,386$ (cortex). Puissance variable $\approx$ \SI{19}{--}\SI{33}{\per\meter} (accommodation). Enveloppé dans une capsule élastique, suspendu par le \cle{ligament suspenseur (zonule de Zinn)} au corps ciliaire.
\item \textbf{Humeur vitrée :} Gel transparent (eau $99\%$, acide hyaluronique, collagène), indice $n \approx 1,336$. Remplit la cavité vitrée (4/5 du volume), maintient la rétine appliquée sur la choroïde, transparence optique.
\end{itemize}
\end{propriete}

\begin{aretenir}[L'iris : diaphragme vivant]
L'iris est un diaphragme circulaire pigmenté (mélanine = couleur de l'œil) percé en son centre de la \cle{pupille}. Il contient deux muscles lisses antagonistes innervés par le système nerveux autonome :
\begin{itemize}
\item \textbf{Sphincter de la pupille} (circulaire) : innervation \textbf{parasympathique} (nerf oculomoteur III) $\to$ \cle{myosis} (rétrécissement) en lumière forte (réflexe photomoteur direct et consensuel).
\item \textbf{Dilatateur de la pupille} (radiaire) : innervation \textbf{sympathique} $\to$ \cle{mydriase} (dilatation) dans l'obscurité, émotion, douleur.
\end{itemize}
C'est le \cle{réflexe pupillaire} : adaptation de la quantité de lumière atteignant la rétine et augmentation de la profondeur de champ (petite pupille = meilleur piqué, comme diaphragme photo).
\end{aretenir}

\courssub{La rétine : siège de la réception lumineuse}

La rétine est une fine membrane (\SI{\approx 200}{\micro\meter}) transparente, inversée (la lumière traverse les couches nerveuses avant d'atteindre les photorécepteurs au fond).

\begin{popfigure}
\begin{tikzpicture}[scale=1.3, every node/.style={font=\small}]
% Couches rétiniennes (schéma simplifié vertical)
% Lumière arrive par le HAUT (depuis le corps vitré)
% 1. Limite interne (membrane basale)
\draw[thick, black] (-2,2.5) -- (2,2.5) node[midway, above] {Corps vitré / Humour vitrée};
% 2. Couche des fibres nerveuses (axones ganglionnaires)
\draw[fill=popblue!20, draw=popblue] (-2,2.3) -- (2,2.3) -- (2,2.1) -- (-2,2.1) -- cycle;
\node at (0,2.2) {Couche des fibres nerveuses (axones $\to$ Nerf optique)};
% 3. Couche des cellules ganglionnaires
\draw[fill=popgreen!20, draw=popgreen] (-2,2.0) -- (2,2.0) -- (2,1.7) -- (-2,1.7) -- cycle;
\foreach \x in {-1.5,-0.5,0.5,1.5} \draw[fill=popgreen] (\x,1.85) circle (0.08);
\node at (0,1.85) {Corps cellulaires des \textbf{cellules ganglionnaires}};
% 4. Couche plexiforme interne (synapses bipolaires-ganglionnaires)
\draw[fill=poporange!20, draw=poporange] (-2,1.6) -- (2,1.6) -- (2,1.4) -- (-2,1.4) -- cycle;
\node at (0,1.5) {Couche plexiforme interne};
% 5. Couche des cellules bipolaires
\draw[fill=poppurple!20, draw=poppurple] (-2,1.3) -- (2,1.3) -- (2,1.0) -- (-2,1.0) -- cycle;
\foreach \x in {-1.7,-0.7,0.3,1.3} \draw[fill=poppurple] (\x,1.15) circle (0.07);
\node at (0,1.15) {Corps cellulaires des \textbf{cellules bipolaires}};
% 6. Couche plexiforme externe (synapses photorécepteurs-bipolaires)
\draw[fill=poppink!20, draw=poppink] (-2,0.9) -- (2,0.9) -- (2,0.7) -- (-2,0.7) -- cycle;
\node at (0,0.8) {Couche plexiforme externe};
% 7. Couche nucléaire externe (corps photorécepteurs)
\draw[fill=popgold!30, draw=popgold] (-2,0.6) -- (2,0.6) -- (2,0.0) -- (-2,0.0) -- cycle;
% Bâtonnets (nombreux, fins)
\foreach \x in {-1.8,-1.4,-1.0,-0.6,-0.2,0.2,0.6,1.0,1.4,1.8} \draw[fill=black] (\x,0.3) rectangle ++(0.15,0.25);
% Cônes (moins nombreux, larges)
\foreach \x in {-1.6,-0.8,0,0.8,1.6} \draw[fill=poporange] (\x,0.3) rectangle ++(0.25,0.25);
\node at (0,0.3) {Couche nucléaire externe : \textbf{Bâtonnets} (noirs, fins) \& \textbf{Cônes} (orange, larges)};
% 8. Segments externes (photoréception)
\draw[fill=popred!30, draw=popred] (-2,-0.1) -- (2,-0.1) -- (2,-0.6) -- (-2,-0.6) -- cycle;
% Disques bâtonnets
\foreach \x in {-1.8,-1.4,-1.0,-0.6,-0.2,0.2,0.6,1.0,1.4,1.8} {
 \foreach \y in {-0.15,-0.25,-0.35,-0.45} \draw[fill=black] (\x,\y) rectangle ++(0.15,0.05);
}
% Segments cônes
\foreach \x in {-1.6,-0.8,0,0.8,1.6} {
 \foreach \y in {-0.15,-0.25,-0.35,-0.45} \draw[fill=poporange] (\x,\y) rectangle ++(0.25,0.05);
}
\node at (0,-0.35) {Segments externes : disques membranaires (rhodopsine / iodopsines)};
% 9. Épithélium pigmentaire (RPE)
\draw[fill=choroid, draw=black, thick] (-2,-0.7) -- (2,-0.7) -- (2,-0.9) -- (-2,-0.9) -- cycle;
\node at (0,-0.8) {\textbf{Épithélium pigmentaire rétinien (RPE)} : mélanine, phagocytose disques usagés, barrière hémato-rétinienne};
% Flèche lumière
\draw[->, very thick, yellow!80!orange] (0,2.7) -- (0,-0.65) node[midway, right, black] {Lumière};
% Repères Macula / Tache aveugle (schéma latéral mini)
\begin{scope}[shift={(3.5,0.5)}, scale=0.6]
\draw[fill=retina] (-1.5,0) arc (180:0:1.5 and 0.5) -- (1.5,-2) arc (0:-180:1.5 and 0.5) -- cycle;
\draw[fill=popgreen!50] (0,0) circle (0.3) node[above] {Macula / Fovéa};
\draw[fill=white, draw=black, thick] (-1.2,-1) rectangle (-0.6,-1.6) node[midway, align=center] {Tache\\aveugle};
\draw[->, thick] (-1.5,-1.3) -- (-2.2,-1.3) node[left] {Nerf optique};
\node at (0,1.2) {\textbf{Topographie rétinienne}};
\end{scope}
\end{tikzpicture}
\legende{Structure histologique simplifiée de la rétine (coupe transverse). La lumière traverse les couches nerveuses (haut) pour atteindre les segments externes des photorécepteurs (bas). Encadré : topographie avec macula (vision fine, cônes seulement) et tache aveugle (sortie nerf optique, pas de photorécepteurs).}
\end{popfigure}

\begin{definition}[Photorécepteurs : bâtonnets et cônes]
Ce sont les cellules sensorielles spécialisées dans la \cle{transduction} signal lumineux $\to$ signal électrique (potentiel de récepteur).
\begin{itemize}
\item \textbf{Bâtonnets (\SI{\approx 120}{millions}) :} Très sensibles (vision \cle{scotopique}, nuit, seuil bas), achromatiques (une seule photopigment : \cle{rhodopsine} ou pourpre rétinien), faible acuité (convergence forte : plusieurs bâtonnets $\to$ 1 bipolaire), absents de la fovéa, maxi densité à \SI{20}{\degree} de la fovéa (vision latérale nuit).
\item \textbf{Cônes (\SI{\approx 6}{millions}) :} Moins sensibles (vision \cle{photopique}, jour, couleur), trois types (S, M, L = bleu, vert, rouge) $\to$ \cle{vision trichromatique}, haute acuité (convergence 1:1 en fovéa), concentrés en \cle{macula} (zone jaune, \SI{\approx 5,5}{mm} de diamètre) et surtout \cle{fovéa} (centre, \SI{\approx 1,5}{mm}, cônes seulement, acuité max).
\end{itemize}
\end{definition}

\begin{propriete}[Macula et Tache aveugle]
\begin{itemize}
\item \textbf{Macula lutea (tache jaune) :} Zone centrale de la rétine, riche en xanthophylles (pigments jaunes filtrant le bleu), centre = \textbf{fovéa} (dépression, cônes, vascularisation nulle $\to$ transparence max). C'est la zone de \textbf{vision précise, colorée, diurne} (lecture, reconnaissance faciale).
\item \textbf{Tache aveugle (point de Mariotte / disque optique) :} Zone d'émergence du nerf optique et des vaisseaux sanguins centraux. \textbf{Absence totale de photorécepteurs} (ni cônes, ni bâtonnets) $\to$ \textbf{zone insensible à la lumière}. Le cerveau « comble » le trou par interpolation (on ne la perçoit pas spontanément).
\end{itemize}
\end{propriete}

\begin{savaistu}[L'expérience de la tache aveugle (Mariotte, 1668)]
Dessine une croix et un point espacés de \SI{\approx 7}{cm} sur une feuille blanche. Ferme l'œil gauche, fixe la croix de l'œil droit à \SI{\approx 30}{cm}. Approche lentement la feuille : le point \textbf{disparaît} quand son image tombe sur la tache aveugle, puis réapparaît. Preuve anatomique directe d'une zone rétinienne sans photorécepteurs. Essaie chez toi ce soir !
\end{savaistu}

\courssub{Le nerf optique et les annexes oculomotrices}

\begin{propriete}[Nerf optique (IIème paire crânienne)]
Regroupe les axones des \SI{\approx 1,2}{million} de cellules ganglionnaires de la rétine. Traverse la sclérotique au niveau de la \cle{lamelle cribrosa} (zone de faiblesse, siège du glaucome). Parcourt l'orbite, passe par le \cle{canal optique} (os sphénoïde), atteint le \cle{chiasma optique} (décroisement partiel : fibres nasales croisent, fibres temporales restent) $\to$ \cle{tractus optique} $\to$ \cle{corps genouillé latéral} (thalamus) $\to$ \cle{radiation optique} $\to$ \cle{cortex visuel primaire (aire 17, lobe occipital)}. Myélinisé (blanc) après la sortie du globe (pas de myéline dans la rétine $\to$ transparence).
\end{propriete}

\begin{propriete}[Muscles oculomoteurs (motricité extrinsèque)]
Six muscles squelettiques (striés, volontaires/réflexes) insérés sur la sclérotique, innervés par 3 nerfs crâniens :
\begin{itemize}
\item \textbf{4 muscles droits :} Supérieur (III), Inférieur (III), Médial / Interne (III), Latéral / Externe (VI). Mouvements : élévation, dépression, adduction (vers nez), abduction (vers tempe).
\item \textbf{2 muscles obliques :} Supérieur (IV - trochléaire, passe par une poulie/trochlée), Inférieur (III). Mouvements : rotation (intorsion/extorsion), aide à élévation/dépression.
\end{itemize}
Permettent les \cle{mouvements conjugués} (versions : les deux yeux même direction) et \cle{vergences} (mouvements désjugués : convergence/divergence pour la vision de près/loin). Coordination par noyaux du tronc cérébral (PPRF, noyau MLF).
\end{propriete}

\begin{propriete}[Appareil lacrymal et film lacrymal]
\begin{itemize}
\item \textbf{Glande lacrymale principale :} Fosse lacrymale (os frontal, haut-externe). Sécrétion continue (basale) et réflexe (irritation, émotion).
\item \textbf{Film lacrymal (3 couches) :} Lipide (externe, glandes de Meibomius, évite évaporation) / Aqueuse (moyenne, glande lacrymale, O$_2$, nutriments, lysozyme) / Mucine (interne, cellules caliciformes conjonctives, adhérence cornée).
\item \textbf{Drainage :} Points lacrymaux (sup/inf, bord libre paupières) $\to$ Canaux lacrymaux $\to$ Sac lacrymal $\to$ Canal lacrymo-nasal $\to$ Fosses nasales (écoulement nasal).
\end{itemize}
Rôles : \textbf{Optique} (surface lisse cornée), \textbf{Nutritif} (cornée avasculaire), \textbf{Protection} (lavage, lysozyme, IgA), \textbf{Mécanique} (lubrification paupières).
\end{propriete}

\begin{exoresolu}[Identification des structures sur un schéma de coupe sagittale]
\textbf{Énoncé :} Le document ci-contre représente une coupe sagittale de l'œil humain. Les flèches 1 à 10 indiquent des structures anatomiques. Complète le tableau en associant chaque numéro au nom de la structure et à sa fonction principale.
\textbf{Document :} (Schéma type examen : œil vu de côté, flèches sur Cornée, Iris, Pupille, Cristallin, Rétine, Choroïde, Sclérotique, Nerf optique, Humour aqueuse, Humour vitrée).

\begin{center}
\begin{tabularx}{\linewidth}{|c|X|X|}\hline
\rowcolor{popblueL}
\textbf{N°} & \textbf{Structure} & \textbf{Fonction principale} \\ \hline
1 & Cornée & Première lentille convergente (puissance \SI{\approx 43}{D}), protection, transparence. \\ \hline
2 & Humour aqueuse & Maintien pression intraoculaire, nutrition cornée/cristallin, transparence. \\ \hline
3 & Iris & Diaphragme pigmenté, régule diamètre pupillaire (quantité lumière). \\ \hline
4 & Pupille & Orifice central de l'iris, laisse passer la lumière. \\ \hline
5 & Cristallin & Lentille convergente \textbf{variable} (accommodation), mise au point. \\ \hline
6 & Humour vitrée & Maintien forme globe, transparence, appui rétine sur choroïde. \\ \hline
7 & Rétine & Réception lumineuse (photorécepteurs), traitement signal, message nerveux. \\ \hline
8 & Choroïde & Vascularisation (nutrition rétine externe), pigmentation (absorption lumière). \\ \hline
9 & Sclérotique & Protection, maintien forme, insertion muscles oculomoteurs. \\ \hline
10 & Nerf optique & Transmission message nerveux rétine $\to$ cerveau. \\ \hline
\end{tabularx}
\end{center}

\textbf{Solution détaillée :} Le raisonnement suit le parcours de la lumière (Méthode p. X) et la connaissance des tuniques (de l'extérieur vers l'intérieur).
\begin{itemize}
\item Flèches antérieures (devant le cristallin) : Cornée (1, bombée, transparente), Humour aqueuse (2, liquide clair), Iris (3, coloré, musculaire), Pupille (4, trou noir central).
\item Cristallin (5) : lentille biconvexe centrale, suspendue par zonule.
\item Cavité postérieure : Humour vitrée (6, gel remplissant le volume).
\item Paroi postérieure (tuniques) : Rétine (7, interne, fine, rosée), Choroïde (8, moyenne, noire, vascularisée), Sclérotique (9, externe, blanche, fibreuse).
\item Sortie : Nerf optique (10, blanc, postéro-médian).
\end{itemize}
\textbf{Piège :} Ne pas inverser Choroïde (noire, vasculaire) et Rétine (translucide, nerveuse). La Rétine est \emph{interne} (côté vitrée).
\end{exoresolu}

\begin{exercice}[Entraînement : Anatomie et fonctions]
\begin{enumerate}
\item Sur un schéma d'œil non légendé fourni à part, place et nomme : Cornée, Cristallin, Rétine, Choroïde, Sclérotique, Nerf optique, Iris, Pupille, Humour aqueuse, Humour vitrée, Muscle droit médial, Glande lacrymale.
\item Complète : La lumière traverse dans l'ordre : \ldots $\to$ \ldots $\to$ \ldots $\to$ \ldots $\to$ Rétine.
\item Pourquoi la choroïde est-elle pigmentée en noir (mélanine) ?
\item Quelle est la différence fondamentale entre la fovéa et la tache aveugle en termes de photorécepteurs et d'acuité visuelle ?
\item Un patient a une pression intraoculaire de \SI{28}{mmHg}. Quel liquide est en cause ? Par où s'écoule-t-il normalement ? Quelle pathologie redoute-t-on ?
\end{enumerate}
\end{exercice}

\begin{corrige}[Exercice Anatomie et fonctions]
\begin{enumerate}
\item Vérifie ton schéma avec la figure de cours (p. X). Respecte l'ordre des tuniques (Sclérotique ext. $\to$ Choroïde $\to$ Rétine int.) et des milieux (Cornée $\to$ Aqueuse $\to$ Cristallin $\to$ Vitrée).
\item Cornée $\to$ Humour aqueuse $\to$ Cristallin $\to$ Humour vitrée.
\item Pour absorber la lumière qui a traversé la rétine et éviter les réflexions parasites (phénomène de « chambre noire »), améliorant le contraste de l'image.
\item Fovéa : cônes (densité max \SI{\approx 150000}{mm^{-2}}), acuité max (vision fine, lecture), vision des couleurs. Tache aveugle : 0 photorécepteur (sortie nerf optique + vaisseaux), acuité nulle (zone aveugle physiologique).
\item L'humeur aqueuse. Écoulement normal : Chambre postérieure $\to$ Pupille $\to$ Chambre antérieure $\to$ Réseau trabéculaire (angle irido-cornéen) $\to$ Canal de Schlemm $\to$ Veines sclérales $\to$ Circulation veineuse générale. Pathologie : \cle{Glaucome} (hypertonie oculaire $\to$ compression fibres nerf optique à la lamelle cribrosa $\to$ excavation papillaire $\to$ perte champ visuel $\to$ cécité).
\end{enumerate}
\end{corrige}

\coursec{Formation de l'image et accommodation}

Nous venons d'étudier l'anatomie de l'œil : ses enveloppes, ses milieux transparents et ses structures internes. Nous savons maintenant \emph{où} se trouvent la cornée, le cristallin, l'iris, la rétine et le nerf optique. Mais comment cet ensemble transforme-t-il la lumière en une image nette que le cerveau peut interpréter ? C'est ce que nous allons découvrir en suivant le trajet de la lumière, en comprenant le rôle crucial du cristallin comme lentille vivante et déformable, et en analysant le mécanisme d'accommodation.

\courssub{Trajet de la lumière et formation de l'image sur la rétine}

La lumière émise ou réfléchie par un objet traverse les milieux transparents de l'œil dans un ordre précis avant d'atteindre les cellules photosensibles.

\begin{experience}[Trajet des rayons lumineux dans l'œil et formation de l'image]
\textbf{Observation (modèle de la chambre noire) :} On place une lentille convergente (modélisant le système cornée-cristallin) devant un écran (modélisant la rétine) dans une boîte obscure. Un objet lumineux est placé devant la lentille.
\textbf{Protocole :} On déplace l'écran jusqu'à obtenir une image nette de l'objet.
\textbf{Observations :}
\begin{itemize}
    \item Une image de l'objet se forme sur l'écran.
    \item Cette image est \textbf{renversée} (haut en bas, gauche à droite) par rapport à l'objet.
    \item L'image est \textbf{réelle} (elle se forme sur un écran) et \textbf{rétrécie}.
\end{itemize}
\textbf{Interprétation :} L'œil fonctionne comme une chambre noire. Le système optique (cornée + cristallin) agit comme une lentille convergente. Les rayons lumineux provenant d'un point de l'objet sont réfractés (déviés) par la cornée puis par le cristallin pour converger vers un point unique sur la rétine.
\textbf{Conclusion :} L'image d'un objet se forme sur la rétine, elle est renversée, réelle et rétrécie.
\end{experience}

\begin{propriete}[Formation de l'image sur la rétine]
L'image d'un objet se forme \textbf{renversée} sur la rétine, comme sur l'écran d'une chambre noire. Le système optique de l'œil (cornée + cristallin) se comporte comme une \textbf{lentille convergente} qui fait converger les rayons lumineux vers la rétine.
\end{propriete}

\begin{popfigure}
\begin{tikzpicture}[scale=1.1, >=Stealth]
% Coordonnées clés
\def\corneaX{0}
\def\pupilX{0.5}
\def\lensCenterX{1.8}
\def\retinaX{5.5}
\def\objX{-4}
\def\objH{1.5}
\def\imgH{-1.0} % Inverted, smaller

% Object (Arrow up)
\draw[popblue, very thick] (\objX,0) -- (\objX,\objH) node[above, black] {Objet};
\draw[popblue, very thick] (\objX,\objH) -- (\objX-0.2, \objH-0.3) (\objX,\objH) -- (\objX+0.2, \objH-0.3);

% Eye structure outlines
% Cornea (bulge)
\draw[popdark, thick] (\corneaX, -1.8) .. controls (\corneaX+0.5, -1.8) and (\corneaX+0.5, 1.8) .. (\corneaX, 1.8);
% Sclera / Wall
\draw[popdark, thick] (\corneaX, 1.8) -- (\retinaX, 1.5);
\draw[popdark, thick] (\corneaX, -1.8) -- (\retinaX, -1.5);
% Retina
\draw[poporange, very thick] (\retinaX, -1.5) -- (\retinaX, 1.5) node[right, black, align=left] {Rétine\\(Photorécepteurs)};

% Lens (Crystalline) - Biconvex shape
\draw[poppurple, thick] (\lensCenterX, -0.8) .. controls (\lensCenterX+0.3, -0.8) and (\lensCenterX+0.3, 0.8) .. (\lensCenterX, 0.8);
\draw[poppurple, thick] (\lensCenterX, -0.8) .. controls (\lensCenterX-0.3, -0.8) and (\lensCenterX-0.3, 0.8) .. (\lensCenterX, 0.8);
\node[poppurple] at (\lensCenterX, -1.3) {Cristallin};

% Iris / Pupil
\draw[popgreen, thick] (\pupilX, -0.4) -- (\pupilX, 0.4);
\node[popgreen] at (\pupilX, 0.9) {Iris / Pupille};

% Cornea label
\node[popblue] at (\corneaX, 2.2) {Cornée};

% Optic Nerve
\draw[popdark, thick, ->] (\retinaX, 0) -- (\retinaX+1, 0) node[right] {Nerf optique $\to$ Cerveau};

% Principal Rays (from top of object to bottom of image)
% Ray 1: Parallel to axis -> refracts through focal point (approx center of lens for thin lens approx)
% We draw rays hitting the lens center (optical center) -> straight line
% Ray from top of object through optical center of lens system (approx lensCenterX)
\draw[poppink, dashed, thick] (\objX, \objH) -- (\lensCenterX, 0) -- (\retinaX, \imgH);
% Ray from top of object parallel to axis -> refracts through focal point (simplified: crosses axis at focal point)
% For simplicity, draw ray hitting cornea top -> refract to image point
\draw[poppink, thick] (\objX, \objH) -- (\corneaX, 1.2) -- (\retinaX, \imgH);
% Ray from bottom of object (tip of arrow base) through center
\draw[poppink, dashed, thick] (\objX, 0) -- (\lensCenterX, 0) -- (\retinaX, 0);
% Ray from bottom parallel
\draw[poppink, thick] (\objX, 0) -- (\corneaX, 0) -- (\retinaX, 0);

% Image (Arrow down on retina)
\draw[poporange, very thick] (\retinaX, 0) -- (\retinaX, \imgH) node[below, black] {Image renversée};
\draw[poporange, very thick] (\retinaX, \imgH) -- (\retinaX-0.15, \imgH+0.2) (\retinaX, \imgH) -- (\retinaX+0.15, \imgH+0.2);

% Labels for media
\node[align=center, text width=2cm] at (0.8, -2.2) {Humeur\\ aqueuse};
\node[align=center, text width=2cm] at (3.5, -2.2) {Humeur\\ vitrée};

% Optical Axis
\draw[popmuted, dashed] (\objX-0.5, 0) -- (\retinaX+1.5, 0);

% Focal points (approximate)
\node[popmuted, font=\small] at (\lensCenterX-1.2, -0.3) {F};
\node[popmuted, font=\small] at (\lensCenterX+1.2, -0.3) {F'};

\end{tikzpicture}
\legende{Schéma du trajet des rayons lumineux dans l'œil. Le système cornée-cristallin (lentille convergente) forme une image réelle, renversée et rétrécie sur la rétine. Le nerf optique transmet le message nerveux au cortex visuel qui redresse l'image.}
\end{popfigure}

\begin{methode}[Interpréter un schéma de trajet lumineux dans l'œil]
\begin{enumerate}
    \item \textbf{Identifier le système optique :} Repérer la cornée et le cristallin (lentille convergente équivalente).
    \item \textbf{Repérer l'objet :} Noter sa position (loin ou près) et sa taille (flèche orientée vers le haut).
    \item \textbf{Tracer les rayons caractéristiques :}
    \begin{itemize}
        \item Rayon passant par le centre optique (non dévié).
        \item Rayon incident parallèle à l'axe optique (réfracté vers le foyer image F').
    \end{itemize}
    \item \textbf{Localiser l'image :} L'intersection des rayons réfractés détermine la position de l'image sur la rétine.
    \item \textbf{Caractériser l'image :} Vérifier qu'elle est \textbf{renversée} (flèche vers le bas), \textbf{réelle} (sur la rétine) et \textbf{rétrécie}.
\end{enumerate}
\end{methode}

\begin{attention}[Piège fréquent : L'iris ne fait pas la mise au point !]
Beaucoup d'élèves confondent le rôle de l'iris et celui du cristallin.
\begin{itemize}
    \item L'\textbf{iris} (muscles circulaire et radial) modifie le \textbf{diamètre de la pupille} pour réguler la \textbf{quantité de lumière} entrant dans l'œil (réflexe pupillaire). C'est un diaphragme.
    \item Le \textbf{cristallin} modifie sa \textbf{courbure} pour faire converger les rayons exactement sur la rétine. C'est la \textbf{mise au point} (accommodation).
\end{itemize}
\cle{Ne jamais écrire : « L'iris accommode » ou « La pupille fait la mise au point ».}
\end{attention}

\courssub{Le cristallin : une lentille convergente déformable}

Contrairement à une lentille en verre dont la forme est figée, le cristallin est une lentille \textbf{vivante, transparente et élastique}. Sa courbure peut varier grâce à l'action des muscles ciliaires. C'est cette propriété unique qui permet à l'œil de voir net aussi bien un livre à 25 cm qu'une montagne à plusieurs kilomètres.

\begin{definition}[Accommodation]
L'\textbf{accommodation} est la modification de la courbure du cristallin (et donc de sa puissance de convergence) permettant la formation d'une image nette sur la rétine pour des objets situés à \textbf{différentes distances}.
\end{definition}

\courssub{Mécanisme de l'accommodation : vision de loin et vision de près}

Le mécanisme repose sur l'antagonisme entre l'élasticité propre du cristallin (qui tend à se bomber) et la traction exercée par les ligaments suspenseurs (zonule de Zinn) tirés par l'étirement du muscle ciliaire.

\begin{experience}[Mise en évidence de l'accommodation (Expérience subjective)]
\textbf{Matériel :} Un crayon (ou le doigt), un objet lointain (fenêtre, arbre).
\textbf{Protocole :}
\begin{enumerate}
    \item Tends le bras, fixe le bout du crayon (vision de près). Observe l'arrière-plan : il est flou.
    \item Sans bouger le crayon, fixe l'objet lointain (vision de loin). Observe le crayon : il devient flou (et souvent double).
    \item Alterne la fixation : près $\leftrightarrow$ loin. Tu peux ressentir une légère tension musculaire interne à l'œil lors du passage près $\to$ loin.
\end{enumerate}
\textbf{Interprétation :} L'œil ne peut pas voir net simultanément un objet proche et un objet lointain. Il doit \textbf{changer la puissance de convergence de son cristallin} pour déplacer le foyer sur la rétine selon la distance de l'objet.
\textbf{Conclusion :} L'accommodation est un acte actif (vision de près) ou passif (vision de loin) modifiant la forme du cristallin.
\end{experience}

\begin{popfigure}
\begin{tikzpicture}[scale=1.0, >=Stealth]
% Common style
\tikzset{
    eyeWall/.style={popdark, thick},
    lensRelaxed/.style={poppurple!60, thick},
    lensAccom/.style={poppurple, very thick},
    ligament/.style={popmuted, thin, densely dashed},
    muscle/.style={popred, thick},
    ray/.style={poppink, thick},
    label/.style={font=\small, popdark}
}

% ---------- OEIL AU REPOS (VISION DE LOIN) ----------
\begin{scope}[xshift=0cm]
% Coordinates
\def\cx{0} % Center of lens approx
\def\retX{4.5}
\def\cornX{-1.5}
\def\muscY{1.6}
\def\ligY{0.7}

% Sclera / Cornea
\draw[eyeWall] (\cornX, -1.5) .. controls (\cornX+0.5, -1.5) and (\cornX+0.5, 1.5) .. (\cornX, 1.5);
\draw[eyeWall] (\cornX, 1.5) -- (\retX, 1.2);
\draw[eyeWall] (\cornX, -1.5) -- (\retX, -1.2);
% Retina
\draw[poporange, very thick] (\retX, -1.2) -- (\retX, 1.2) node[right, label] {Rétine};

% Ciliary Muscle (Relaxed = stretched flat)
\draw[muscle] (\cornX+0.8, \muscY) .. controls (\cx+1.5, \muscY) and (\cx+1.5, -\muscY) .. (\cornX+0.8, -\muscY);
\node[label, popred, align=center] at (\cx+1.5, \muscY+0.4){Muscle ciliaire\\(relâché/étiré)};

% Ligaments (Taut)
\draw[ligament] (\cx, \ligY) -- (\cornX+0.8, \muscY);
\draw[ligament] (\cx, -\ligY) -- (\cornX+0.8, -\muscY);
\draw[ligament] (\cx, 0.9) -- (\cornX+0.8, \muscY-0.1);
\draw[ligament] (\cx, -0.9) -- (\cornX+0.8, -\muscY+0.1);
\node[label, popmuted, align=center] at (\cx+0.5, \ligY+0.5){Ligaments\\(tendus)};

% Lens (Flat)
\draw[lensRelaxed] (\cx, -0.5) .. controls (\cx+0.2, -0.5) and (\cx+0.2, 0.5) .. (\cx, 0.5);
\draw[lensRelaxed] (\cx, -0.5) .. controls (\cx-0.2, -0.5) and (\cx-0.2, 0.5) .. (\cx, 0.5);
\node[label, poppurple, align=center] at (\cx, -1.0){Cristallin\\(aplati)};

% Rays from infinity (Parallel)
\draw[ray] (-3, 1.0) -- (\cornX, 1.0) -- (\retX, 0.2);
\draw[ray] (-3, -1.0) -- (\cornX, -1.0) -- (\retX, -0.2);
\draw[ray] (-3, 0) -- (\cornX, 0) -- (\retX, 0);
\node[label, align=center] at (-3.5, 0){Objet lointain\\ (rayons $\parallel$)};

% Title
\node[font=\bfseries, popblue] at (\cx, 2.2) {\textbf{Œil au repos : Vision de loin}};
\end{scope}

% ---------- OEIL ACCOMMODE (VISION DE PRES) ----------
\begin{scope}[xshift=11cm]
% Coordinates
\def\cx{0}
\def\retX{4.5}
\def\cornX{-1.5}
\def\muscY{1.2} % Muscle contracts -> moves forward/inward
\def\ligY{0.7}

% Sclera / Cornea
\draw[eyeWall] (\cornX, -1.5) .. controls (\cornX+0.5, -1.5) and (\cornX+0.5, 1.5) .. (\cornX, 1.5);
\draw[eyeWall] (\cornX, 1.5) -- (\retX, 1.2);
\draw[eyeWall] (\cornX, -1.5) -- (\retX, -1.2);
% Retina
\draw[poporange, very thick] (\retX, -1.2) -- (\retX, 1.2) node[right, label] {Rétine};

% Ciliary Muscle (Contracted = rounded, moves forward)
\draw[muscle] (\cornX+0.5, \muscY) .. controls (\cx+1.0, \muscY+0.3) and (\cx+1.0, -\muscY-0.3) .. (\cornX+0.5, -\muscY);
\node[label, popred, align=center] at (\cx+1.0, \muscY+0.6){Muscle ciliaire\\(contracté)};

% Ligaments (Slack)
\draw[ligament] (\cx, \ligY) .. controls (\cx+0.5, \ligY) .. (\cornX+0.5, \muscY);
\draw[ligament] (\cx, -\ligY) .. controls (\cx+0.5, -\ligY) .. (\cornX+0.5, -\muscY);
\draw[ligament] (\cx, 0.9) .. controls (\cx+0.5, 0.9) .. (\cornX+0.5, \muscY-0.1);
\draw[ligament] (\cx, -0.9) .. controls (\cx+0.5, -0.9) .. (\cornX+0.5, -\muscY+0.1);
\node[label, popmuted, align=center] at (\cx+0.5, \ligY+0.5){Ligaments\\(relâchés)};

% Lens (Bulged / Round)
\draw[lensAccom] (\cx, -1.0) .. controls (\cx+0.4, -1.0) and (\cx+0.4, 1.0) .. (\cx, 1.0);
\draw[lensAccom] (\cx, -1.0) .. controls (\cx-0.4, -1.0) and (\cx-0.4, 1.0) .. (\cx, 1.0);
\node[label, poppurple, align=center] at (\cx, -1.4){Cristallin\\(bombé)};

% Rays from near object (Diverging)
\draw[ray] (-2, 1.5) -- (\cornX, 1.0) -- (\retX, 0.2);
\draw[ray] (-2, -1.5) -- (\cornX, -1.0) -- (\retX, -0.2);
\draw[ray] (-2, 0) -- (\cornX, 0) -- (\retX, 0);
\node[label, align=center] at (-2.5, 0){Objet proche\\ (rayons divergents)};

% Title
\node[font=\bfseries, popblue] at (\cx, 2.2) {\textbf{Œil accommodé : Vision de près}};
\end{scope}
\end{tikzpicture}
\legende{Comparaison du mécanisme d'accommodation. \textbf{Gauche (Vision de loin) :} Muscle ciliaire relâché $\to$ ligaments tendus $\to$ cristallin aplati (faible convergence). \textbf{Droite (Vision de près) :} Muscle ciliaire contracté $\to$ ligaments relâchés $\to$ élasticité du cristallin $\to$ cristallin bombé (forte convergence).}
\end{popfigure}

\begin{propriete}[Mécanisme de l'accommodation]
\begin{itemize}
    \item \textbf{Vision de loin (œil au repos) :} Le muscle ciliaire est \textbf{relâché} (étiré par la pression intraoculaire). Les ligaments suspenseurs sont \textbf{tendus}. Ils tirent le cristallin vers l'extérieur $\to$ le cristallin s'\textbf{aplatit} (faible convergence). L'image d'un objet lointain (rayons parallèles) se forme sur la rétine.
    \item \textbf{Vision de près (œil accommodé) :} Le muscle ciliaire se \textbf{contracte} (se raccourcit, se bombe vers l'avant). Les ligaments suspenseurs se \textbf{relâchent}. Le cristallin, par sa propre \textbf{élasticité}, se \textbf{bombe} (forte convergence). L'image d'un objet proche (rayons divergents) se forme sur la rétine.
\end{itemize}
\end{propriete}

\begin{aretenir}[Sens de l'accommodation]
L'accommodation est un \textbf{effort} (contraction musculaire) pour voir \textbf{de près}. Au repos, l'œil est fait pour voir \textbf{de loin}. C'est pourquoi une lecture prolongée fatigue les muscles ciliaires (asthénopie).
\end{aretenir}

\courssub{Le réflexe pupillaire : adaptation de la quantité de lumière}

Si le cristallin gère la \textbf{netteté} (mise au point), l'iris gère l'\textbf{éclairement} (quantité de lumière atteignant la rétine). C'est un réflexe autonome, involontaire et consensuel (les deux pupilles réagissent ensemble).

\begin{experience}[Mise en évidence du réflexe pupillaire]
\textbf{Matériel :} Une lampe de poche, un miroir (ou travailler à deux).
\textbf{Protocole :}
\begin{enumerate}
    \item Place-toi dans une pièce faiblement éclairée. Observe la taille de ta pupille dans le miroir (elle est \textbf{dilatée}, grande).
    \item Allume la lampe de poche et dirige le faisceau vers un œil (ou les deux).
    \item Observe la réaction immédiate : la pupille se \textbf{rétrécit} (myosis).
    \item Éteins la lampe : la pupille se \textbf{dilate} à nouveau (mydriase).
\end{enumerate}
\textbf{Interprétation :} L'iris contient deux muscles antagonistes :
\begin{itemize}
    \item Muscle \textbf{sphincter} de la pupille (circulaire) $\to$ contraction = \textbf{rétrécissement} (lumière vive).
    \item Muscle \textbf{dilatateur} de la pupille (radial) $\to$ contraction = \textbf{dilatation} (obscurité).
\end{itemize}
C'est un \textbf{réflexe photomoteur} : Récepteur (rétine) $\to$ Nerf optique $\to$ Centre nerveux (tronc cérébral) $\to$ Nerf oculomoteur (parasympathique) $\to$ Effecteur (muscle sphincter de l'iris).
\end{experience}

\begin{definition}[Réflexe pupillaire (photomoteur)]
C'est la variation involontaire du diamètre de la pupille en fonction de l'intensité lumineuse : \textbf{myosis} (rétrécissement) en lumière vive pour protéger la rétine et augmenter la profondeur de champ ; \textbf{mydriase} (dilatation) dans l'obscurité pour capter un maximum de photons.
\end{definition}

\courssub{De la rétine au cerveau : photorécepteurs et traitement de l'image}

L'image renversée formée sur la rétine n'est qu'une distribution lumineuse. Elle doit être convertie en signal nerveux.

\begin{definition}[Photorécepteurs : Bâtonnets et Cônes]
La rétine contient deux types de cellules photosensibles (photorécepteurs) qui transforment l'énergie lumineuse en influx nerveux (transduction) :
\begin{center}
\begin{tabularx}{\linewidth}{|l|X|X|}\hline
\rowcolor{popblueL}
\textbf{Caractère} & \textbf{Bâtonnets ($\sim$ 120 millions)} & \textbf{Cônes ($\sim$ 6 millions)} \\ \hline
\textbf{Vision} & Crépusculaire (scotopique), noir et blanc & Diurne (photopique), \textbf{couleurs}, détails fins \\ \hline
\textbf{Seuil de sensibilité} & Très bas (quelques photons) & Élevé (lumière intense nécessaire) \\ \hline
\textbf{Acuité visuelle} & Faible (vision floue, périphérique) & \textbf{Maximale} (fovéa, vision centrale) \\ \hline
\textbf{Localisation} & Périphérie de la rétine (absents de la fovéa) & Concentrés à la \textbf{fovéa} (tache jaune) \\ \hline
\textbf{Pigment} & Rhodopsine (pourpre rétinien) & Iodopsines (3 types : Bleu, Vert, Rouge) \\ \hline
\end{tabularx}
\end{center}
\end{definition}

\begin{savaistu}[Pourquoi l'image est-elle nette seulement au centre ?]
Seule la \textbf{fovéa} (dépression centrale de la macula) ne contient \textbf{que des cônes}, très serrés, chacun relié à une fibre du nerf optique (convergence 1:1). C'est la zone de \textbf{vision maximale}. En périphérie, la convergence bâtonnets $\to$ fibre nerveuse est forte (ex: 100:1), ce qui augmente la sensibilité mais baisse la résolution. C'est pourquoi tu dois fixer un mot pour le lire net, mais tu détectes le mouvement sur les côtés.
\end{savaistu}

\begin{propriete}[Du message nerveux à la perception visuelle]
\begin{enumerate}
    \item \textbf{Transduction :} La lumière modifie le pigment (rhéodopsine/iodopsine) $\to$ déclenchement d'un influx nerveux dans le photorécepteur.
    \item \textbf{Traitement rétinien :} Les cellules bipolaires et ganglionnaires intègrent les signaux (contraste, contours, mouvement).
    \item \textbf{Transmission :} Les axones des cellules ganglionnaires forment le \textbf{nerf optique} (II). Ils sortent de l'œil au \textbf{point aveugle} (papille optique, sans photorécepteurs).
    \item \textbf{Chiasma optique :} Les fibres nasales (vision temporale) croisent ; les fibres temporales (vision nasale) ne croisent pas. $\to$ Chaque hémisphère reçoit l'information du champ visuel controlatéral.
    \item \textbf{Cortex visuel (aire 17, lobe occipital) :} Reconstruction de l'image, \textbf{redressement} de l'image renversée, reconnaissance des formes, couleurs, profondeur, mouvement.
\end{enumerate}
\end{propriete}

\begin{attention}[Point aveugle vs Fovéa]
\begin{itemize}
    \item \textbf{Fovéa :} Zone de \textbf{meilleure acuité} (cônes seulement), pas de vaisseaux sanguins en surface. C'est où l'image se forme pour la vision précise.
    \item \textbf{Point aveugle (Papille) :} Sortie du nerf optique + entrée/sortie vaisseaux. \textbf{Zéro photorécepteur} $\to$ \textbf{aucune vision}. Le cerveau « comble » le trou (remplissage perceptuel).
\end{itemize}
\cle{Ne pas confondre : on \textbf{fixe} avec la fovéa ; on \textbf{ne voit pas} au point aveugle.}
\end{attention}

\courssub{Exercice résolu : Analyse d'un défaut d'accommodation}

\begin{exoresolu}[Analyse de la presbytie naissante]
\textbf{Énoncé :} Monsieur K., 45 ans, instituteur à Yaoundé, se plaint de devoir « allonger les bras » pour lire son journal le matin. Il voit parfaitement net le tableau noir à 5 m. Son ophtalmologue lui prescrit des verres convergents de +1,50 dioptries pour la lecture de près.
\begin{enumerate}
    \item Rappeler le mécanisme normal de l'accommodation pour la vision de près.
    \item Expliquer pourquoi M. K. doit éloigner le journal pour le lire net.
    \item Justifier la correction par des verres convergents.
\end{enumerate}
\tcblower
\textbf{Solution détaillée :}
\begin{enumerate}
    \item \textbf{Mécanisme normal (vision de près) :} Pour voir net un objet proche (ex: 25 cm), les rayons lumineux arrivent \textbf{divergents} sur l'œil. Le muscle ciliaire se \textbf{contracte}, les ligaments suspenseurs se \textbf{relâchent}, le cristallin \textbf{se bombe} (augmentation de sa puissance de convergence). La convergence accrue permet de faire converger les rayons divergents précisément sur la rétine.
    \item \textbf{Cause de la gêne (Presbytie) :} Avec l'âge (dès 40-45 ans), le cristallin \textbf{perd de son élasticité} (durcit, se sclérose). Même si le muscle ciliaire se contracte normalement, le cristallin ne peut plus se bomber suffisamment. Sa \textbf{puissance de convergence maximale diminue}. Le point rapproché (distance minimale de vision nette) \textbf{recule} (ex: passe de 25 cm à 50 cm ou plus). M. K. éloigne le journal pour placer l'objet \textbf{au-delà de son point rapproché}, là où les rayons sont moins divergents, compatibles avec la convergence résiduelle de son cristallin durci.
    \item \textbf{Justification de la correction :} Les verres convergents (puissance +1,50 D) \textbf{convergent au préalable} les rayons lumineux avant qu'ils n'entrent dans l'œil. Ils transforment les rayons divergents (provenant de l'objet à 30-40 cm) en rayons \textbf{moins divergents} (ou parallèles si l'objet est au foyer du verre). L'œil presbyte, qui a une convergence insuffisante pour le près, reçoit ainsi des rayons qu'il peut converger sur la rétine sans effort d'accommodation (ou avec un effort résiduel faible). Le verre \textbf{compense le manque de convergence} du cristallin vieillissant.
\end{enumerate}
\end{exoresolu}

\begin{exercice}[Application : Trajet lumineux et accommodation]
Un observateur regarde un oiseau perché sur une branche à 10 m (Situation A), puis regarde une fourmi sur sa main à 20 cm (Situation B).
\begin{enumerate}
    \item Pour chaque situation, décrire l'état du muscle ciliaire, des ligaments suspenseurs et la forme du cristallin.
    \item Tracer schématiquement le trajet de deux rayons lumineux émis par le sommet de l'objet (oiseau puis fourmi) jusqu'à la rétine en précisant la nature de l'image formée.
    \item Expliquer pourquoi la vision de l'oiseau devient floue si l'observateur continue de fixer la fourmi (accommodation maintenue pour le près).
\end{enumerate}
\end{exercice}

\begin{corrige}[Exercice : Trajet lumineux et accommodation]
\begin{enumerate}
    \item \textbf{Situation A (Loin, 10 m) :} Muscle ciliaire \textbf{relâché} $\to$ Ligaments \textbf{tendus} $\to$ Cristallin \textbf{aplati} (faible convergence). Rayons incidents $\approx$ \textbf{parallèles}. Image sur rétine : \textbf{renversée, réelle, rétrécie}.
    \item \textbf{Situation B (Près, 20 cm) :} Muscle ciliaire \textbf{contracté} $\to$ Ligaments \textbf{relâchés} $\to$ Cristallin \textbf{bombé} (forte convergence). Rayons incidents \textbf{divergents}. Image sur rétine : \textbf{renversée, réelle, rétrécie}.
    \item Si l'œil reste accommodé pour le près (cristallin bombé, forte convergence) alors que l'objet est loin (rayons parallèles), la convergence est \textbf{trop forte}. Les rayons convergent \textbf{en avant de la rétine} (myopie artificielle temporaire). L'image sur la rétine est une \textbf{tache floue} (cercle de diffusion). Il faut relâcher l'accommodation (muscle ciliaire au repos) pour aplatir le cristallin et ramener le foyer sur la rétine.
\end{enumerate}
\end{corrige}

\begin{savaistu}[L'accommodation chez les animaux]
Tous les vertébrés n'accommodent pas de la même manière !
\begin{itemize}
    \item \textbf{Mammifères (dont l'homme) et Oiseaux :} Accommodation par \textbf{déformation du cristallin} (changement de courbure). Rapide, précis.
    \item \textbf{Poissons et Amphibiens :} Cristallin sphérique et dur (haute convergence fixe pour l'eau). Accommodation par \textbf{déplacement du cristallin} (avant/arrière) grâce à un muscle rétracteur (poissons) ou protracteur (amphibiens). Plus lent.
    \item \textbf{Serpents :} Combinaison des deux (déformation + déplacement).
\end{itemize}
C'est une adaptation au milieu de vie (air vs eau) : l'eau a un indice de réfraction proche de la cornée ($n \approx 1,33$), la cornée ne converge presque plus ; tout le pouvoir optique repose sur le cristallin sphérique.
\end{savaistu}

\coursec{Anomalies de la vision : myopie, hypermétropie, presbytie}

Dans la section précédente, nous avons vu comment l'œil \textbf{emmétrope} (normal) forme une image nette d'un objet situé à l'infini exactement sur la rétine, sans effort d'accommodation, et comment l'accommodation permet de voir net de près en augmentant la convergence du cristallin. Mais la nature n'est pas toujours parfaite : la forme du globe oculaire ou le vieillissement du cristallin peuvent perturber cette mise au point idéale. Nous allons étudier les trois principales \textbf{anomalies de la réfraction} (amétropies) : la myopie, l'hypermétropie et la presbytie, en comprenant leur origine optique, leurs symptômes et leur correction par des verres ophtalmiques.

\courssub{Rappel : l'œil emmétrope, référence de la vision normale}

\begin{definition}[Œil emmétrope]
Un œil est dit \cle{emmétrope} lorsque, au repos (accommodation nulle), l'image d'un objet situé à l'infini se forme \textbf{exactement sur la rétine}. La puissance optique totale de l'œil (cornée + cristallin) et la longueur axiale du globe oculaire sont en parfait équilibre.
\end{definition}

\begin{popfigure}
\begin{tikzpicture}[scale=0.9, >=Stealth]
% Eye globe
\draw[popdark, thick] (0,0) circle (2.5cm);
% Cornea bulge
\draw[popblue, thick] (-2.5,0) .. controls (-3,0.5) and (-3,1.5) .. (-2.5,2) .. controls (-2,2.5) and (-1.5,2.5) .. (-0.5,2) .. controls (0.5,2) and (1.5,2.5) .. (2.5,2) .. controls (3,1.5) and (3,0.5) .. (2.5,0);
% Lens
\draw[poppurple, thick] (-0.5,-0.8) .. controls (-0.5,-1.2) and (0.5,-1.2) .. (0.5,-0.8) .. controls (0.5,-0.4) and (-0.5,-0.4) .. (-0.5,-0.8);
% Retina
\draw[popgreen, very thick] (1.5, -2) arc (0:180:2.5cm and 2cm);
% Optic nerve
\draw[popmuted, thick] (0,-2.5) -- (0,-3.5);
% Object at infinity (parallel rays)
\draw[poporange, thick, ->] (-5, 1.5) -- (-2.5, 1.5);
\draw[poporange, thick, ->] (-5, 0.5) -- (-2.5, 0.5);
\draw[poporange, thick, ->] (-5, -0.5) -- (-2.5, -0.5);
% Refraction at cornea/lens (simplified convergence)
\draw[poporange, thick] (-2.5, 1.5) -- (0, 0);
\draw[poporange, thick] (-2.5, 0.5) -- (0, 0);
\draw[poporange, thick] (-2.5, -0.5) -- (0, 0);
% Focus on retina
\fill[poporange] (1.8, 0) circle (2pt);
% Labels
\node[popblue] at (-3.5, 2.5) {Cornée};
\node[poppurple] at (-1.5, -1.5) {Cristallin};
\node[popgreen] at (3.5, 0) {Rétine};
\node[popmuted] at (0, -4) {Nerf optique};
\node[poporange] at (-5.5, 1.5) {Objet $\infty$};
\node[popdark, align=center] at (0, 3.2) {\textbf{Œil emmétrope}\\(Image nette sur la rétine)};
\end{tikzpicture}
\legende{Schéma simplifié de l'œil emmétrope : les rayons parallèles provenant d'un objet à l'infini convergent exactement sur la rétine après réfraction par la cornée et le cristallin.}
\end{popfigure}

\courssub{La myopie : l'œil « trop long » ou « trop convergent »}

\begin{experience}[Observation clinique : le myope et le tableau noir]
Un élève de 3ème, Assan, se plaint de ne pas voir distinctement les écritures au tableau noir placé à \SI{5}{m}, alors qu'il lit son cahier de près sans difficulté. L'ophtalmologiste mesure sa \cle{réfraction} et prescrit des verres de \textbf{-2,50 dioptries}.
\par \textbf{Interprétation :} Assan est myope. Son œil forme l'image d'un objet lointain \textbf{en avant de la rétine}. De près, l'objet étant plus proche, ses rayons divergent davantage à l'entrée de l'œil ; le point de convergence recule et peut coïncider avec la rétine sans accommodation (ou avec un faible effort).
\end{experience}

\begin{definition}[Myopie]
La \cle{myopie} est une anomalie de la réfraction où l'image d'un objet éloigné (à l'infini) se forme \textbf{en avant de la rétine}. La vision de loin est floue, la vision de près reste nette (jusqu'à un certain point, le \cle{point lointain}).
\end{definition}

\begin{propriete}[Causes optiques de la myopie]
Deux causes anatomiques principales, souvent combinées :
\begin{enumerate}
    \item \textbf{Axiale (la plus fréquente) :} Le globe oculaire est \textbf{trop long} (longueur axiale $> \SI{24}{mm}$ environ). La rétine est « trop loin » du plan focal du système optique.
    \item \textbf{Indicielle / Cornéenne :} Le système optique (cornée trop bombée ou cristallin trop puissant) est \textbf{trop convergent}. La focale est trop courte pour la longueur de l'œil.
\end{enumerate}
\end{propriete}

\begin{popfigure}
\begin{tikzpicture}[scale=0.9, >=Stealth]
% Myopic Eye (Uncorrected)
\begin{scope}[xshift=-7cm]
\draw[popdark, thick] (0,0) circle (2.5cm);
\draw[popblue, thick] (-2.5,0) .. controls (-3,0.5) and (-3,1.5) .. (-2.5,2) .. controls (-2,2.5) and (-1.5,2.5) .. (-0.5,2) .. controls (0.5,2) and (1.5,2.5) .. (2.5,2) .. controls (3,1.5) and (3,0.5) .. (2.5,0);
\draw[poppurple, thick] (-0.5,-0.8) .. controls (-0.5,-1.2) and (0.5,-1.2) .. (0.5,-0.8) .. controls (0.5,-0.4) and (-0.5,-0.4) .. (-0.5,-0.8);
\draw[popgreen, very thick] (1.5, -2) arc (0:180:2.5cm and 2cm);
\draw[popmuted, thick] (0,-2.5) -- (0,-3.5);
% Rays
\draw[poporange, thick, ->] (-5, 1.5) -- (-2.5, 1.5);
\draw[poporange, thick, ->] (-5, 0.5) -- (-2.5, 0.5);
\draw[poporange, thick, ->] (-5, -0.5) -- (-2.5, -0.5);
% Convergence BEFORE retina
\draw[poporange, thick] (-2.5, 1.5) -- (0.5, 0);
\draw[poporange, thick] (-2.5, 0.5) -- (0.5, 0);
\draw[poporange, thick] (-2.5, -0.5) -- (0.5, 0);
\fill[poporange] (0.5, 0) circle (2pt) node[above right] {Foyer (avant rétine)};
% Divergence after focus
\draw[poporange, dashed, thin] (0.5, 0) -- (1.8, -0.5);
\draw[poporange, dashed, thin] (0.5, 0) -- (1.8, 0.5);
\node[popdark, align=center] at (0, 3.2) {\textbf{Myopie (non corrigée)}\\(Image en avant de la rétine)};
\end{scope}

% Myopic Eye (Corrected)
\begin{scope}[xshift=7cm]
% Diverging lens
\draw[popblue!70!black, very thick] (-3.5, -1.5) -- (-3.5, 1.5);
\draw[popblue!70!black, thick] (-3.5, -1.5) -- (-3.7, -1.5) (-3.5, 1.5) -- (-3.7, 1.5);
\draw[popblue!70!black, thick] (-3.5, -1.5) -- (-3.3, -1.5) (-3.5, 1.5) -- (-3.3, 1.5);
\node[popblue!70!black] at (-3.5, 2) {Verre divergent};
% Eye
\draw[popdark, thick] (0,0) circle (2.5cm);
\draw[popblue, thick] (-2.5,0) .. controls (-3,0.5) and (-3,1.5) .. (-2.5,2) .. controls (-2,2.5) and (-1.5,2.5) .. (-0.5,2) .. controls (0.5,2) and (1.5,2.5) .. (2.5,2) .. controls (3,1.5) and (3,0.5) .. (2.5,0);
\draw[poppurple, thick] (-0.5,-0.8) .. controls (-0.5,-1.2) and (0.5,-1.2) .. (0.5,-0.8) .. controls (0.5,-0.4) and (-0.5,-0.4) .. (-0.5,-0.8);
\draw[popgreen, very thick] (1.5, -2) arc (0:180:2.5cm and 2cm);
\draw[popmuted, thick] (0,-2.5) -- (0,-3.5);
% Rays from infinity
\draw[poporange, thick, ->] (-6, 1.5) -- (-3.5, 1.5);
\draw[poporange, thick, ->] (-6, 0.5) -- (-3.5, 0.5);
\draw[poporange, thick, ->] (-6, -0.5) -- (-3.5, -0.5);
% Divergence by lens
\draw[poporange, thick] (-3.5, 1.5) -- (-2.5, 2.0);
\draw[poporange, thick] (-3.5, 0.5) -- (-2.5, 0.8);
\draw[poporange, thick] (-3.5, -0.5) -- (-2.5, -0.4);
% Convergence by eye ON retina
\draw[poporange, thick] (-2.5, 2.0) -- (1.8, 0);
\draw[poporange, thick] (-2.5, 0.8) -- (1.8, 0);
\draw[poporange, thick] (-2.5, -0.4) -- (1.8, 0);
\fill[poporange] (1.8, 0) circle (2pt) node[above right] {Image nette};
\node[popdark, align=center] at (0, 3.2) {\textbf{Myopie corrigée}\\(Verre divergent ramène l'image sur la rétine)};
\end{scope}
\end{tikzpicture}
\legende{À gauche : œil myope non corrigé, l'image se forme en avant de la rétine (vision floue de loin). À droite : correction par un \textbf{verre divergent} (concave, puissance négative en dioptries) qui diverge les rayons incidents avant qu'ils n'entrent dans l'œil, reculant le foyer sur la rétine.}
\end{popfigure}

\begin{attention}[Piège classique : « Le myope voit mal de près »]
\textbf{FAUX.} C'est l'erreur la plus fréquente en examen.
\begin{itemize}
    \item Le myope \textbf{voit bien de près} (son point proche est souvent plus proche que la normale, ex : \SI{10}{cm} au lieu de \SI{25}{cm}).
    \item Le myope \textbf{voit mal de loin} (son point lointain est à distance finie, ex : \SI{50}{cm}).
\end{itemize}
Rappelle-toi : \emph{Myope} $\rightarrow$ \textbf{M}al voyant de loin, \textbf{B}on voyant de près (Mnémotechnique : \textbf{My}ope = \textbf{My}ope de loin).
\end{attention}

\begin{exemplebox}[Exemple chiffré : Correction d'un myope à -2,00 D]
Un élève myope a une réfraction de \textbf{-2,00 dioptries (D)}.
\begin{itemize}
    \item Son \textbf{point lointain} (distance max où il voit net sans correction) est à $d = \frac{1}{|P|} = \frac{1}{2} = \mathbf{0,50\ m = 50\ cm}$. Au-delà de \SI{50}{cm}, tout est flou. Il ne peut pas lire le tableau à \SI{3}{m}.
    \item Le verre correcteur est un \textbf{verre divergent de -2,00 D}. Il crée une image virtuelle de l'objet lointain (à l'infini) au foyer de la lentille, c'est-à-dire à \SI{50}{cm} de l'œil. L'œil myope, qui voit net à \SI{50}{cm}, reçoit alors une image nette.
\end{itemize}
\end{exemplebox}

\courssub{L'hypermétropie : l'œil « trop court » ou « pas assez convergent »}

\begin{experience}[Observation clinique : l'hypermétrope et la lecture]
Marie, 12 ans, se plaint de maux de tête en fin de journée, de fatigue oculaire et de difficultés à lire longtemps. Elle voit bien le tableau. L'examen révèle une hypermétropie de \textbf{+1,50 D}.
\par \textbf{Interprétation :} L'œil de Marie est « trop court ». Pour voir net de loin, il doit \textbf{accommoder en permanence} (cristallin bombé) pour augmenter sa convergence et ramener l'image sur la rétine. De près, l'effort d'accommodation est encore plus grand (vision de loin + vision de près), d'où la fatigue (asthénopie) et les céphalées.
\end{experience}

\begin{definition}[Hypermétropie]
L'\cle{hypermétropie} est une anomalie où l'image d'un objet éloigné se forme \textbf{en arrière de la rétine} (théoriquement). Pour la ramener sur la rétine, l'œil doit fournir un effort d'accommodation permanent. La vision de loin peut être nette (si l'accommodation compense), mais la vision de près est fatigante, voire floue si l'amplitude d'accommodation est insuffisante.
\end{definition}

\begin{propriete}[Causes de l'hypermétropie]
\begin{enumerate}
    \item \textbf{Axiale (majoritaire) :} Globe oculaire \textbf{trop court} (longueur axiale $< \SI{23}{mm}$).
    \item \textbf{Indicielle :} Système optique \textbf{pas assez convergent} (cornée trop plate, cristallin pas assez puissant).
\end{enumerate}
\end{propriete}

\begin{popfigure}
\begin{tikzpicture}[scale=0.9, >=Stealth]
% Hypermetropic Eye (Uncorrected - Accommodating)
\begin{scope}[xshift=-7cm]
\draw[popdark, thick] (0,0) circle (2.5cm);
\draw[popblue, thick] (-2.5,0) .. controls (-3,0.5) and (-3,1.5) .. (-2.5,2) .. controls (-2,2.5) and (-1.5,2.5) .. (-0.5,2) .. controls (0.5,2) and (1.5,2.5) .. (2.5,2) .. controls (3,1.5) and (3,0.5) .. (2.5,0);
% Thick lens = accommodating
\draw[poppurple, very thick] (-0.5,-1) .. controls (-0.5,-1.5) and (0.5,-1.5) .. (0.5,-1) .. controls (0.5,-0.5) and (-0.5,-0.5) .. (-0.5,-1);
\draw[popgreen, very thick] (1.5, -2) arc (0:180:2.5cm and 2cm);
\draw[popmuted, thick] (0,-2.5) -- (0,-3.5);
% Rays from infinity
\draw[poporange, thick, ->] (-5, 1.5) -- (-2.5, 1.5);
\draw[poporange, thick, ->] (-5, 0.5) -- (-2.5, 0.5);
\draw[poporange, thick, ->] (-5, -0.5) -- (-2.5, -0.5);
% Convergence -> Focus BEHIND retina (virtual extension)
\draw[poporange, thick] (-2.5, 1.5) -- (0, 0);
\draw[poporange, thick] (-2.5, 0.5) -- (0, 0);
\draw[poporange, thick] (-2.5, -0.5) -- (0, 0);
% Extension lines to virtual focus behind
\draw[poporange, dashed, thin] (0, 0) -- (2.5, 0);
\fill[poporange] (2.5, 0) circle (2pt) node[right] {Foyer théorique (arrière rétine)};
% Blur circle on retina
\draw[poporange, fill=poporange!20] (1.5, -0.2) ellipse (0.3 and 0.15);
\node[popdark, align=center] at (0, 3.2) {\textbf{Hypermétropie (non corrigée)}\\(Accommodation forcée pour voir net)};
\node[poppurple] at (-1.5, -2) {Cristallin bombé (effort)};
\end{scope}

% Hypermetropic Eye (Corrected)
\begin{scope}[xshift=7cm]
% Converging lens
\draw[poporange!80!black, very thick] (-3.5, -1.5) -- (-3.5, 1.5);
\draw[poporange!80!black, thick] (-3.5, -1.5) .. controls (-3.2, -1.5) and (-3.2, 1.5) .. (-3.5, 1.5);
\draw[poporange!80!black, thick] (-3.5, -1.5) .. controls (-3.8, -1.5) and (-3.8, 1.5) .. (-3.5, 1.5);
\node[poporange!80!black] at (-3.5, 2) {Verre convergent};
% Eye (Relaxed)
\draw[popdark, thick] (0,0) circle (2.5cm);
\draw[popblue, thick] (-2.5,0) .. controls (-3,0.5) and (-3,1.5) .. (-2.5,2) .. controls (-2,2.5) and (-1.5,2.5) .. (-0.5,2) .. controls (0.5,2) and (1.5,2.5) .. (2.5,2) .. controls (3,1.5) and (3,0.5) .. (2.5,0);
\draw[poppurple, thick] (-0.5,-0.8) .. controls (-0.5,-1.2) and (0.5,-1.2) .. (0.5,-0.8) .. controls (0.5,-0.4) and (-0.5,-0.4) .. (-0.5,-0.8);
\draw[popgreen, very thick] (1.5, -2) arc (0:180:2.5cm and 2cm);
\draw[popmuted, thick] (0,-2.5) -- (0,-3.5);
% Rays
\draw[poporange, thick, ->] (-6, 1.5) -- (-3.5, 1.5);
\draw[poporange, thick, ->] (-6, 0.5) -- (-3.5, 0.5);
\draw[poporange, thick, ->] (-6, -0.5) -- (-3.5, -0.5);
% Convergence by lens
\draw[poporange, thick] (-3.5, 1.5) -- (-2.5, 0.8);
\draw[poporange, thick] (-3.5, 0.5) -- (-2.5, 0.2);
\draw[poporange, thick] (-3.5, -0.5) -- (-2.5, -0.4);
% Convergence by eye ON retina
\draw[poporange, thick] (-2.5, 0.8) -- (1.8, 0);
\draw[poporange, thick] (-2.5, 0.2) -- (1.8, 0);
\draw[poporange, thick] (-2.5, -0.4) -- (1.8, 0);
\fill[poporange] (1.8, 0) circle (2pt) node[above right] {Image nette};
\node[popdark, align=center] at (0, 3.2) {\textbf{Hypermétropie corrigée}\\(Verre convergent + Œil au repos)};
\node[poppurple] at (-1.5, -2) {Cristallin au repos};
\end{scope}
\end{tikzpicture}
\legende{À gauche : œil hypermétrope non corrigé. Le foyer théorique est derrière la rétine. L'œil doit accommoder (cristallin bombé, flèches violettes) pour ramener l'image sur la rétine, causant fatigue. À droite : correction par un \textbf{verre convergent} (convexe, puissance positive) qui pré-converge les rayons, permettant à l'œil de rester au repos pour voir net de loin.}
\end{popfigure}

\begin{attention}[Hypermétropie latente vs manifeste]
Chez l'enfant et le jeune adulte, l'hypermétropie est souvent \textbf{latente} : l'accommodation compense totalement le défaut, la vision de loin est 10/10. Mais l'effort permanent cause des maux de tête, larmoiements, refus de la lecture. Seule une examen sous \textbf{cycloplégie} (gouttes paralysant l'accommodation, ex : collyre à la cyclopentolate) révèle le défaut total. Au BEPC, on retient : \emph{hypermétropie = effort d'accommodation permanent = fatigue de près}.
\end{attention}

\courssub{La presbytie : le vieillissement du cristallin}

\begin{savaistu}[La presbytie, une fatalité universelle]
Contrairement à la myopie ou l'hypermétropie qui sont des « erreurs de construction » de l'œil (forme, longueur), la \cle{presbytie} est un \textbf{processus physiologique normal} lié à l'âge. Elle touche \textbf{100\% de la population} à partir de 40-45 ans. Même un myope deviendra presbyte : il devra ôter ses lunettes de loin pour lire de près (ou porter des verres progressifs). C'est pourquoi nos grands-parents écartent le journal à bout de bras : leur « point proche » a reculé au-delà de la longueur de leurs bras !
\end{savaistu}

\begin{definition}[Presbytie]
La \cle{presbytie} est la \textbf{perte progressive de l'élasticité du cristallin} avec l'âge, entraînant une \textbf{diminution de la puissance d'accommodation}. Le point proche (distance minimale de vision nette) s'éloigne. Elle apparaît vers 40-45 ans et se stabilise vers 60-65 ans (accommodation résiduelle nulle).
\end{definition}

\begin{propriete}[Mécanisme et évolution]
\begin{itemize}
    \item \textbf{Cause :} Durcissement du noyau du cristallin (dépôt de protéines insolubles) et perte d'élasticité de la capsule. Le muscle ciliaire se contracte, mais le cristallin ne se bombe plus assez.
    \item \textbf{Évolution typique de l'amplitude d'accommodation (AA) :}
    \begin{center}
    \begin{tabular}{|c|c|c|}\hline
    \textbf{Âge} & \textbf{Amplitude d'accommodation (D)} & \textbf{Point proche (cm)} \\ \hline
    10 ans & $\sim 14$ D & $\sim 7$ cm \\ \hline
    20 ans & $\sim 10$ D & $\sim 10$ cm \\ \hline
    40 ans & $\sim 4,5$ D & $\sim 22$ cm \\ \hline
    50 ans & $\sim 2,0$ D & $\sim 50$ cm \\ \hline
    60 ans & $\sim 1,0$ D & $\sim 100$ cm \\ \hline
    70 ans & $\sim 0$ D & $\infty$ (plus d'accommodation) \\ \hline
    \end{tabular}
    \end{center}
\end{itemize}
\end{propriete}

\begin{popfigure}
\begin{tikzpicture}[scale=0.9, >=Stealth]
\begin{axis}[
    width=0.9\linewidth, height=8cm,
    xlabel={Âge (ans)}, ylabel={Amplitude d'accommodation (Dioptries)},
    xmin=10, xmax=70, ymin=0, ymax=15,
    grid=major, grid style={popline},
    tick label style={font=\small},
    label style={font=\small},
    axis line style={popdark},
    tick style={popdark},
]
\addplot[popcurve, thick, domain=10:70, samples=100] {18.5 - 0.3*exp(0.05*x)}; % Approximation courbe classique
\addplot[only marks, mark=*, poporange, thick] coordinates {
    (10,14) (20,10) (30,7) (40,4.5) (50,2) (60,1) (70,0)
};
\node[popdark, align=left, anchor=south west] at (axis cs:12,13) {
    \textbf{Chute de l'accommodation avec l'âge}\\
    Seuil presbytie génant ($\approx$ 3-4 D) vers 40-45 ans.
};
\end{axis}
\end{tikzpicture}
\legende{Évolution de l'amplitude d'accommodation en fonction de l'âge. La presbytie devient génante (AA $< 4$ D) vers 40-45 ans, imposant une correction pour la lecture (distance habituelle \SI{33}{cm}-\SI{40}{cm}).}
\end{popfigure}

\begin{experience}[Correction de la presbytie : la loupe de lecture]
Un presbyte de 55 ans (accommodation résiduelle $\approx 1,5$ D) veut lire à \SI{33}{cm} (3 D nécessaires). Son œil ne peut fournir que 1,5 D. Il lui manque 1,5 D.
\par \textbf{Solution :} Un \textbf{verre convergent} (loupe) de \textbf{+1,50 D} (ou +2,00 D pour plus de confort) porté \textbf{uniquement pour la vision de près}. Ce verre ajoute la convergence manquante. De loin, le presbyte retire ses lunettes (ou utilise la partie haute de verres progressifs).
\end{experience}

\courssub{Distinction fondamentale : Hypermétropie vs Presbytie}

\begin{aretenir}[Ne pas confondre : Hypermétropie $\neq$ Presbytie]
\begin{center}
\begin{tabularx}{\linewidth}{|l|X|X|}\hline
\rowcolor{popblueL}
\textbf{Critère} & \textbf{Hypermétropie} & \textbf{Presbytie} \\ \hline
\textbf{Origine} & Morphologique : œil trop court / pas assez convergent (souvent congénital). & Physiologique : vieillissement du cristallin (perte élasticité), inévitable. \\ \hline
\textbf{Âge d'apparition} & Souvent dépistée dans l'enfance (dépistage scolaire). & Apparition vers \textbf{40-45 ans}, évolution jusqu'à 60-65 ans. \\ \hline
\textbf{Mécanisme optique} & Puissance de l'œil insuffisante pour sa longueur. \textit{Accommodation normale mais sursollicitée}. & Puissance d'accommodation \textbf{insuffisante} (cristallin rigide). Œil souvent emmétrope au loin. \\ \hline
\textbf{Symptôme principal} & Fatigue, maux de tête, vision floue de près \textit{malgré effort}. & \textbf{Impossibilité} de voir net de près (bras trop courts), soulagement en écartant l'objet. \\ \hline
\textbf{Correction} & Verres \textbf{convergents} portés \textbf{en permanence} (loin + près). & Verres \textbf{convergents} portés \textbf{uniquement pour le près} (lecture, téléphone). \\ \hline
\end{tabularx}
\end{center}
\end{aretenir}

\courssub{Complément : L'astigmatisme (notion de programme)}

\begin{definition}[Astigmatisme]
L'\cle{astigmatisme} est un défaut de réfraction dû à une \textbf{déformation de la cornée} (ou plus rarement du cristallin) qui n'est plus sphérique mais \textbf{toroïdale} (forme de ballon de rugby : courbures différentes selon les méridiens).
\end{definition}

\begin{propriete}[Conséquences et correction]
\begin{itemize}
    \item \textbf{Symptôme :} Vision déformée, floue \textbf{à toutes les distances} (loin et près). Confusion entre lettres proches (H/M, E/B, 8/0). Fatigue visuelle, maux de tête.
    \item \textbf{Correction :} \textbf{Verres cylindriques (toriques)} de puissance adaptée selon l'axe du défaut (ex : $-1,00$ D $\times$ 90°). Ils compensent la différence de courbure cornéenne.
\end{itemize}
\end{propriete}

\courssub{Méthode : Analyser un schéma de défaut visuel}

\begin{methode}[Diagnostiquer un défaut visuel sur un schéma de rayons]
\begin{enumerate}
    \item \textbf{Repérer le foyer :} Identifier où les rayons incidents (parallèles = objet à l'infini) convergent \textbf{après} avoir traversé le système optique (cornée + cristallin), \textbf{sans tenir compte d'un éventuel verre correcteur}.
    \item \textbf{Comparer à la rétine :} Localiser la rétine (couche pigmentée au fond de l'œil).
    \begin{itemize}
        \item Foyer \textbf{en avant} de la rétine $\rightarrow$ \textbf{Myopie}.
        \item Foyer \textbf{sur} la rétine $\rightarrow$ \textbf{Emmétropie} (ou hypermétropie corrigée/accommodée).
        \item Foyer \textbf{en arrière} de la rétine (ou rayons convergents qui se croiseraient derrière) $\rightarrow$ \textbf{Hypermétropie}.
    \end{itemize}
    \item \textbf{Déduire le verre correcteur :}
    \begin{itemize}
        \item Myopie : Verre \textbf{divergent} (concave, puissance $< 0$) pour \textbf{reculer} le foyer.
        \item Hypermétropie : Verre \textbf{convergent} (convexe, puissance $> 0$) pour \textbf{avancer} le foyer.
        \item Presbytie : Verre \textbf{convergent} pour le \textbf{près seulement} (ajoute la convergence manquante pour la lecture).
    \end{itemize}
    \item \textbf{Vérifier la cohérence :} Le myope voit net de près sans lunettes ; l'hypermétrope fatigue de près ; le presbyte écarte le texte.
\end{enumerate}
\end{methode}

\courssub{Tableau de synthèse comparatif}

\begin{center}
\begin{tabularx}{\linewidth}{|l|X|X|X|X|}\hline
\rowcolor{popblueL}
\textbf{Anomalie} & \textbf{Cause principale} & \textbf{Siège de l'image (œil au repos, objet $\infty$)} & \textbf{Symptôme fonctionnel} & \textbf{Correction (Verres)} \\ \hline
\textbf{Myopie} & Œil trop long \textbf{ou} Système trop convergent & \textbf{En avant} de la rétine & Flou de \textbf{loin}, Net de \textbf{près} & \textbf{Divergents} (Concaves, $P < 0$) \\ \hline
\textbf{Hypermétropie} & Œil trop court \textbf{ou} Système pas assez convergent & \textbf{En arrière} de la rétine (théorique) & Fatigue, maux de tête, Flou de \textbf{près} (effort constant) & \textbf{Convergents} (Convexes, $P > 0$) \textbf{permanents} \\ \hline
\textbf{Presbytie} & \textbf{Vieillissement} : Cristallin rigide, perte élasticité & Sur la rétine (loin), \textbf{Impossible de converger assez pour le près} & Bras trop courts, \textbf{Flou de près} (dès 40-45 ans) & \textbf{Convergents} (Convexes, $P > 0$) \textbf{pour le près seulement} \\ \hline
\textbf{Astigmatisme} & Cornée torique (courbures inégales) & \textbf{Deux foyers} (lignes verticales/horizontales séparées) & Vision déformée \textbf{toutes distances}, confusion lettres & \textbf{Cylindriques} (Toriques, axe précis) \\ \hline
\end{tabularx}
\end{center}

\courssub{Exercice résolu : Application intégrée}

\begin{exoresolu}[Analyse de cas cliniques au dispensaire scolaire]
\textbf{Énoncé :} Lors d'une visite médicale au lycée de Ngaoundéré, trois élèves présentent les symptômes suivants :
\begin{enumerate}
    \item \textbf{Kamga (14 ans) :} « Je ne vois pas le tableau au fond de la classe, mais je lis mon livre sans problème. »
    \item \textbf{Amina (13 ans) :} « J'ai mal à la tête tous les soirs après les devoirs, mes yeux piquent, mais je vois bien le tableau. »
    \item \textbf{M. Njoya (52 ans, surveillant) :} « Depuis deux ans, je dois tenir le journal à \SI{50}{cm} pour le lire. De loin, je vois très bien sans lunettes. »
\end{enumerate}
Pour chaque cas : a) Identifier l'anomalie probable. b) Préciser le siège de l'image d'un objet à l'infini sur l'œil non corrigé au repos. c) Indiquer le type de verre correcteur et son usage (permanent / lecture seule).

\tcblower
\textbf{Solution détaillée :}

\textbf{Cas 1 : Kamga (14 ans)}
\begin{itemize}
    \item \textbf{a) Anomalie :} \textbf{Myopie}. Vision floue de loin, nette de près = signe pathognomonique.
    \item \textbf{b) Siège de l'image :} \textbf{En avant de la rétine}. L'œil est trop long ou trop convergent.
    \item \textbf{c) Correction :} \textbf{Verres divergents} (puissance négative, ex : -1,50 D). \textbf{Port permanent} (pour voir le tableau, la TV, conduire plus tard). Il les ôtera peut-être pour lire de très près si sa myopie est forte.
\end{itemize}

\textbf{Cas 2 : Amina (13 ans)}
\begin{itemize}
    \item \textbf{a) Anomalie :} \textbf{Hypermétropie}. Vision de loin bonne (accommodation compense), mais \textbf{asthénopie} (fatigue, maux de tête, picotements) lors du travail de près (sursollicitation de l'accommodation).
    \item \textbf{b) Siège de l'image :} \textbf{En arrière de la rétine} (foyer théorique). Au repos, l'image serait floue ; l'accommodation permanente la ramène sur la rétine.
    \item \textbf{c) Correction :} \textbf{Verres convergents} (puissance positive, ex : +1,00 D). \textbf{Port permanent} (soulage l'accommodation pour la vision de loin ET de près). Attention : si elle ne les met que pour lire, elle fatiguera toujours en regardant le tableau.
\end{itemize}

\textbf{Cas 3 : M. Njoya (52 ans)}
\begin{itemize}
    \item \textbf{a) Anomalie :} \textbf{Presbytie}. Âge $> 40$ ans, vision de loin intacte (emmétropie supposée), \textbf{recul du point proche} (bras trop courts).
    \item \textbf{b) Siège de l'image :} \textbf{Sur la rétine} pour la vision de loin (œil emmétrope au repos). Pour la vision de près, l'accommodation résiduelle est insuffisante : l'image se formerait \textbf{en arrière de la rétine} si on forçait la mise au point à \SI{33}{cm}.
    \item \textbf{c) Correction :} \textbf{Verres convergents} (puissance positive, ex : +2,00 D à 52 ans). \textbf{Port pour la lecture seule} (vision de près : téléphone, journal, couture). Il ne les met \textbf{pas} pour regarder au loin (cinéma, conduite, tableau).
\end{itemize}

\textbf{Conclusion pédagogique :} Ce cas illustre parfaitement la distinction \textbf{Hypermétropie (enfant/jeune, port permanent, fatigue)} vs \textbf{Presbytie (adulte $>40$ ans, port lecture seule, bras courts)}. Tous deux se corrigent par des verres convergents, mais la \textbf{cause} et le \textbf{mode de port} diffèrent radicalement.
\end{exoresolu}

\courssub{Activité d'intégration : Exercice d'entraînement}

\begin{exercice}[Analyse de schémas et calcul de correction]
\begin{enumerate}
    \item On observe un schéma de rayons parallèles incidents sur un œil. Après réfraction, les rayons convergent en un point F situé \textbf{entre le cristallin et la rétine}.
    \begin{enumerate}
        \item Quel est le défaut visuel ? Justifier.
        \item Quel type de verre corrige ce défaut ? Dessiner schématiquement l'effet de ce verre sur les rayons incidents.
        \item Un myope porte une correction de \textbf{-3,00 D}. Calculer son point lointain (en cm). Interpréter ce résultat pour sa vie quotidienne (tableau à \SI{4}{m}, lecture à \SI{30}{cm}).
    \end{enumerate}
    \item Un élève hypermétrope de +2,50 D ne porte pas ses lunettes. Il lit un livre à \SI{25}{cm}.
    \begin{enumerate}
        \item Quelle est l'accommodation totale nécessaire pour cet élève (vision de loin + vision de près) ? Rappel : $A = \frac{1}{d}$ (d en m).
        \item Expliquer pourquoi il risque des maux de tête en fin de journée.
    \end{enumerate}
    \item \textbf{Vrai ou Faux ?} Justifier brièvement.
    \begin{enumerate}
        \item « Un myope voit flou de près. »
        \item « La presbytie se corrige par des verres divergents. »
        \item « L'hypermétropie est due à un cristallin qui a perdu son élasticité. »
        \item « Un astigmate confond le H et le M. »
    \end{enumerate}
\end{enumerate}
\end{exercice}

\begin{corrige}[Exercice d'entraînement]
\begin{enumerate}
    \item \begin{enumerate}
        \item \textbf{Myopie.} Le foyer F est en avant de la rétine.
        \item \textbf{Verre divergent (concave).} Il diverge les rayons incidents ; l'œil myope les converge ensuite sur la rétine. (Schéma : lentille mince évasée aux bords, rayons sortants divergents).
        \item $P = -3,00$ D. Point lointain $PL = \frac{1}{|P|} = \frac{1}{3} \approx 0,33$ m = \textbf{33 cm}. Interprétation : Kamga voit net spontanément jusqu'à \SI{33}{cm}. Le tableau à \SI{4}{m} est flou (nécessite lunettes). La lecture à \SI{30}{cm} est \textbf{nette sans lunettes} (\SI{30}{cm} $<$ \SI{33}{cm}).
    \end{enumerate}
    \item \begin{enumerate}
        \item Vision de loin : accommodation pour corriger l'hypermétropie = +2,50 D. Vision de près (\SI{25}{cm} = 0,25 m) : accommodation nécessaire = $\frac{1}{0,25} = 4,00$ D. \textbf{Total = 2,50 + 4,00 = 6,50 D.}
        \item Une accommodation de 6,50 D est énorme pour un adolescent (amplitude max $\sim$ 10-12 D, mais effort soutenu intenable). C'est une \textbf{sursollicitation permanente} du muscle ciliaire $\rightarrow$ fatigue musculaire (asthénopie) $\rightarrow$ céphalées de tension, yeux rouges, larmoiements.
    \end{enumerate}
    \item \begin{enumerate}
        \item \textbf{Faux.} Le myope voit \textbf{net de près} (jusqu'à son point lointain) et flou de loin.
        \item \textbf{Faux.} La presbytie (manque de convergence pour le près) se corrige par des \textbf{verres convergents} (convexes, puissance positive).
        \item \textbf{Faux.} C'est la définition de la \textbf{presbytie}. L'hypermétropie est due à un œil trop court ou pas assez convergent (morphologie).
        \item \textbf{Vrai.} L'astigmatisme déforme les images (étirement selon un axe), rendant difficiles la discrimination des lettres formées de traits verticaux/horizontaux similaires (H/M, E/F/B/P, 8/0).
    \end{enumerate}
\end{enumerate}
\end{corrige}

\begin{aretenir}[Essentiel à retenir pour le BEPC - Section 3]
\begin{itemize}
    \item \textbf{Myopie :} Image \textbf{avant} rétine $\rightarrow$ Verre \textbf{Divergent} ($-$) $\rightarrow$ Flou \textbf{loin}, Net \textbf{près}.
    \item \textbf{Hypermétropie :} Image \textbf{arrière} rétine $\rightarrow$ Verre \textbf{Convergent} ($+$) \textbf{permanent} $\rightarrow$ Fatigue \textbf{près}, Maux de tête.
    \item \textbf{Presbytie :} Cristallin \textbf{rigide} (âge $>40$ ans) $\rightarrow$ Verre \textbf{Convergent} ($+$) \textbf{lecture seule} $\rightarrow$ Bras \textbf{trop courts}.
    \item \textbf{Méthode schéma :} 1) Foyer ? 2) vs Rétine ? 3) Défaut ? 4) Verre ? (Divergent = recule foyer / Convergent = avance foyer).
    \item \textbf{Distinction clé :} Hypermétropie = défaut morphologique (jeune, port permanent) ; Presbytie = vieillissement physiologique (adulte, port près seul).
\end{itemize}
\end{aretenir}

\coursec{Maladies de la vision : cataracte et glaucome}

Après avoir étudié les anomalies de la vision (myopie, hypermétropie, presbytie), qui sont des défauts de mise au point de l'image et qui se corrigent par des verres de lunettes, nous allons maintenant aborder deux \textbf{maladies de l'œil} plus graves : la \cle{cataracte} et le \cle{glaucome}. Contrairement aux anomalies précédentes, il s'agit de maladies qui touchent les organes de l'œil (le cristallin ou le nerf optique) et qui peuvent conduire à la \textbf{cécité} (perte totale de la vue) si elles ne sont pas soignées à temps. Au Cameroun, ces deux maladies sont les premières causes de cécité : il est donc important de bien les connaître pour favoriser le dépistage précoce.

\courssub{La cataracte : opacification du cristallin}

\textbf{Observation et constat.}\par
Une personne âgée se plaint d'une baisse progressive de la vue : elle voit « comme à travers un voile » ou un brouillard, elle est gênée par les lumières vives (soleil, phares des voitures la nuit) et elle perçoit mal les couleurs. À l'examen, le médecin constate que le cristallin, normalement transparent, est devenu \textbf{blanchâtre et opaque}.

\begin{experience}[Observation du cristallin]
\textbf{Matériel :} appareil d'examen de l'œil (lampe à fente) chez l'ophtalmologiste.
\textbf{Protocole :} on éclaire l'œil avec un fin faisceau de lumière et on observe le cristallin à travers la pupille.
\textbf{Observations :}
\begin{itemize}
    \item \textbf{Œil sain :} le cristallin est transparent et incolore ; il laisse passer la lumière, qui atteint la rétine.
    \item \textbf{Œil atteint de cataracte :} le cristallin présente des zones blanchâtres (opacités) qui bloquent ou diffusent la lumière ; la rétine reçoit mal la lumière.
\end{itemize}
\textbf{Interprétation :} la transparence du cristallin est indispensable à la formation d'une image nette sur la rétine. Si le cristallin devient opaque, les rayons lumineux ne traversent plus correctement.
\textbf{Conclusion :} la cataracte est une opacification du cristallin.
\end{experience}

\begin{definition}[Cataracte]
La \cle{cataracte} est une \textbf{opacification progressive et indolore du cristallin} qui empêche la lumière d'atteindre correctement la rétine. Elle entraîne une baisse progressive de la vision, pouvant aller jusqu'à la cécité si elle n'est pas traitée.
\end{definition}

\textbf{Causes et mécanisme.}\par
Le cristallin est une lentille naturelle transparente située derrière l'iris. Avec l'âge (cataracte de la personne âgée, la plus fréquente), mais aussi à cause de certains facteurs (diabète, coup ou blessure de l'œil, exposition excessive au soleil), le cristallin perd progressivement sa transparence : il devient trouble, comme une vitre sale. La lumière est alors diffusée dans tous les sens au lieu d'être convergée vers la rétine : l'image devient floue.

\begin{popfigure}
\begin{tikzpicture}[scale=0.9, every node/.style={font=\small}]
    % --- Œil sain ---
    \begin{scope}[xshift=-4.5cm]
        \draw[popdark, thick] (0,0) ellipse (2cm and 1.5cm);
        \draw[popgreen, thick] (0,0) ellipse (1.6cm and 1.1cm);
        \draw[popblue!50!white, thick, fill=popblue!10] (-2,0) arc (180:360:2cm and 0.5cm) -- cycle;
        % Cristallin sain (transparent)
        \draw[popgreen!70!black, thick, fill=popgreen!10] (-0.3,-0.6) rectangle (0.7,0.6);
        \node[popgreen!70!black] at (0.2,0.95) {\textbf{Cristallin clair}};
        % Nerf optique
        \draw[poporange, thick] (2,0) -- (3.2,0);
        \node[poporange] at (2.6,-0.7) {\textbf{Nerf optique}};
        % Rayons lumineux qui traversent
        \draw[->, popgold, thick] (-3.8,0.5) -- (-2,0.3);
        \draw[->, popgold, thick] (-3.8,-0.5) -- (-2,-0.3);
        \draw[->, popgold, thick] (0.7,0.2) -- (1.55,0.1);
        \draw[->, popgold, thick] (0.7,-0.2) -- (1.55,-0.1);
        \node[popgold] at (-3,0.85) {Lumière};
        \node[align=center] at (0,-2.3) {\textbf{Œil sain}\\cristallin transparent};
    \end{scope}

    % --- Œil cataracté ---
    \begin{scope}[xshift=4.5cm]
        \draw[popdark, thick] (0,0) ellipse (2cm and 1.5cm);
        \draw[popgreen, thick] (0,0) ellipse (1.6cm and 1.1cm);
        \draw[popblue!50!white, thick, fill=popblue!10] (-2,0) arc (180:360:2cm and 0.5cm) -- cycle;
        % Cristallin opaque
        \draw[poporange, thick, fill=poporange!40, pattern=north east lines, pattern color=poporange!60] (-0.3,-0.6) rectangle (0.7,0.6);
        \node[poporange] at (0.2,0.95) {\textbf{Cristallin opaque}};
        \draw[poporange, thick] (2,0) -- (3.2,0);
        \node[poporange] at (2.6,-0.7) {\textbf{Nerf optique}};
        % Lumière diffusée
        \draw[->, popgold, thick] (-3.8,0.5) -- (-2,0.3);
        \draw[->, popgold, thick] (-3.8,-0.5) -- (-2,-0.3);
        \draw[popgold, densely dashed, thick] (0.7,0.3) -- (1.55,0.7);
        \draw[popgold, densely dashed, thick] (0.7,-0.3) -- (1.55,-0.7);
        \draw[popgold, densely dashed, thick] (0.7,0) -- (1.55,0.2);
        \node[popgold] at (-3,0.85) {Lumière};
        \node[popmuted, align=center] at (0.2,0) {\tiny diffusion};
        \node[align=center] at (0,-2.3) {\textbf{Cataracte}\\la lumière est diffusée};
    \end{scope}
\end{tikzpicture}
\legende{Comparaison schématique : à gauche, œil sain (le cristallin transparent laisse converger la lumière vers la rétine) ; à droite, œil atteint de cataracte (le cristallin opaque diffuse la lumière : l'image sur la rétine est floue).}
\end{popfigure}

\textbf{Symptômes et évolution.}\par
La cataracte évolue lentement, sur plusieurs années. Les signes sont : une vision voilée, une gêne face à la lumière vive (éblouissement), une baisse progressive de la vision de loin et de près. Au stade avancé, la pupille apparaît blanchâtre et la personne peut devenir aveugle.

\begin{propriete}[Traitement de la cataracte]
Le \textbf{seul traitement efficace est la chirurgie} : aucun collyre ni médicament ne peut rendre sa transparence à un cristallin opacifié.
\begin{itemize}
    \item Le chirurgien \textbf{retire le cristallin opaque}.
    \item Il le remplace par une \textbf{lentille artificielle} (implant) placée dans l'œil.
    \item Le résultat est excellent : le patient \textbf{retrouve la vue} dans la grande majorité des cas. La cataracte est donc une cécité \textbf{curable}.
\end{itemize}
\end{propriete}

\begin{savaistu}[Campagnes de chirurgie de la cataracte au Cameroun]
Chaque année, le Ministère de la Santé Publique (MINSANTÉ), avec l'aide d'organisations non gouvernementales, organise des \textbf{campagnes gratuites de chirurgie de la cataracte} dans les hôpitaux du pays (par exemple à Douala, Yaoundé ou Garoua). Des centaines de personnes âgées, souvent pauvres, retrouvent ainsi la vue. C'est une action importante de santé publique dans notre pays.
\end{savaistu}

\courssub{Le glaucome : destruction du nerf optique par une pression trop élevée}

\textbf{Observation.}\par
Un homme de 55 ans consulte parce qu'il heurte souvent les objets situés sur les côtés, alors qu'il voit encore bien droit devant lui. Il ne ressent \textbf{aucune douleur}. Le médecin mesure la \textbf{pression à l'intérieur de l'œil} (pression intra-oculaire) : elle est trop élevée. L'examen du fond de l'œil montre que le \textbf{nerf optique est abîmé}.

\begin{experience}[Mesure de la pression intra-oculaire]
\textbf{Matériel :} appareil de mesure de la pression de l'œil (tonomètre) chez l'ophtalmologiste.
\textbf{Protocole :} on mesure la pression régnant à l'intérieur du globe oculaire, puis on examine le nerf optique au fond de l'œil.
\textbf{Observations :}
\begin{itemize}
    \item \textbf{Œil sain :} la pression intra-oculaire est normale ; le nerf optique a un aspect normal.
    \item \textbf{Œil glaucomateux :} la pression intra-oculaire est \textbf{élevée} ; le nerf optique est creusé et abîmé ; le champ visuel (l'espace que l'on voit sans bouger les yeux) est réduit sur les côtés.
\end{itemize}
\textbf{Interprétation :} l'intérieur de l'œil contient un liquide, l'\textbf{humeur aqueuse}, qui est produit en permanence et doit s'évacuer en permanence. Si l'évacuation se fait mal, le liquide s'accumule et la pression augmente. Cette pression excessive comprime et détruit peu à peu les fibres du nerf optique.
\textbf{Conclusion :} le glaucome est une maladie due à une pression intra-oculaire trop élevée, qui détruit progressivement le nerf optique.
\end{experience}

\begin{definition}[Glaucome]
Le \cle{glaucome} est une maladie de l'œil caractérisée par une \textbf{augmentation de la pression intra-oculaire} (due à une mauvaise évacuation de l'humeur aqueuse), qui entraîne une \textbf{destruction progressive et irréversible du nerf optique}.
\end{definition}

\textbf{Mécanisme.}\par
L'humeur aqueuse est un liquide transparent qui remplit la partie avant de l'œil. Elle est produite en permanence et s'évacue normalement par de petits canaux situés à la jonction entre l'iris et la cornée.
\begin{itemize}
    \item \textbf{Œil sain :} production = évacuation, donc la pression reste normale et stable.
    \item \textbf{Œil glaucomateux :} les canaux d'évacuation fonctionnent mal ; le liquide s'accumule, la pression monte et écrase le nerf optique à l'arrière de l'œil. Les fibres nerveuses détruites \textbf{ne se régénèrent pas}.
\end{itemize}

\begin{popfigure}
\begin{tikzpicture}[scale=1.0, every node/.style={font=\small}]
    % Champ visuel normal
    \begin{scope}[xshift=-4.5cm]
        \draw[popdark, thick, fill=popgreen!20] (0,0) circle (2.2cm);
        \fill[popgreen!60] (0,0) circle (0.25cm);
        \fill[popdark] (1.1,-0.3) circle (0.2cm);
        \node[popdark] at (1.1,-0.75) {\tiny tache aveugle};
        \node[align=center, popdark] at (0,-2.9) {\textbf{Champ visuel normal}\\on voit large, sur les côtés};
    \end{scope}

    % Champ visuel glaucome : vision en tunnel
    \begin{scope}[xshift=4.5cm]
        \draw[popdark, thick, fill=popmuted!40] (0,0) circle (2.2cm);
        \fill[popgreen!50] (0,0) circle (0.7cm);
        \draw[poppink, thick, <->] (-0.7,0.9) -- (0.7,0.9) node[midway, above] {\tiny vision restante};
        \node[align=center, poppink] at (0,-2.9) {\textbf{Glaucome avancé}\\« vision en tunnel »};
    \end{scope}

    \draw[poporange, thick, ->] (-1.8,-1.6) -- (1.8,-1.6) node[midway, above] {\textbf{Évolution sans traitement}};
\end{tikzpicture}
\legende{Évolution du champ visuel dans le glaucome : à gauche, champ visuel normal (vision large) ; à droite, glaucome avancé : la vision des côtés est perdue, il ne reste qu'une vision centrale étroite, dite « en tunnel ».}
\end{popfigure}

\textbf{Évolution : un mal « silencieux ».}\par
\begin{itemize}
    \item \textbf{Au début :} aucun symptôme, aucune douleur ; le malade ne se rend compte de rien.
    \item \textbf{Ensuite :} la vision des côtés (vision périphérique) diminue peu à peu.
    \item \textbf{Stade avancé :} « vision en tunnel » : on ne voit plus que droit devant soi.
    \item \textbf{Stade final :} cécité totale et \textbf{irréversible}.
\end{itemize}

\begin{savaistu}[Le « voleur silencieux de la vue »]
Le glaucome est surnommé le \textbf{« voleur silencieux de la vue »} car il détruit le nerf optique pendant des années \textbf{sans douleur et sans signe visible} au début. Au Cameroun, beaucoup de personnes atteintes l'ignorent. La seule protection est le \textbf{dépistage régulier chez l'ophtalmologiste à partir de 40 ans} (mesure de la pression de l'œil et examen du fond de l'œil), surtout si des membres de la famille sont atteints.
\end{savaistu}

\begin{propriete}[Traitement du glaucome]
Le but du traitement est de \textbf{faire baisser la pression intra-oculaire} pour arrêter la destruction du nerf optique. \textbf{Attention : les dégâts déjà causés sont irréversibles} ; le traitement ne fait que stabiliser la maladie.
\begin{enumerate}
    \item \textbf{Traitement médical :} des \textbf{collyres} (gouttes pour les yeux) à instiller chaque jour, \textbf{à vie}, qui diminuent la production d'humeur aqueuse ou facilitent son évacuation.
    \item \textbf{Traitement chirurgical ou laser :} création d'une voie d'évacuation pour le liquide, quand les collyres ne suffisent pas.
\end{enumerate}
La difficulté principale est l'\textbf{observance} : le malade doit prendre son traitement tous les jours, sans interruption, même s'il ne ressent rien.
\end{propriete}

\courssub{Comparaison cataracte et glaucome : points clés pour le BEPC}

Il est essentiel de ne pas confondre ces deux grandes causes de cécité. Le tableau suivant résume les différences à maîtriser.

\begin{center}
\begin{tabularx}{\linewidth}{|l|X|X|}
\hline
\rowcolor{popblueL}
\textbf{Critère} & \textbf{Cataracte} & \textbf{Glaucome} \\
\hline
\textbf{Organe atteint} & Le cristallin (il devient opaque) & Le nerf optique (détruit par la pression) \\
\hline
\textbf{Cause} & Perte de transparence du cristallin (vieillissement, diabète, blessure) & Pression intra-oculaire trop élevée (mauvaise évacuation de l'humeur aqueuse) \\
\hline
\textbf{Signes} & Vision voilée, éblouissement, baisse progressive de la vue & Aucun signe au début, puis perte de la vision des côtés (« vision en tunnel ») \\
\hline
\textbf{Douleur} & Non & Non (forme chronique) \\
\hline
\textbf{Traitement} & Chirurgie : ablation du cristallin + implant & Collyres à vie, parfois laser ou chirurgie \\
\hline
\textbf{Résultat} & \textbf{Guérison} : la vue revient & \textbf{Stabilisation seulement} : la vue perdue ne revient pas (irréversible) \\
\hline
\textbf{Prévention} & Opérer dès que la gêne est importante & Dépistage régulier après 40 ans \\
\hline
\end{tabularx}
\end{center}

\begin{attention}[Erreurs fréquentes au BEPC : ne pas confondre cataracte et glaucome]
\begin{itemize}
    \item \textbf{Ne pas écrire :} « Le glaucome se guérit par une opération qui rend la vue. » $\rightarrow$ \textbf{Faux.} Le traitement du glaucome fait baisser la pression mais \textbf{ne répare pas} le nerf optique détruit : la vue perdue ne revient jamais.
    \item \textbf{Ne pas écrire :} « La cataracte est due à une pression trop forte dans l'œil. » $\rightarrow$ \textbf{Faux.} La cataracte est une opacification du cristallin ; la pression de l'œil est normale dans la cataracte.
    \item \textbf{Mots-clés à retenir :} \cle{cataracte} = \textbf{cristallin opaque / chirurgie / guérison / voile}. \cle{glaucome} = \textbf{pression élevée / nerf optique / irréversible / dépistage / vision en tunnel}.
\end{itemize}
\end{attention}

\courssub{Activités d'intégration et exercice résolu}

\begin{exoresolu}[Analyse de cas : cataracte ou glaucome ?]
\textbf{Énoncé :} Deux patients consultent à l'hôpital de district. Pour chaque cas, identifie la maladie probable (cataracte ou glaucome) et justifie ta réponse avec \textbf{deux arguments}.

\textbf{Cas 1 :} M. K., 72 ans, cultivateur. Il se plaint d'une baisse progressive de la vue depuis 2 ans et il est gêné par le soleil. À l'examen : le cristallin est blanchâtre et opaque ; la pression intra-oculaire est normale.

\textbf{Cas 2 :} Mme M., 55 ans, institutrice. Elle consulte pour un simple contrôle, sans plainte particulière. À l'examen : la pression intra-oculaire est élevée ; le nerf optique est creusé et abîmé ; le champ visuel est réduit sur les côtés ; le cristallin est transparent.

\tcblower
\textbf{Solution détaillée :}

\textbf{Cas 1 : il s'agit d'une cataracte.}
\begin{itemize}
    \item \textbf{Argument 1 (les signes) :} personne âgée, baisse progressive et indolore de la vue, gêne face à la lumière vive : c'est le tableau typique de la cataracte.
    \item \textbf{Argument 2 (l'examen) :} le cristallin est opaque (organe atteint dans la cataracte) et la pression intra-oculaire est normale, ce qui élimine le glaucome.
\end{itemize}
\textit{Conduite à tenir :} opération chirurgicale (ablation du cristallin opaque et pose d'un implant) ; le pronostic est excellent : le patient retrouvera la vue.

\textbf{Cas 2 : il s'agit d'un glaucome.}
\begin{itemize}
    \item \textbf{Argument 1 (la pression et le nerf) :} la pression intra-oculaire est élevée et le nerf optique est abîmé : ce sont les deux signes principaux du glaucome.
    \item \textbf{Argument 2 (le champ visuel) :} la vision des côtés est réduite alors que la vision centrale est encore bonne (la patiente ne se plaint de rien) : c'est l'évolution typique, « silencieuse », du glaucome. De plus, le cristallin est transparent, ce qui élimine la cataracte.
\end{itemize}
\textit{Conduite à tenir :} collyres tous les jours à vie pour faire baisser la pression, avec des contrôles réguliers. Les dégâts déjà présents sont malheureusement irréversibles : d'où l'importance du dépistage précoce.
\end{exoresolu}

\begin{exercice}[Entraînement BEPC : QCM et Vrai/Faux]
\textbf{1. (QCM) La cataracte est une opacification :}
\begin{enumerate}
    \item de la cornée ;
    \item du cristallin ;
    \item du nerf optique ;
    \item de la rétine.
\end{enumerate}

\textbf{2. (QCM) Le traitement efficace de la cataracte est :}
\begin{enumerate}
    \item des collyres à vie ;
    \item le port de lunettes de soleil ;
    \item la chirurgie avec pose d'un implant ;
    \item le repos des yeux.
\end{enumerate}

\textbf{3. (Vrai ou Faux ? Justifie ta réponse.)}
\begin{enumerate}
    \item Le glaucome provoque dès le début une forte douleur de l'œil.
    \item Dans le glaucome, la perte de la vision commence par les côtés.
    \item La cataracte est une cécité curable.
    \item Le dépistage du glaucome repose sur la mesure de la pression intra-oculaire et l'examen du fond de l'œil.
\end{enumerate}

\textbf{4. (Problème) Un homme de 45 ans n'a jamais consulté d'ophtalmologiste. Son père est devenu aveugle à cause d'un glaucome. Que lui conseilles-tu ? Justifie ta réponse.}
\end{exercice}

\begin{corrige}[Exercice n]
\textbf{1. Réponse 2 :} la cataracte est une opacification du cristallin.

\textbf{2. Réponse 3 :} le seul traitement efficace est la chirurgie (ablation du cristallin opaque et pose d'un implant). Les collyres servent à traiter le glaucome, pas la cataracte.

\textbf{3.}
\begin{enumerate}
    \item \textbf{Faux.} Le glaucome évolue au début \textbf{sans douleur ni signe apparent} : c'est pour cela qu'on l'appelle le « voleur silencieux de la vue ».
    \item \textbf{Vrai.} La vision périphérique (des côtés) disparaît d'abord, ce qui donne au stade avancé une « vision en tunnel » ; la vision centrale n'est atteinte qu'à la fin.
    \item \textbf{Vrai.} L'opération chirurgicale (ablation du cristallin opaque + implant) permet de retrouver la vue dans la grande majorité des cas.
    \item \textbf{Vrai.} La mesure de la pression intra-oculaire et l'examen du nerf optique au fond de l'œil permettent de dépister le glaucome avant l'apparition des symptômes.
\end{enumerate}

\textbf{4.} Je lui conseille de \textbf{consulter régulièrement un ophtalmologiste pour un dépistage du glaucome} (mesure de la pression intra-oculaire et examen du fond de l'œil). Justification : il a plus de 40 ans et son père était glaucomateux, ce qui augmente son risque ; or le glaucome ne provoque aucun signe au début, et seul un dépistage précoce permet de commencer le traitement avant la destruction du nerf optique, destruction qui serait irréversible.
\end{corrige}

\courssub{Synthèse : ce qu'il faut retenir pour l'examen}

\begin{aretenir}[Maladies de la vision : cataracte et glaucome]
\begin{enumerate}
    \item \textbf{Cataracte} = opacification du \textbf{cristallin}. Cause principale : le vieillissement. Signes : vision voilée, éblouissement, baisse progressive de la vue. \textbf{Traitement : chirurgie (ablation du cristallin + implant) $\rightarrow$ guérison.} Cécité \textbf{curable}.
    \item \textbf{Glaucome} = destruction du \textbf{nerf optique} par une \textbf{pression intra-oculaire trop élevée} (mauvaise évacuation de l'humeur aqueuse). Aucun signe au début (« voleur silencieux »), puis perte de la vision des côtés (« vision en tunnel »). \textbf{Traitement : collyres à vie, parfois chirurgie $\rightarrow$ stabilisation seulement.} Cécité \textbf{irréversible}.
    \item \textbf{Distinction essentielle :} la cataracte touche un milieu transparent de l'œil (le cristallin) et se guérit ; le glaucome touche le nerf qui transmet les messages au cerveau (le nerf optique) et ne se guérit pas.
    \item \textbf{Hygiène et prévention :} dépistage du glaucome chez l'ophtalmologiste à partir de 40 ans (surtout en cas d'antécédent familial) ; opérer la cataracte dès qu'elle gêne la vie quotidienne ; protéger les yeux du soleil et des blessures.
    \item \textbf{Contexte camerounais :} la cataracte et le glaucome sont les premières causes de cécité dans notre pays ; les campagnes gratuites de chirurgie de la cataracte permettent à de nombreux patients de retrouver la vue.
\end{enumerate}
\end{aretenir}

\coursec{Hygiène de la vision}

Après avoir étudié les maladies de l'œil comme la cataracte et le glaucome, qui surviennent souvent avec l'âge ou à la suite de pathologies, nous abordons maintenant un volet essentiel : la \cle{prévention}. Beaucoup de troubles visuels peuvent être évités, ou leur apparition retardée, par des gestes simples au quotidien. L'hygiène de la vision, c'est l'ensemble des règles de conduite qui visent à préserver l'intégrité de l'œil et la qualité de la vue tout au long de la vie.

\begin{definition}[Hygiène de la vision]
L'\cle{hygiène de la vision} est l'ensemble des règles de conduite et des comportements quotidiens visant à préserver l'œil contre les agressions physiques, chimiques, infectieuses et contre la fatigue visuelle, afin de maintenir une bonne acuité visuelle le plus longtemps possible.
\end{definition}

\courssub{Éclairage adapté et distances de travail}

\textbf{Observation.} Un élève qui lit dans un coin sombre de la classe ou sous un arbre par temps couvert plisse les yeux, rapproche le livre de son visage et se plaint de maux de tête en fin de journée. Un autre, installé près de la fenêtre avec un livre tenu à 35~cm, lit confortablement pendant une heure.

\textbf{Interprétation.} Un éclairage insuffisant oblige le cristallin à \cle{accommoder} davantage (bomber au maximum pour voir net de près) et la pupille à se dilater (\cle{mydriase}) pour laisser entrer plus de lumière ; cette dilatation réduit la profondeur de champ et l'effort d'accommodation prolongé fatigue les muscles ciliaires. À l'inverse, un éclairage trop violent ou mal orienté crée un \cle{éblouissement} (reflet sur la page) ou un \cle{contre-jour} qui diminue le contraste de l'image perçue. La distance de lecture idéale (30 à 40~cm) correspond à la zone où l'effort d'accommodation est modéré pour un œil normal ou correctement corrigé.

\begin{aretenir}[Règles d'éclairage et de distance]
\begin{itemize}
\item \textbf{Éclairage suffisant et homogène :} travailler sous une lumière d'environ 300 à 500~lux (lampe de bureau orientable + éclairage général de la pièce). Éviter les zones d'ombre portées par la main qui écrit.
\item \textbf{Pas d'éblouissement :} la source lumineuse ne doit pas être dans le champ visuel direct ni se refléter sur le papier ou l'écran. Placer la lampe du côté opposé à la main qui écrit (à gauche pour un droitier).
\item \textbf{Pas de contre-jour :} ne pas se placer dos à une fenêtre claire sans store ; l'écran ou la page ne doit pas être plus sombre que l'arrière-plan.
\item \textbf{Distance de lecture :} 30 à 40~cm entre les yeux et le livre ou l'écran. Ne pas lire couché (distance variable, éclairage mauvais) ni dans un véhicule en mouvement (secousses $\Rightarrow$ accommodation instable).
\end{itemize}
\end{aretenir}

\begin{popfigure}
\begin{tikzpicture}[scale=0.95, every node/.style={font=\small}]
% --- MAUVAISE POSTURE (gauche) ---
\begin{scope}[xshift=-4.2cm]
% Table
\draw[popdark, thick] (-2.2,0) -- (2.2,0);
\draw[popdark, thick] (-2.0,0) -- (-2.0,-1.2);
\draw[popdark, thick] (2.0,0) -- (2.0,-1.2);
% Livre
\draw[popmuted, thick] (0.3,0) rectangle (1.5,0.15);
% Personnage penché en avant
\draw[poporange, thick] (-1.2,0) -- (-0.9,0.9); % dos courbé
\draw[poporange, fill=poporange!25, thick] (-0.55,1.15) circle (0.35); % tête penchée
\draw[poporange, thick] (-0.9,0.7) -- (0.2,0.25); % bras vers le livre
% Distance trop courte
\draw[<->, poppink, thick] (-0.35,1.0) -- (0.55,0.35) node[midway, above left, poppink] {15 cm : trop près};
% Lampe mal placée (dans le champ visuel)
\draw[popgold, thick] (1.8,0) -- (1.8,1.0);
\draw[popgold, fill=popgold!40, thick] (1.8,1.2) circle (0.25);
\draw[popgold, ->, thick] (1.6,1.1) -- (0.2,0.9) node[midway, above, popgold] {éblouissement};
\node[align=center, poppink, font=\small\bfseries] at (0,-1.7) {Mauvaise posture\\dos courbé, trop près, éblouissement};
\end{scope}
% --- BONNE POSTURE (droite) ---
\begin{scope}[xshift=4.2cm]
% Table
\draw[popdark, thick] (-2.2,0) -- (2.2,0);
\draw[popdark, thick] (-2.0,0) -- (-2.0,-1.2);
\draw[popdark, thick] (2.0,0) -- (2.0,-1.2);
% Livre
\draw[popblue, thick] (0.4,0) rectangle (1.7,0.15);
% Personnage droit
\draw[popgreen, thick] (-1.2,0) -- (-1.2,1.0); % dos droit
\draw[popgreen, fill=popgreen!25, thick] (-1.2,1.35) circle (0.35); % tête droite
\draw[popgreen, thick] (-1.2,0.7) -- (0.3,0.25); % bras vers le livre
% Distance correcte
\draw[<->, popgreen, thick] (-0.85,1.2) -- (0.75,0.3) node[midway, above, popgreen] {35 cm : correct};
% Lampe sur le côté opposé
\draw[popblue, thick] (-2.0,0) -- (-2.0,1.0);
\draw[popblue, fill=popblue!30, thick] (-2.0,1.2) circle (0.25);
\draw[popblue, ->, thick] (-1.8,1.1) -- (0.6,0.4) node[midway, above, popblue] {éclairage latéral};
\node[align=center, popgreen, font=\small\bfseries] at (0,-1.7) {Bonne posture\\dos droit, 30--40 cm, lampe latérale};
\end{scope}
\end{tikzpicture}
\legende{Schéma comparatif des postures de lecture. À gauche : erreurs fréquentes (distance trop courte de 15~cm, dos courbé, lampe dans le champ visuel provoquant un éblouissement). À droite : posture correcte (dos droit, distance de 30 à 40~cm, lampe placée latéralement du côté opposé à la main qui écrit).}
\end{popfigure}

\courssub{Gestion des écrans et pauses visuelles}

Avec la généralisation des téléphones portables, des tablettes et des ordinateurs, la \cle{fatigue visuelle} liée aux écrans devient fréquente chez les jeunes. Les symptômes : yeux secs, picotements, vision floue passagère, maux de tête frontaux, douleurs de la nuque.

\begin{experience}[Effet de la fixation prolongée d'un écran]
\textbf{Question :} pourquoi la lecture sur écran fatigue-t-elle plus que la lecture sur papier ?\\
\textbf{Matériel :} un téléphone portable, un livre, un chronomètre, un questionnaire de symptômes (échelle de 0 à 5 : sécheresse, flou, douleur).\\
\textbf{Protocole :}
\begin{enumerate}
\item Lire un texte sur papier pendant 20~min (distance 35~cm, bon éclairage). Noter les symptômes.
\item Après 10~min de pause, lire le même texte sur téléphone (même distance, luminosité moyenne). Noter les symptômes.
\item Recommencer en appliquant la règle des pauses (voir méthode ci-dessous) et comparer les résultats.
\end{enumerate}
\textbf{Observations attendues :} les scores de sécheresse, de flou et de douleur sont plus élevés après la lecture sur écran sans pause. L'application de la règle des pauses réduit nettement les symptômes.\\
\textbf{Interprétation :} devant un écran, le \cle{clignement des paupières} diminue fortement (d'environ 15 à 20 clignements par minute à 5 à 7 par minute) ; le film lacrymal qui protège et humidifie la cornée s'évapore alors davantage, d'où la sécheresse. De plus, l'écran impose un effort d'accommodation permanent.\\
\textbf{Conclusion :} l'écran n'est pas dangereux en soi, mais son usage prolongé sans pause dessèche la surface de l'œil et fatigue l'accommodation.
\end{experience}

\begin{methode}[Règle des pauses visuelles (20-20-20) et réglages de l'écran]
\begin{enumerate}
\item \textbf{Toutes les 20 minutes :} lever les yeux de l'écran.
\item \textbf{Regarder au loin :} fixer un objet situé à au moins 6~mètres (relaxation de l'accommodation) pendant \textbf{20 secondes}.
\item \textbf{Cligner délibérément des yeux :} 5 à 10 clignements complets pour réétaler le film lacrymal sur la cornée.
\item \textbf{Régler l'écran :} luminosité adaptée à l'éclairage de la pièce (pas plus claire que le mur derrière), contraste élevé, taille des caractères confortable.
\item \textbf{Respecter les distances :} 50 à 70~cm pour un écran d'ordinateur (distance d'un bras tendu), 30 à 40~cm pour un téléphone ou une tablette.
\item \textbf{Limiter la durée totale :} ne pas dépasser environ 2~h par jour d'écran de loisir en dehors du temps scolaire.
\end{enumerate}
\end{methode}

\courssub{Protection contre les agressions extérieures}

L'œil est exposé à des agressions physiques (poussières, projectiles), chimiques (produits ménagers, chaux, pesticides) et lumineuses (rayons ultraviolets du soleil, arc de soudure). Au Cameroun, l'harmattan (vent sec chargé de poussières fines) et les pistes non goudronnées augmentent le risque d'irritation et de corps étrangers dans l'œil. Les travaux de soudure, de menuiserie et l'agriculture (manipulation de produits phytosanitaires) sont des activités à haut risque.

\begin{exemplebox}[Le soudeur du quartier : le « coup d'arc »]
Un jeune soudeur de Douala travaille sans masque de protection, avec de simples lunettes de soleil. Après une matinée de soudure à l'arc, il ressent une douleur intense aux deux yeux, une gêne insupportable à la lumière et des larmes abondantes. Diagnostic : brûlure de la cornée par les rayons ultraviolets intenses émis par l'arc électrique (\cle{kératite} d'origine lumineuse). Traitement : collyre cicatrisant prescrit par le médecin, repos des yeux pendant 48~h. La guérison est complète en quelques jours, mais le risque de récidive et de cataracte précoce est réel si l'œil n'est pas protégé.\\
\textbf{Leçon :} seul un masque de soudeur à verre filtrant spécial protège contre les rayonnements de l'arc de soudure. Les lunettes de soleil ordinaires ne suffisent pas.
\end{exemplebox}

\begin{aretenir}[Protections selon l'activité]
\begin{center}
\begin{tabularx}{\linewidth}{|l|X|}\hline
\rowcolor{popblueL}\textbf{Situation à risque} & \textbf{Protection adaptée} \\ \hline
Soudure à l'arc & Masque de soudeur à verre filtrant spécial \\ \hline
Menuiserie, meulage, bricolage & Lunettes de protection enveloppantes résistant aux chocs \\ \hline
Agriculture (produits phytosanitaires) & Lunettes étanches, voire visière de protection \\ \hline
Soleil intense, harmattan, plage & Lunettes de soleil de bonne qualité, filtrant 100\,\% des ultraviolets (marquage CE) \\ \hline
Natation & Lunettes de piscine étanches (protection contre le chlore) \\ \hline
\end{tabularx}
\end{center}
\end{aretenir}

\courssub{Hygiène oculaire et comportements à risque}

\begin{attention}[Erreurs fréquentes et idées reçues]
\begin{itemize}
\item \textbf{« Se frotter les yeux soulage. »} FAUX. Se frotter les yeux avec les mains sales introduit des microbes (risque de \cle{conjonctivite} ou d'infection de la cornée) et peut irriter la cornée. \textbf{Conduite correcte :} se laver les mains, puis rincer l'œil au sérum physiologique ou à l'eau propre.
\item \textbf{« Porter des lunettes fatigue les yeux ou abîme la vue. »} FAUX. Ne \textbf{pas} porter sa correction oblige l'œil à forcer en permanence (fatigue d'accommodation, maux de tête). La correction \textbf{repose} l'œil. La vue évolue naturellement avec l'âge (croissance de l'œil, presbytie), et non à cause des lunettes.
\item \textbf{« Les collyres traditionnels (infusions de feuilles, eau de riz...) soignent tout. »} DANGEREUX. Ces préparations ne sont pas stériles : elles peuvent provoquer des infections graves de l'œil, et leur composition inadaptée peut brûler la cornée. \textbf{Règle :} ne jamais mettre de collyre sans avis médical ; un collyre entamé se jette après quelques semaines (risque de contamination).
\item \textbf{« Je vois bien, pas besoin de consulter. »} PIÈGE. Le glaucome évolue longtemps sans aucun symptôme. Un contrôle régulier, surtout après 40~ans ou en cas d'antécédents familiaux, est indispensable.
\end{itemize}
\end{attention}

\courssub{Alimentation et vision : le rôle clé de la vitamine A}

La rétine contient des cellules sensibles à la lumière : les \cle{bâtonnets} (vision en lumière faible, vision nocturne) et les \cle{cônes} (vision des couleurs en pleine lumière). Les bâtonnets contiennent un pigment visuel qui se décompose sous l'effet de la lumière et doit être régénéré en permanence : cette régénération exige de la \cle{vitamine A}. En cas de carence en vitamine A, la régénération du pigment est bloquée : c'est la \cle{cécité nocturne} (l'enfant ne voit plus dans la pénombre, trébuche le soir). Si la carence persiste, la conjonctive et la cornée se dessèchent, puis la cornée se ramollit et peut se perforer : la cécité devient alors irréversible.

\begin{savaistu}[Vitamine A et huile de palme rouge au Cameroun]
La carence en vitamine A est une cause majeure de cécité infantile évitable dans le monde. Au Cameroun, l'\textbf{huile de palme rouge} (non raffinée, extraite du fruit du palmier à huile) est une source exceptionnelle de provitamine A (bêta-carotène), bien plus riche que la carotte. Une cuillère à soupe par jour couvre largement les besoins quotidiens d'un enfant. Les \textbf{légumes-feuilles} (feuilles de manioc, folong, amarante, épinards) et les \textbf{mangues} bien mûres en sont aussi de bonnes sources. La vitamine A étant soluble dans les graisses, sa absorption est meilleure lorsque ces aliments sont préparés avec un peu de matière grasse. \textbf{Action :} consommer chaque jour ces aliments locaux, peu coûteux et disponibles sur tous les marchés.
\end{savaistu}

\begin{exoresolu}[Conseil nutritionnel justifié]
\textbf{Énoncé :} une mère amène son fils de 6~ans qui « ne voit plus le soir depuis 3~mois ». L'examen révèle une conjonctive sèche et de petites taches grisâtres sur le blanc de l'œil ; l'acuité visuelle est normale en plein jour. Le régime alimentaire de l'enfant est à base de bâton de manioc et de riz, sans huile de palme rouge ni légumes-feuilles.\\
\textbf{Question :} établir le diagnostic probable et proposer deux conseils nutritionnels justifiés.\\
\tcblower
\textbf{Solution :}
\begin{enumerate}
\item \textbf{Diagnostic :} carence en vitamine A, responsable d'une cécité nocturne. Les bâtonnets de la rétine ne peuvent plus régénérer leur pigment visuel, d'où une vision impossible en lumière faible, alors que la vision diurne (assurée par les cônes) reste normale.
\item \textbf{Conseil 1 :} introduire \textbf{l'huile de palme rouge} dans la sauce quotidienne (une cuillère à soupe par jour pour l'enfant). \textbf{Justification :} elle est très riche en bêta-carotène, transformé en vitamine A dans l'organisme ; la régénération du pigment des bâtonnets redevient possible et la vision nocturne se rétablit en quelques semaines si la cornée n'est pas encore atteinte.
\item \textbf{Conseil 2 :} ajouter chaque jour \textbf{des légumes-feuilles vert foncé} (feuilles de manioc, folong) cuits brièvement avec un peu d'huile. \textbf{Justification :} ils apportent un supplément de provitamine A ; la présence de matière grasse améliore l'absorption de la vitamine A, qui est soluble dans les graisses.
\end{enumerate}
\end{exoresolu}

\courssub{Suivi médical et conduite à tenir en cas de traumatisme}

\begin{methode}[Quand consulter un ophtalmologiste ?]
\begin{enumerate}
\item \textbf{Contrôle systématique :} à l'entrée au collège, puis tous les 2~ans si l'on ne porte pas de correction ; chaque année si l'on porte des lunettes ou des lentilles.
\item \textbf{Consultation rapide (dans les jours qui suivent) :} baisse de la vue brutale ou progressive, maux de tête frontaux répétés en fin de journée, vision double, cercles colorés autour des lumières (signe possible de glaucome aigu).
\item \textbf{Antécédents familiaux :} un parent proche atteint de glaucome impose un dépistage précoce et régulier (mesure de la tension de l'œil, examen du fond de l'œil).
\item \textbf{Maladies générales :} diabète ou hypertension artérielle $\Rightarrow$ examen régulier du fond de l'œil (la rétine peut être atteinte).
\end{enumerate}
\end{methode}

\begin{aretenir}[Conduite à tenir face à un traumatisme oculaire]
\begin{center}
\begin{tabularx}{\linewidth}{|l|X|}\hline
\rowcolor{popblueL}\textbf{Type de traumatisme} & \textbf{Conduite à tenir (premiers secours)} \\ \hline
Corps étranger (poussière, sable, insecte) & \textbf{Ne pas frotter.} Cligner fortement des yeux pour faire couler les larmes. Rincer au sérum physiologique ou à l'eau propre. Si la gêne persiste : consulter. \\ \hline
Choc sur l'œil (balle, coup, chute) & \textbf{Ne pas appuyer sur l'œil.} Protéger avec une coque rigide (gobelet en plastique découpé) sans comprimer. Consulter en urgence (risque d'hémorragie interne). \\ \hline
Brûlure chimique (chaux, acide, soude, pesticide) & \textbf{URGENCE ABSOLUE.} Rincer \textbf{immédiatement et abondamment} à l'eau courante pendant \textbf{15 à 20~min}, paupières bien écartées. Aller ensuite aux urgences. \textbf{Ne mettre ni collyre, ni huile, ni lait.} \\ \hline
Brûlure par la lumière (arc de soudure) & Fermer les yeux, appliquer une compresse froide propre, consulter rapidement. \\ \hline
Plaie de l'œil (objet pointu, éclat) & \textbf{Ne pas retirer l'objet.} Protéger avec une coque rigide sans pression. Allonger le blessé, tête surélevée. Urgence chirurgicale. \\ \hline
\end{tabularx}
\end{center}
\end{aretenir}

\courssub{Synthèse : les règles d'or de l'hygiène visuelle}

\begin{popfigure}
\begin{tikzpicture}[scale=1, every node/.style={font=\small}]
% Ligne 1 : 4 règles
\foreach \x/\titre/\desc/\couleur in {
  0/{1. Éclairage}/{300--500 lux, sans éblouissement}/popgold,
  4.6/{2. Distance}/{30--40 cm (livre), 50--70 cm (écran)}/popblue,
  9.2/{3. Pauses}/{règle 20-20-20, cligner des yeux}/popgreen,
  13.8/{4. Écrans}/{luminosité adaptée, durée limitée}/poppurple}{
  \draw[\couleur, fill=\couleur!15, thick, rounded corners=0.15cm] (\x-2.0,0) rectangle (\x+2.0,1.6);
  \node[\couleur, font=\small\bfseries, align=center] at (\x,1.2) {\titre};
  \node[popdark, font=\footnotesize, align=center] at (\x,0.5) {\desc};
}
% Ligne 2 : 4 règles
\foreach \x/\titre/\desc/\couleur in {
  0/{5. Protection}/{lunettes adaptées (soleil, soudure, bricolage)}/poporange,
  4.6/{6. Hygiène}/{mains propres, pas de collyre sans avis médical}/poppink,
  9.2/{7. Alimentation}/{vitamine A : huile de palme rouge, légumes-feuilles}/popgreen,
  13.8/{8. Suivi}/{contrôle régulier, urgence si traumatisme}/popdark}{
  \draw[\couleur, fill=\couleur!15, thick, rounded corners=0.15cm] (\x-2.0,-2.2) rectangle (\x+2.0,-0.6);
  \node[\couleur, font=\small\bfseries, align=center] at (\x,-1.0) {\titre};
  \node[popdark, font=\footnotesize, align=center] at (\x,-1.7) {\desc};
}
\end{tikzpicture}
\legende{Les 8 règles d'or de l'hygiène visuelle : (1) éclairage suffisant sans éblouissement ; (2) distances de lecture et d'écran respectées ; (3) pauses visuelles régulières ; (4) usage raisonné des écrans ; (5) protections adaptées aux activités à risque ; (6) hygiène des mains et des collyres ; (7) alimentation riche en vitamine A ; (8) suivi médical régulier.}
\end{popfigure}

\begin{exercice}[Mise en situation : rédiger une fiche-conseil]
\textbf{Consigne :} tu es élève, chargé(e) d'une causerie éducative dans un club de santé scolaire. Tu dois concevoir une affiche « 5 gestes pour protéger mes yeux » destinée aux élèves de 3ème. Pour chaque geste : (1) énoncer la règle clairement, (2) expliquer en une phrase le mécanisme biologique qu'elle protège, (3) donner un exemple concret de situation à risque au Cameroun. Utilise la méthode « règle + justification » vue en cours.
\end{exercice}

\begin{corrige}[Exercice : fiche-conseil « 5 gestes pour protéger mes yeux »]
\textbf{Geste 1 — Éclairage :} « Je lis et j'écris sous une lumière suffisante, avec la lampe placée du côté opposé à ma main. » --- \emph{Justification :} un éclairage suffisant évite l'effort maximal d'accommodation et la dilatation excessive de la pupille, sources de fatigue visuelle. --- \emph{Exemple :} faire ses devoirs au salon éclairé seulement par la télévision (contre-jour et lumière insuffisante).

\textbf{Geste 2 — Distance :} « Je garde mon livre ou mon téléphone à 35--40~cm de mes yeux. » --- \emph{Justification :} à cette distance, l'effort d'accommodation du cristallin reste modéré pour un œil normal ou corrigé. --- \emph{Exemple :} lire couché sur le ventre, le livre posé sur l'oreiller à 15~cm des yeux.

\textbf{Geste 3 — Pauses écrans :} « Toutes les 20 minutes, je regarde au loin (au moins 6~m) pendant 20 secondes et je cligne plusieurs fois des yeux. » --- \emph{Justification :} regarder au loin relâche l'accommodation (le cristallin se repose) et cligner réétale le film lacrymal qui protège la cornée du dessèchement. --- \emph{Exemple :} passer 2~h d'affilée sur les réseaux sociaux sans lever les yeux, puis avoir les yeux rouges et brûlants le soir.

\textbf{Geste 4 — Protection :} « Je porte des lunettes de soleil filtrant les ultraviolets en saison sèche et un masque de soudeur à l'atelier. » --- \emph{Justification :} les rayons ultraviolets agressent la cornée et le cristallin (risque de cataracte précoce) ; le rayonnement intense de l'arc de soudure brûle la cornée. --- \emph{Exemple :} rouler en moto-taxi à Yaoundé pendant l'harmattan sans lunettes $\Rightarrow$ irritation et conjonctivite ; souder sans masque $\Rightarrow$ « coup d'arc ».

\textbf{Geste 5 — Alimentation et suivi :} « Je mange chaque jour de l'huile de palme rouge et des légumes-feuilles, et je fais contrôler ma vue régulièrement. » --- \emph{Justification :} la vitamine A est indispensable à la régénération du pigment visuel des bâtonnets (vision nocturne) ; le dépistage précoce du glaucome, maladie longtemps sans symptôme, permet de sauver la vue. --- \emph{Exemple :} un enfant nourri uniquement de foufou et de riz sans sauce rouge développe une cécité nocturne ; un adulte atteint de glaucome non dépisté perd définitivement une partie de son champ visuel.
\end{corrige}

\coursec{Synthèse et activités d'intégration}

Nous venons de voir que l'hygiène de la vision repose sur des gestes simples : bien s'éclairer, ménager ses yeux, consulter régulièrement un ophtalmologiste. Mais pourquoi ces conseils sont-ils si importants ? Pour le comprendre, il faut rassembler tout ce que nous avons appris dans cette leçon : l'anatomie de l'œil, la formation de l'image, les anomalies et les maladies de la vision. Cette dernière partie te donne une vision d'ensemble, puis t'apprend à résoudre les problèmes de synthèse comme on te les posera au BEPC.

\courssub{Bilan structuré : la chaîne de la vision}

\textbf{Voir, c'est recevoir la lumière, la transformer en message nerveux et l'interpréter.}\par
Ce processus fonctionne comme une \cle{chaîne} : chaque maillon doit être intact pour que la vision soit nette. Si un seul maillon est défaillant, toute la vision est perturbée.

\begin{definition}[Chaîne de la vision]
La \cle{chaîne de la vision} est la succession des étapes qui conduisent de la lumière à la perception de l'image :
\begin{enumerate}
\item la \cle{lumière} réfléchie par l'objet pénètre dans l'œil ;
\item elle traverse les \cle{milieux transparents} (cornée, humeur aqueuse, cristallin, corps vitré) qui la réfractent et la font converger ;
\item l'image se forme sur la \cle{rétine}, où les cellules sensorielles (cônes et bâtonnets) transforment la lumière en \cle{messages nerveux} ;
\item le \cle{nerf optique} transporte ces messages jusqu'au \cle{cerveau}, qui les analyse : c'est la sensation visuelle.
\end{enumerate}
\end{definition}

\begin{popfigure}
\begin{center}
\begin{tikzpicture}[font=\small, >=Stealth,
etape/.style={draw=popblue, fill=popblueL, rounded corners, minimum height=0.9cm, align=center, text width=2.1cm},
trouble/.style={draw=poporange, fill=poporangeL, rounded corners, align=center, text width=2.6cm, font=\footnotesize}]
\node[etape, align=center] (l) at (0,0){Lumière\\ (objet)};
\node[etape, align=center] (m) at (2.9,0){Milieux\\ transparents};
\node[etape, align=center] (r) at (5.8,0){Rétine\\ (image)};
\node[etape, align=center] (n) at (8.7,0){Nerf\\ optique};
\node[etape, align=center] (c) at (11.6,0){Cerveau\\ (perception)};
\draw[->, thick, popdark] (l) -- (m);
\draw[->, thick, popdark] (m) -- (r);
\draw[->, thick, popdark] (r) -- (n);
\draw[->, thick, popdark] (n) -- (c);
\node[trouble] (t1) at (2.9,-2.4) {Cataracte (cristallin opaque) ; myopie, hypermétropie, presbytie (focalisation)};
\node[trouble] (t2) at (5.8,-2.4) {Atteinte des cellules sensorielles};
\node[trouble] (t3) at (8.7,-2.4) {Glaucome (destruction des fibres du nerf optique)};
\draw[->, dashed, poporange] (t1) -- (m);
\draw[->, dashed, poporange] (t2) -- (r);
\draw[->, dashed, poporange] (t3) -- (n);
\end{tikzpicture}
\end{center}
\legende{Schéma-bilan de la chaîne de la vision. Les flèches pleines indiquent le trajet normal de l'information visuelle ; les flèches pointillées orange montrent où se situent les principaux troubles étudiés.}
\end{popfigure}

\begin{propriete}[Localisation du trouble et traitement (synthèse admise)]
Toute perturbation d'un \cle{milieu transparent}, de la \cle{rétine} ou du \cle{nerf optique} altère la vision. La \cle{localisation} du trouble détermine le traitement :
\begin{itemize}
\item trouble de la \emph{focalisation} (globe oculaire trop long ou trop court, cristallin peu déformable) $\Rightarrow$ correction par des \textbf{lentilles} (lunettes) ;
\item trouble de la \emph{transparence} (cataracte) $\Rightarrow$ \textbf{chirurgie} (remplacement du cristallin) ;
\item destruction des fibres du \emph{nerf optique} (glaucome) $\Rightarrow$ traitement précoce pour \textbf{stopper l'évolution}, car les fibres détruites ne se régénèrent pas.
\end{itemize}
\end{propriete}

\courssub{Tableau récapitulatif général}

Le tableau ci-dessous rassemble l'essentiel de la leçon. Apprends-le : il résume les cinq troubles et maladies au programme.

\begin{center}
\small
\begin{tabularx}{\linewidth}{|l|X|X|X|X|}
\hline
\rowcolor{popblueL} \textbf{Trouble} & \textbf{Cause / siège} & \textbf{Symptômes} & \textbf{Correction / traitement} & \textbf{Prévention} \\
\hline
\textbf{Myopie} & Globe oculaire trop long : l'image se forme \emph{en avant} de la rétine & Vision floue de loin, nette de près & Lentilles \textbf{divergentes} & Dépistage scolaire, lunettes adaptées \\
\hline
\textbf{Hypermétropie} & Globe oculaire trop court : l'image se formerait \emph{en arrière} de la rétine & Vision floue de près, fatigue oculaire & Lentilles \textbf{convergentes} & Dépistage, lunettes adaptées \\
\hline
\textbf{Presbytie} & Cristallin qui perd son élasticité avec l'âge : accommodation insuffisante & Difficulté à lire de près après 40--45 ans & Lentilles \textbf{convergentes} pour la lecture & Correction adaptée à l'âge \\
\hline
\textbf{Cataracte} & Cristallin qui devient \emph{opaque} (âge, traumatisme, diabète) & Vision voilée, éblouissement, baisse progressive de la vue & \textbf{Chirurgie} : remplacement du cristallin & Protection solaire, contrôle du diabète, consultation \\
\hline
\textbf{Glaucome} & Excès de pression dans l'œil qui détruit les fibres du \emph{nerf optique} & Souvent sans symptôme au début ; rétrécissement du champ visuel, puis cécité si non traité & Collyres, parfois chirurgie : \textbf{stoppe} la maladie mais ne répare pas & \textbf{Dépistage régulier} (pression oculaire) dès 40 ans \\
\hline
\end{tabularx}
\end{center}

\begin{attention}[Anomalie ou maladie ?]
Ne confonds pas les \cle{anomalies optiques} (myopie, hypermétropie, presbytie), qui sont des défauts de focalisation corrigés par des \textbf{lunettes}, et les \cle{maladies} (cataracte, glaucome), qui sont des atteintes des structures de l'œil nécessitant un \textbf{traitement médical ou chirurgical}. Le glaucome, en particulier, ne se corrige pas avec des lunettes !
\end{attention}

\courssub{Méthode : résoudre un problème de synthèse type BEPC}

Au BEPC, on ne te demandera pas seulement de réciter : on te présentera une \emph{situation concrète} (un élève qui voit flou, un examen médical, un schéma d'œil) et tu devras raisonner. Voici la démarche à suivre.

\begin{methode}[Résoudre un problème de synthèse en 5 étapes]
\begin{enumerate}
\item \textbf{Analyser la situation et les documents} : lis l'énoncé en entier, souligne les données utiles (symptômes, résultats d'examen, schéma), repère la question posée.
\item \textbf{Mobiliser les savoirs} : demande-toi quelle partie de la leçon est concernée (anatomie ? accommodation ? anomalies ? maladies ? hygiène ?).
\item \textbf{Formuler une hypothèse ou un diagnostic} : propose une réponse précise (ex. : « ce patient est myope »).
\item \textbf{Justifier par des arguments scientifiques} : appuie-toi sur les documents et sur tes connaissances (ex. : « l'image se forme en avant de la rétine, caractéristique de la myopie »).
\item \textbf{Conclure} : réponds clairement à la question et, si demandé, propose une solution (correction, traitement, mesure de prévention).
\end{enumerate}
\end{methode}

\begin{attention}[Erreur fréquente au BEPC]
Répondre par une simple liste de mots (« myopie, lentille divergente ») fait \textbf{perdre des points}. Au BEPC, toute réponse doit être \cle{rédigée en phrases complètes} et \cle{justifiée} : énonce le diagnostic, donne l'argument scientifique qui l'appuie, puis conclus. Une réponse juste mais non justifiée n'est qu'à moitié réussie.
\end{attention}

\courssub{Interprétation de documents : s'entraîner comme à l'examen}

\textbf{Trois types de documents reviennent souvent dans les sujets :}\par

\begin{itemize}
\item \textbf{Schéma d'œil défectueux} : repère où se forme l'image par rapport à la rétine. Image \emph{en avant} de la rétine $\Rightarrow$ myopie ; image \emph{en arrière} $\Rightarrow$ hypermétropie. La correction se déduit : lentille divergente pour la myopie, convergente pour l'hypermétropie.
\item \textbf{Cliché de fond d'œil} : c'est la photographie de la rétine. Un fond d'œil normal montre la papille (point de départ du nerf optique) et les vaisseaux. Une papille creusée et élargie évoque un glaucome ; un cristallin opaque empêche d'observer le fond d'œil (cataracte).
\item \textbf{Résultat d'acuité visuelle} : l'acuité se mesure en dixièmes (10/10 = vision normale). Une acuité de 3/10 de loin redevenue 10/10 avec des verres divergents confirme une myopie.
\end{itemize}

\begin{exoresolu}[L'élève qui plisse les yeux]
\textbf{Énoncé.} Étienne, élève de 3ème dans un collège de Yaoundé, s'assoit de plus en plus près du tableau. De loin, il voit flou ; de près, il lit sans difficulté. Le médecin scolaire mesure son acuité visuelle de loin : 4/10 aux deux yeux. Un schéma de son œil montre que l'image d'un objet éloigné se forme \emph{en avant} de la rétine.
\begin{enumerate}
\item Identifier le trouble de la vision dont souffre Étienne. Justifier.
\item Proposer une correction et expliquer son rôle.
\item Citer une conséquence possible de l'absence de correction sur sa scolarité.
\end{enumerate}
\tcblower
\textbf{Solution rédigée.}
\begin{enumerate}
\item Étienne est atteint de \textbf{myopie}. En effet, il voit flou de loin mais net de près, et le schéma montre que l'image d'un objet éloigné se forme \emph{en avant} de la rétine : c'est la signature d'un globe oculaire trop long, caractéristique de la myopie.
\item On corrige la myopie par des lunettes à verres \textbf{divergents}. Ces lentilles écartent légèrement les rayons lumineux avant leur entrée dans l'œil, ce qui déplace le point de formation de l'image \emph{en arrière}, exactement sur la rétine : la vision de loin redevient nette.
\item Sans correction, Étienne ne lit plus ce qui est écrit au tableau : il peut prendre du retard, se fatiguer les yeux et voir ses résultats scolaires baisser. Il faut donc qu'il porte ses lunettes en permanence en classe.
\end{enumerate}
\end{exoresolu}

\courssub{Application à des situations concrètes camerounaises}

\begin{exemplebox}[Trois situations de la vie courante]
\begin{itemize}
\item \textbf{Conseiller un élève myope non corrigé} : lui expliquer que ses maux de tête et sa vision floue viennent d'un défaut de focalisation facilement corrigé ; l'encourager à consulter un ophtalmologiste (par exemple dans un hôpital de district) et à porter ses lunettes sans honte.
\item \textbf{Sensibiliser un parent au dépistage du glaucome} : le glaucome est sournois, il ne provoque aucun symptôme au début mais détruit peu à peu le nerf optique. Après 40 ans, un contrôle régulier de la pression oculaire permet de le détecter tôt et de stopper son évolution par des collyres. Un parent qui attend d'avoir mal aux yeux consulte souvent trop tard.
\item \textbf{Choisir un éclairage d'étude} : travailler dans une pièce bien éclairée, avec une lampe placée de manière à ne pas créer d'éblouissement ni d'ombre sur le cahier, à une distance de lecture d'environ 30 cm. Éviter d'étudier à la lueur faible d'une lampe-tempête ou d'un téléphone, qui fatigue les yeux.
\end{itemize}
\end{exemplebox}

\begin{savaistu}[80 \% des cécités sont évitables ou curables]
Selon l'Organisation mondiale de la santé, environ \textbf{80 \% des cécités} dans le monde pourraient être \cle{évitées ou guéries} : la cataracte se soigne par une opération simple, le glaucome se contrôle s'il est dépisté tôt, et la plupart des anomalies optiques se corrigent avec de simples lunettes. Au Cameroun, des campagnes de dépistage et de chirurgie de la cataracte sont régulièrement organisées dans les hôpitaux. D'où l'importance capitale de l'\cle{hygiène de la vision} et du \cle{dépistage} : consulter tôt, c'est souvent sauver sa vue.
\end{savaistu}

\courssub{Activités d'intégration : rédiger une démarche complète}

\begin{exercice}[Le cas de Mme Ngo Bell]
Mme Ngo Bell, 62 ans, commerçante à Douala, se plaint depuis deux ans de voir « comme à travers un brouillard », même avec ses lunettes. Elle est de plus en plus gênée par la lumière du soleil. L'ophtalmologiste observe son œil : la rétine et le nerf optique sont normaux, mais le cristallin est trouble et opaque.
\begin{enumerate}
\item Rédiger un diagnostic argumenté (maladie + justification à partir des données).
\item Proposer une solution et expliquer son principe.
\item Citer une mesure de prévention de cette maladie.
\end{enumerate}
\end{exercice}

\begin{corrige}[Le cas de Mme Ngo Bell]
\begin{enumerate}
\item Mme Ngo Bell est atteinte de \textbf{cataracte}. En effet, sa vision est voilée (« brouillard ») et elle est éblouie par la lumière, symptômes typiques de cette maladie. De plus, l'examen montre que le cristallin est devenu \emph{opaque} alors que la rétine et le nerf optique sont sains : le trouble se situe donc dans un milieu transparent, et non dans la réception ou la transmission du message nerveux. Or l'opacification du cristallin est précisément la définition de la cataracte.
\item La solution est la \textbf{chirurgie} : le chirurgien retire le cristallin opaque et le remplace par un cristallin artificiel transparent. La lumière peut alors à nouveau traverser le cristallin et former une image nette sur la rétine : la vision est restaurée. Les lunettes ne suffisent pas, car elles ne peuvent pas rendre transparent un cristallin devenu opaque.
\item Pour prévenir ou retarder la cataracte, on peut se \textbf{protéger les yeux du soleil} (chapeau, lunettes de soleil), contrôler les maladies comme le diabète qui la favorisent, et consulter régulièrement un ophtalmologiste à partir d'un certain âge.
\end{enumerate}
\end{corrige}

\begin{exercice}[À toi de jouer : le dépistage en famille]
Ton oncle, 45 ans, n'a jamais consulté d'ophtalmologiste car « il voit très bien ». Rédige un court texte (5 à 8 lignes) pour le convaincre de faire contrôler sa pression oculaire, en utilisant au moins trois savoirs de la leçon.
\end{exercice}

\begin{corrige}[À toi de jouer : le dépistage en famille]
Exemple de réponse attendue : le glaucome est une maladie due à un excès de pression dans l'œil, qui détruit progressivement les fibres du nerf optique. Son danger est qu'elle ne provoque \emph{aucun symptôme} au début : on peut « très bien voir » tout en perdant lentement son champ visuel. Les fibres nerveuses détruites ne se régénèrent jamais, mais un traitement précoce (collyres) stoppe l'évolution de la maladie. Le dépistage par simple mesure de la pression oculaire est donc recommandé dès 40 ans. Consulter maintenant, c'est protéger sa vue pour toute la vie.
\end{corrige}

\begin{aretenir}[L'essentiel de la leçon]
\begin{itemize}
\item La vision suit une chaîne : lumière $\rightarrow$ milieux transparents $\rightarrow$ rétine $\rightarrow$ nerf optique $\rightarrow$ cerveau.
\item Myopie, hypermétropie et presbytie sont des \textbf{anomalies de focalisation} corrigées par des lentilles ; cataracte et glaucome sont des \textbf{maladies} traitées par chirurgie ou médicaments.
\item La localisation du trouble détermine le traitement ; le glaucome exige un dépistage précoce car il est d'abord silencieux.
\item Au BEPC : analyser, mobiliser ses savoirs, diagnostiquer, justifier, conclure — toujours en phrases rédigées.
\end{itemize}
\end{aretenir}

\newpage

\coursec{Activité d'intégration}

\coursec{Leçon 10 : L'œil et la vision}

\courssub{Objectifs et compétences visées (Programme MINESEC 3ème)}
\begin{aretenir}[Compétences du programme]
Cette leçon s'inscrit dans la famille de situations « \textbf{Comprendre le fonctionnement du corps humain et en prendre soin} ». Elle vise les compétences suivantes :
\begin{itemize}
    \item \textbf{Savoirs :} Anatomie de l'œil ; Formation de l'image et accommodation ; Anomalies de la vision (myopie, hypermétropie, presbytie) ; Maladies de la vision (cataracte, glaucome) ; Hygiène de la vision.
    \item \textbf{Savoir-faire :} Observer et identifier les structures de l'œil sur un schéma ou un modèle ; Expliquer le mécanisme de la vision et de l'accommodation ; Caractériser les anomalies et maladies de la vision (causes, conséquences, corrections/traitements) ; Argumenter des conseils d'hygiène visuelle ; Mobiliser les connaissances dans une situation d'intégration (diagnostic simulé).
    \item \textbf{Savoir-être :} Adopter une attitude responsable face à sa santé visuelle (écrans, éclairage, protection UV, dépistage) ; Faire preuve de rigueur dans le raisonnement scientifique (observation $\rightarrow$ hypothèse $\rightarrow$ conclusion).
\end{itemize}
\end{aretenir}

\courssub{Anatomie de l'œil}
\begin{definition}[L'œil, organe de la vision]
L'œil est un organe sensoriel sphérique (globe oculaire) logé dans l'orbite, permettant la réception de la lumière et sa conversion en influx nerveux transmis au cerveau par le nerf optique.
\end{definition}

\begin{popfigure}
\begin{tikzpicture}[scale=1.2, every node/.style={font=\small}]
    % Sclère (blanc de l'œil)
    \draw[popdark, thick, fill=popcream!20] (0,0) circle (2.5cm);
    % Cornée (avant, plus courbe)
    \draw[popblue, thick, fill=popblue!10] (-2.5,0) arc (180:360:2.5cm and 0.8cm) -- cycle;
    % Iris
    \draw[popgreen, thick, fill=popgreen!30] (-1.2,0) arc (180:0:1.2cm and 0.6cm) -- cycle;
    % Pupille (trou au centre de l'iris)
    \fill[black] (0,0) circle (0.4cm);
    % Cristallin
    \draw[poporange, thick, fill=poporange!20] (0.3,0) ellipse (0.8cm and 0.5cm);
    % Corps ciliaire / Zonule
    \draw[popmuted, thin] (-1.2,0.6) -- (0.3,0.5);
    \draw[popmuted, thin] (-1.2,-0.6) -- (0.3,-0.5);
    % Rétine (arrière)
    \draw[poppink, very thick] (1.5,1.8) arc (90:-90:0.5cm and 1.8cm);
    % Nerf optique
    \draw[poppurple, thick, fill=poppurple!20] (2.5,0) -- (3.5,0) -- (3.5,-0.3) -- (2.5,-0.3) -- cycle;
    % Humor vitré (zone centrale)
    \fill[popblue!5] (0,0) ellipse (1.5cm and 2.0cm);
    % Légendes (flèches)
    \draw[->, popdark] (-3.5, 1.5) -- (-2.5, 0.8) node[left, align=right] {\textbf{Cornée}\\(transparente, convergente)};
    \draw[->, popdark] (-3.5, 0) -- (-1.2, 0) node[left, align=center]{\textbf{Iris}\\(diaphragme, contrôle pupille)};
    \draw[->, popdark] (-3.5, -1.5) -- (-1.2, -0.6) node[left, align=right] {\textbf{Pupille}\\(orifice variable)};
    \draw[->, popdark] (1.5, 1.2) -- (0.5, 0.5) node[right, align=center]{\textbf{Cristallin}\\(lentille déformable)};
    \draw[->, popdark] (1.5, -1.2) -- (0.5, -0.5) node[right, align=center]{\textbf{Zonule / Muscle ciliaire}\\(accommodation)};
    \draw[->, popdark] (3.8, 1.5) -- (2.5, 1.5) node[right, align=center]{\textbf{Rétine}\\(photorécepteurs : cônes/bâtonnets)};
    \draw[->, popdark] (3.8, -1.5) -- (2.5, -0.3) node[right, align=center]{\textbf{Nerf optique}\\(vers le cerveau)};
    \draw[->, popdark] (-0.5, 2.2) -- (0, 1.8) node[above, align=center]{\textbf{Humor aqueux}\\(antérieur)};
    \draw[->, popdark] (0.5, -2.2) -- (0, -1.8) node[below, align=center]{\textbf{Humor vitré}\\(postérieur, gélatineux)};
\end{tikzpicture}
\legende{Coupe schématique de l'œil humain (vue sagittale médiane). Les structures transparentes (cornée, humor aqueux, cristallin, humor vitré) forment le système optique. La rétine est le plan récepteur.}
\end{popfigure}

\begin{propriete}[Rôle des principales structures]
\begin{center}
\begin{tabularx}{\linewidth}{|l|X|}\hline
\rowcolor{popblueL} \textbf{Structure} & \textbf{Rôle optique / physiologique} \\\hline
\textbf{Cornée} & Première lentille convergente (puissance $\approx 43$ D), fixe, transparente, avasculaire.\\\hline
\textbf{Iris / Pupille} & Diaphragme réglant la quantité de lumière entrante (réflexe photomoteur).\\\hline
\textbf{Cristallin} & Lentille convergente \textbf{variables} (puissance $\approx 19$ à $33$ D), transparente, élastique. Siège de l'accommodation.\\\hline
\textbf{Humor aqueux / vitré} & Milieux transparents maintenant la forme du globe et la pression intra-oculaire.\\\hline
\textbf{Rétine} & Couche neurosensorielle : photorécepteurs (cônes : vision couleur, détail ; bâtonnets : vision crépusculaire) $\rightarrow$ cellules bipolaires $\rightarrow$ cellules ganglionnaires (axones = nerf optique).\\\hline
\textbf{Choroïde} & Couche vasculaire (nutrition) et pigmentée (absorbe la lumière parasite).\\\hline
\textbf{Sclère} & Enveloppe fibreuse blanche, protectrice, maintien de la forme.\\\hline
\end{tabularx}
\end{center}
\end{propriete}

\courssub{Formation de l'image et accommodation}
\begin{experience}[Modélisation optique de l'œil (Schéma de principe)]
L'œil se comporte comme un système optique convergent unique (équivalent : cornée + cristallin) projetant une image réelle, inversée et réduite sur la rétine.
\begin{popfigure}
\begin{tikzpicture}[scale=0.8, every node/.style={font=\small}]
    % Objet (loin)
    \draw[popdark, thick] (-6, 1.5) -- (-6, -1.5) node[midway, left] {Objet AB (loin)};
    % Rayons incidents parallèles
    \draw[->, popblue, thick] (-6, 1.5) -- (-1, 1.5);
    \draw[->, popblue, thick] (-6, -1.5) -- (-1, -1.5);
    \draw[->, popblue, thick] (-6, 0) -- (-1, 0);
    % Système optique (Cornée + Cristallin)
    \draw[poporange, very thick] (-1, -1) -- (-1, 1) node[midway, above, rotate=90] {Système optique (C + Cr)};
    % Foyer image F' sur rétine
    \draw[popblue, thick] (-1, 1.5) -- (1.5, 0);
    \draw[popblue, thick] (-1, -1.5) -- (1.5, 0);
    \draw[popblue, thick] (-1, 0) -- (1.5, 0);
    % Rétine
    \draw[poppink, very thick] (1.5, -1.8) -- (1.5, 1.8) node[midway, right] {Rétine};
    % Image A'B' (inversée)
    \draw[popdark, thick] (1.5, -0.5) -- (1.5, 0.5) node[midway, right] {Image A'B'};
    % Foyer F'
    \fill[popred] (1.5, 0) circle (2pt) node[below] {$F'$};
    \draw[<->, popmuted] (-1, -2) -- (1.5, -2) node[midway, below] {Longueur axiale $\approx 24$ mm};
\end{tikzpicture}
\legende{Œil \textbf{emétrope} (normal) : L'image d'un objet lointain (rayons parallèles) se forme exactement sur la rétine ($F'$ confondu avec le plan rétinien).}
\end{popfigure}
\end{experience}

\begin{definition}[Accommodation]
Mécanisme permettant de voir net un objet \textbf{rapproché}. 
\begin{enumerate}
    \item \textbf{Stimulus :} Objet proche $\rightarrow$ rayons divergents.
    \item \textbf{Réflexe :} Contraction du \textbf{muscle ciliaire} $\rightarrow$ relâchement de la \textbf{zonule} (ligaments suspenseurs).
    \item \textbf{Effet :} Le cristallin, élastique, \textbf{se bombe} (augmentation de sa courbure) $\rightarrow$ sa \textbf{vergence augmente} (puissance max $\approx 33$ D).
    \item \textbf{Résultat :} Le foyer image $F'$ reste sur la rétine malgré la divergence des rayons incidents.
\end{enumerate}
\end{definition}

\begin{propriete}[Puissance d'accommodation et Point proche]
\begin{itemize}
    \item \textbf{Pouvoir d'accommodation ($P_{acc}$) :} Différence entre la vergence max (vision de près) et la vergence min (vision de loin, repos). $P_{acc} = V_{max} - V_{min}$ (en Dioptries, D).
    \item \textbf{Point proche (Punctum proximum) :} Distance minimale de vision nette. $d_{pp} = \frac{1}{P_{acc}}$ (en mètres).
    \item \textbf{Évolution avec l'âge :} $P_{acc}$ diminue inéluctablement (cristallin se sclérose/durcit) : $\approx 14$ D (10 ans) $\rightarrow$ $\approx 4$ D (40 ans) $\rightarrow$ $\approx 2,5$ D (50 ans) $\rightarrow$ $0$ D (60-65 ans). Le point proche recule.
\end{itemize}
\end{propriete}

\courssub{Anomalies de la vision (Ametropies)}
Ce sont des erreurs de réfraction : l'image ne se forme pas sur la rétine au repos (accommodation nulle pour la vision de loin). Elles se corrigent par des lunettes (verres sphériques).

\begin{popfigure}
\begin{tikzpicture}[scale=0.7, every node/.style={font=\small}]
    % MYOPIE
    \begin{scope}[xshift=-8cm]
        \draw[popdark, thick] (-5, 1.5) -- (-5, -1.5);
        \draw[->, popblue] (-5, 1.5) -- (-1, 1.5);
        \draw[->, popblue] (-5, -1.5) -- (-1, -1.5);
        \draw[->, popblue] (-5, 0) -- (-1, 0);
        \draw[poporange, very thick] (-1, -1) -- (-1, 1);
        \draw[popblue, thick] (-1, 1.5) -- (0.5, 0);
        \draw[popblue, thick] (-1, -1.5) -- (0.5, 0);
        \draw[popblue, thick] (-1, 0) -- (0.5, 0);
        \draw[poppink, very thick] (1.5, -1.5) -- (1.5, 1.5);
        \fill[popred] (0.5, 0) circle (2pt) node[below] {$F'$};
        \draw[popred, dashed] (0.5, 0) -- (1.5, 0);
        \node[align=center] at (-1, -2.5) {\textbf{MYOPIE}\\œil trop long / puissant\\Foyer \textbf{avant} rétine};
    \end{scope}
    % CORRECTION MYOPIE
    \begin{scope}[xshift=-8cm, yshift=-5cm]
        \draw[->, popblue] (-5, 1.5) -- (-2, 1.5);
        \draw[->, popblue] (-5, -1.5) -- (-2, -1.5);
        \draw[popgreen, very thick] (-2, -1) -- (-2, 1) node[midway, above, rotate=90] {Verge DIVERGENT ($V<0$)};
        \draw[->, popblue, dashed] (-2, 1.5) -- (-5, 3); % rayons virtuels
        \draw[->, popblue, dashed] (-2, -1.5) -- (-5, -3);
        \draw[->, popblue] (-2, 1.5) -- (0.5, 0);
        \draw[->, popblue] (-2, -1.5) -- (0.5, 0);
        \draw[->, popblue] (-2, 0) -- (0.5, 0);
        \draw[poporange, very thick] (-1, -1) -- (-1, 1);
        \draw[poppink, very thick] (1.5, -1.5) -- (1.5, 1.5);
        \fill[popred] (1.5, 0) circle (2pt) node[below] {$F'$};
        \node[align=center] at (-1, -2.5) {\textbf{Correction : Verre concave}\\Recule l'image virtuelle au Point Lointain};
    \end{scope}
    % HYPERMETROPIE
    \begin{scope}[xshift=4cm]
        \draw[popdark, thick] (-5, 1.5) -- (-5, -1.5);
        \draw[->, popblue] (-5, 1.5) -- (-1, 1.5);
        \draw[->, popblue] (-5, -1.5) -- (-1, -1.5);
        \draw[->, popblue] (-5, 0) -- (-1, 0);
        \draw[poporange, very thick] (-1, -1) -- (-1, 1);
        \draw[popblue, thick] (-1, 1.5) -- (2.5, 0);
        \draw[popblue, thick] (-1, -1.5) -- (2.5, 0);
        \draw[popblue, thick] (-1, 0) -- (2.5, 0);
        \draw[poppink, very thick] (1.5, -1.5) -- (1.5, 1.5);
        \fill[popred] (2.5, 0) circle (2pt) node[below] {$F'$};
        \draw[popred, dashed] (1.5, 0) -- (2.5, 0);
        \node[align=center] at (-1, -2.5) {\textbf{HYPERMÉTROPIE}\\œil trop court / faible\\Foyer \textbf{arrière} rétine};
    \end{scope}
    % CORRECTION HYPERMETROPIE
    \begin{scope}[xshift=4cm, yshift=-5cm]
        \draw[->, popblue] (-5, 1.5) -- (-2, 1.5);
        \draw[->, popblue] (-5, -1.5) -- (-2, -1.5);
        \draw[popgreen, very thick] (-2, -1) -- (-2, 1) node[midway, above, rotate=90] {Verge CONVERGENT ($V>0$)};
        \draw[->, popblue] (-2, 1.5) -- (0.5, 0);
        \draw[->, popblue] (-2, -1.5) -- (0.5, 0);
        \draw[->, popblue] (-2, 0) -- (0.5, 0);
        \draw[poporange, very thick] (-1, -1) -- (-1, 1);
        \draw[poppink, very thick] (1.5, -1.5) -- (1.5, 1.5);
        \fill[popred] (1.5, 0) circle (2pt) node[below] {$F'$};
        \node[align=center] at (-1, -2.5) {\textbf{Correction : Verre convexe}\\Avance le foyer sur la rétine};
    \end{scope}
\end{tikzpicture}
\legende{Schémas optiques des ametropies axiles et de leur correction par verres sphériques. En haut : Myopie (foyer antérieur) corrigée par verre divergent. En bas : Hypermétropie (foyer postérieur) corrigée par verre convergent.}
\end{popfigure}

\begin{aretenir}[Synthèse des anomalies de la vision de loin]
\begin{center}
\begin{tabularx}{\linewidth}{|l|X|X|X|}\hline
\rowcolor{popblueL} \textbf{Anomalie} & \textbf{Mécanisme optique} & \textbf{Cause anatomique principale} & \textbf{Correction (Verre)} \\\hline
\textbf{Myopie} & Foyer image \textbf{en avant} de la rétine. Vision nette de près, floue de loin. & Œil \textbf{trop long} (axile) ou \textbf{trop puissant} (cornée bombée). & \textbf{Divergent} (concave, $V < 0$).\\Recule l'image virtuelle au Point Lointain (PL). \\\hline
\textbf{Hypermétropie} & Foyer image \textbf{en arrière} de la rétine. Vision de loin possible par accommodation (fatigue), près difficile. & Œil \textbf{trop court} (axile) ou \textbf{pas assez puissant}. & \textbf{Convergent} (convexe, $V > 0$).\\Avance le foyer sur la rétine. \\\hline
\textbf{Presbytie} & \textbf{Perte du pouvoir d'accommodation} (cristallin sclérosé). Foyer de près \textbf{arrière} rétine. Vision de loin normale (émétrope). & \textbf{Vieillissement} du cristallin (perte élasticité). Inéluctable après 40-45 ans. & \textbf{Convergent pour la vision de près} (Addition $Add > 0$).\\Lunettes de près, bifocaux ou progressifs. \\\hline
\end{tabularx}
\end{center}
\end{aretenir}

\courssub{Maladies de la vision : Cataracte et Glaucome}
Contrairement aux ametropies, ce sont des \textbf{pathologies} (lésions tissulaires) entraînant une baisse de la vision non corrigeable par de simples verres.

\courssub{La Cataracte}
\begin{definition}[Cataracte]
Opacification partielle ou totale du \textbf{cristallin} (perte de transparence), le plus souvent liée à l'âge (cataracte sénile), mais possible dès la naissance (congénitale) ou secondaire (traumatisme, diabète, corticoïdes, UV).
\end{definition}

\begin{popfigure}
\begin{tikzpicture}[scale=1, every node/.style={font=\small}]
    % Œil normal
    \begin{scope}[xshift=-5cm]
        \draw[popdark, thick] (0,0) circle (2cm);
        \draw[popblue, thick, fill=popblue!10] (-2,0) arc (180:360:2cm and 0.6cm) -- cycle;
        \draw[poporange, thick, fill=poporange!20] (0,0) ellipse (0.7cm and 0.4cm);
        \draw[poppink, very thick] (1.2,1.5) arc (90:-90:0.4cm and 1.5cm);
        \node[align=center] at (0, -2.8){\textbf{Œil sain}\\Cristallin transparent};
        \draw[->, popblue] (-4, 1) -- (-2, 0.5);
        \draw[->, popblue] (-4, -1) -- (-2, -0.5);
        \draw[->, popblue] (-4, 0) -- (-2, 0);
        \draw[popblue] (-2, 0.5) -- (0.5, 0);
        \draw[popblue] (-2, -0.5) -- (0.5, 0);
        \draw[popblue] (-2, 0) -- (0.5, 0);
    \end{scope}
    % Œil cataracte
    \begin{scope}[xshift=5cm]
        \draw[popdark, thick] (0,0) circle (2cm);
        \draw[popblue, thick, fill=popblue!10] (-2,0) arc (180:360:2cm and 0.6cm) -- cycle;
        % Cristallin opaque (hachuré)
        \draw[poporange, thick, fill=poporange!50, pattern=north east lines, pattern color=popdark] (0,0) ellipse (0.7cm and 0.4cm);
        \draw[poppink, very thick] (1.2,1.5) arc (90:-90:0.4cm and 1.5cm);
        \node[align=center] at (0, -2.8){\textbf{Cataracte nucléaire}\\Cristallin opaque (jaunâtre)};
        % Diffusion lumière
        \draw[->, popblue] (-4, 1) -- (-2, 0.5);
        \draw[->, popblue] (-4, -1) -- (-2, -0.5);
        \draw[->, popblue] (-4, 0) -- (-2, 0);
        \draw[popblue, loosely dashed] (-2, 0.5) -- (0, 0.3) -- (1.5, 0.8);
        \draw[popblue, loosely dashed] (-2, -0.5) -- (0, -0.3) -- (1.5, -0.8);
        \draw[popblue, loosely dashed] (-2, 0) -- (0, 0.1) -- (1.5, 0.2);
        \draw[popblue, loosely dashed] (-2, 0) -- (0, -0.1) -- (1.5, -0.2);
        \node[popred, align=center] at (3, 0) {Diffusion\\de la lumière\\(Voile, éblouissement)};
    \end{scope}
\end{tikzpicture}
\legende{Principes optiques de la cataracte. Le cristallin opaque diffuse la lumière (rayons pointillés) : baisse de l'acuité visuelle (moins de lumière sur la rétine), voile lumineux (perte de contraste), éblouissement, altération des couleurs (jaunissement nucléaire).}
\end{popfigure}

\begin{propriete}[Prise en charge de la cataracte]
\begin{itemize}
    \item \textbf{Traitement unique : Chirurgie.} Aucune goutte ne rend la transparence.
    \item \textbf{Technique standard : Phacoémulsification.} Incision cornéenne minime ($\approx 2,2$ mm), fragmentation du cristallin opaque par ultrasons, aspiration.
    \item \textbf{Implant intraoculaire (IOL) :} Lentille artificielle souple (acrylique) placée dans le sac capsulaire (enveloppe du cristallin). Puissance calculée par biométrie (longueur axiale, courbure cornée). Récupération visuelle rapide (jours). Succès $> 95\%$.
    \item \textbf{Prévention :} Protection UV (lunettes cat. 3), arrêt tabac, équilibre glycémique (diabète), alimentation antioxydante (vit C, E, lutéine).
\end{itemize}
\end{propriete}

\courssub{Le Glaucome}
\begin{definition}[Glaucome chronique à angle ouvert (GCAO) - Forme la plus fréquente]
Neuropathie optique chronique, progressive, irréversible, caractérisée par une \textbf{excavation de la papille} (creusement de la tête du nerf optique) et une atteinte du champ visuel, \textbf{souvent associée à une Pression Intra-Oculaire (PIO) élevée} ($> 21$ mmHg).
\end{definition}

\begin{popfigure}
\begin{tikzpicture}[scale=1.2, every node/.style={font=\small}]
    % Fond d'œil normal
    \begin{scope}[xshift=-5cm]
        \draw[popdark, thick, fill=poporange!20] (0,0) circle (1.5cm); % Rétine fond orange
        % Papille normale
        \draw[poppurple, thick, fill=poppurple!30] (0.3, 0) circle (0.4cm);
        \draw[poppurple, thick] (0.3, 0) circle (0.15cm); % Excavation physiologique petite
        \node[align=center] at (0, -2.2){\textbf{Papille normale}\\Excavation/Diame < 0,3};
        \draw[->, popdark] (-2.5, 1.5) -- (-1, 0.8) node[left] {Artères/Veines};
        \draw[->, popdark] (-2.5, -1.5) -- (-1, -0.8) node[left] {Nerf optique};
    \end{scope}
    % Fond d'œil glaucomateux
    \begin{scope}[xshift=5cm]
        \draw[popdark, thick, fill=poporange!20] (0,0) circle (1.5cm);
        % Papille excavée
        \draw[poppurple, thick, fill=poppurple!30] (0.3, 0) circle (0.5cm); % Disque élargi
        \draw[poppurple, thick, fill=white] (0.3, 0) circle (0.35cm); % Grande excavation
        \node[align=center] at (0, -2.2){\textbf{Papille glaucomateuse}\\Excavation/Disque > 0,7\\(Amincissement anneau neurorétinien)};
        \draw[->, popdark] (2.5, 1.5) -- (1, 0.8) node[right, align=center]{Artères/Veines\\(déviation)};
        \draw[->, popdark] (2.5, -1.5) -- (1, -0.8) node[right, align=center]{Perte fibres\\(Scotomes)};
    \end{scope}
\end{tikzpicture}
\legende{Aspect du fond d'œil (papille optique). À gauche : papille saine, excavation physiologique centrale petite. À droite : glaucome, excavation élargie (creusement) due à la destruction des fibres nerveuses, anneau neurorétinien aminci.}
\end{popfigure}

\begin{propriete}[Physiopathologie et Traitement du Glaucome]
\begin{itemize}
    \item \textbf{Mécanisme :} PIO élevée ($\uparrow$ sécrétion humor aqueux ou $\downarrow$ évacuation par le trabéculum) $\rightarrow$ Contrainte mécanique + Ischémie sur la \textbf{lamelle cribreuse} (passage du nerf optique) $\rightarrow$ Compression/Apoptose des axones des cellules ganglionnaires $\rightarrow$ \textbf{Excavation papillaire} $\rightarrow$ \textbf{Perte du champ visuel} (périphérique d'abord $\rightarrow$ vision tubulaire $\rightarrow$ cécité centrale).
    \item \textbf{Caractère insidieux :} Indolore, vision centrale conservée tardivement $\rightarrow$ \textbf{Dépistage indispensable} (Tonométrie, Fond d'œil, Champ visuel) après 40 ans.
    \item \textbf{Traitement (Objectif : Baisser PIO $\approx 20-30\%$) :}
    \begin{enumerate}
        \item \textbf{Médical (1ère ligne) : Collyres hypotonisants quotidiens à vie.}
        \begin{itemize}
            \item Analogues des prostaglandines (Latanoprost, Travoprost) : $\uparrow$ débit uvéo-scléral. (1 fois/jour, soir).
            \item Bêta-bloquants (Timolol) : $\downarrow$ sécrétion aqueuse. (Attention asthme, cœur).
            \item Inhibiteurs anhydrase carbonique (Dorzolamide) : $\downarrow$ sécrétion.
            \item Alpha-2 agonistes (Brimonidine) : $\downarrow$ sécrétion + $\uparrow$ débit.
        \end{itemize}
        \item \textbf{Laser :} Trabéculoplastie (SLT) pour faciliter l'écoulement.
        \item \textbf{Chirurgie :} Trabéculectomie (fistule de filtration) ou implants de drainage (si échec médical/laser).
    \end{enumerate}
\end{itemize}
\end{propriete}

\courssub{Hygiène de la vision}
\begin{methode}[Bonnes pratiques pour préserver sa vision]
\begin{enumerate}
    \item \textbf{Éclairage :} Suffisant ($\ge 300$ lux lecture, $\ge 500$ lux travail de précision/couture). Lampe orientée (gauche pour droitier), sans éblouissement direct. Lumière naturelle privilégiée.
    \item \textbf{Distances et Posture :} Lecture/écran $\approx 35$--$40$ cm. Dos droit. Écran au niveau des yeux (haut de l'écran $\approx$ niveau regard).
    \item \textbf{Règle 20-20-20 (Écrans) :} Toutes les \textbf{20 minutes}, regarder à \textbf{20 pieds ($\approx 6$ m)} pendant \textbf{20 secondes}. Cligner des yeux volontairement (réhydratation cornée).
    \item \textbf{Protection UV :} Lunettes solaires catégorie 3 (norme CE) en extérieur (réverbération sable, eau, neige, altitude). Protège cornée, cristallin (cataracte), rétine (DMLA).
    \item \textbf{Nutrition :} Fruits/légumes colorés (Vit A, C, E, $\beta$-carotène, lutéine, zéaxanthine : épinards, choux, carottes, mangues, papayes). Poisson (Oméga-3). Hydratation.
    \item \textbf{Contrôle médical :} \textbf{Annuel après 50 ans} (PIO, Fond d'œil, Acuité visuelle). Dépistage glaucome, cataracte, DMLA, rétinopathie diabétique/hypertensive. Enfant : dépistage scolaire (amblyopie, strabisme, myopie).
    \item \textbf{Éviter :} Tabac (facteur majeur cataracte, DMLA, glaucome), frottement yeux (kératocône), automédication collyres (corticoïdes $\rightarrow$ PIO $\uparrow$).
\end{enumerate}
\end{methode}

\courssub{Activité d'intégration : Consultation au CSI de Mvog-Ada (Yaoundé)}

\courssub{Situation professionnelle simulée}
\label{sit:consultation}
\textbf{Contexte :} Vous êtes stagiaire assistant(e) au \textbf{Centre de Santé Intégré (CSI) de Mvog-Ada}, à Yaoundé. L'ophtalmologiste de permanence, le Dr \textsc{Atangana}, vous confie l'analyse préliminaire de trois dossiers patients reçus ce matin lors de la consultation « Santé visuelle scolaire et professionnelle ». Votre mission : examiner les documents fournis, poser un diagnostic argumenté pour chaque cas et préparer la fiche de synthèse que le médecin validera avant la prescription.

\vspace{0.3cm}
\noindent\textbf{Dossier n°1 : \textsc{Ekoué}, 13 ans, élève en classe de 4\textsuperscript{e} au Lycée de Nkolbisson.}
\begin{itemize}
    \item \textbf{Anamnèse :} L'élève se plaint de « ne pas voir le tableau » depuis la rentrée. Il doit se lever et avancer jusqu'au premier rang pour copier la leçon. En revanche, il lit son cahier et son manuel « sans problème » à une distance habituelle ($\approx 30$ cm). Il signale des maux de tête en fin de journée de cours.
    \item \textbf{Examen clinique :} Acuité visuelle de loin (échelle de Monoyer) : $AV_{loin} = 3/10$ à l'œil droit (OD) et $3/10$ à l'œil gauche (OG). Acuité visuelle de près (texte Parinaud) : $AV_{près} = P2$ (normal) aux deux yeux.
    \item \textbf{Document 1 (Schéma optique simplifié de l'œil d'Ékoué) :} 
\end{itemize}

\begin{popfigure}
\begin{tikzpicture}[scale=0.8, every node/.style={font=\small}]
    % Œil myope
    \draw[popdark, thick] (-4, 1.5) -- (-4, -1.5) node[midway, left] {Objet loin ($\infty$)};
    \draw[->, popblue, thick] (-4, 1.5) -- (-1, 1.5);
    \draw[->, popblue, thick] (-4, -1.5) -- (-1, -1.5);
    \draw[->, popblue, thick] (-4, 0) -- (-1, 0);
    \draw[poporange, very thick] (-1, -1) -- (-1, 1) node[midway, above, rotate=90] {Cornée + Cristallin};
    % Foyer AVANT rétine
    \draw[popblue, thick] (-1, 1.5) -- (0.2, 0);
    \draw[popblue, thick] (-1, -1.5) -- (0.2, 0);
    \draw[popblue, thick] (-1, 0) -- (0.2, 0);
    \fill[popred] (0.2, 0) circle (2pt) node[below] {$F'$};
    % Rétine plus loin
    \draw[poppink, very thick] (1.5, -1.5) -- (1.5, 1.5) node[midway, right] {Rétine};
    % Cercle de diffusion sur rétine
    \draw[popred, thick, fill=popred!20] (1.5, -0.4) -- (1.5, 0.4);
    \node[popred, align=center] at (2.5, 0) {Cercle de diffusion\\(Image floue)};
    \node[align=center] at (-1, -2.5) {\textbf{Myopie : Foyer $F'$ antérieur à la rétine}};
\end{tikzpicture}
\legende{Document 1 : Schéma optique de l'œil myope d'Ékoué. Les rayons parallèles (objet à l'infini) convergent en un foyer $F'$ situé \textbf{en avant de la rétine}. La rétine reçoit une tache lumineuse floue (cercle de diffusion).}
\end{popfigure}

\vspace{0.3cm}
\noindent\textbf{Dossier n°2 : \textsc{Mme Essomba}, 52 ans, couturière au Marché Central.}
\begin{itemize}
    \item \textbf{Anamnèse :} Depuis environ deux ans, elle doit « écarter le bras au maximum » pour enfiler une aiguille ou lire les mesures sur son mètre-ruban. Elle rapproche le texte de ses yeux pour le voir net, mais au-delà de $40$--$50$ cm, la vision de près devient floue. Elle ressent une fatigue oculaire (picotements, larmoiement) et des céphalées après deux heures de couture continue. Elle ne porte pas de lunettes. Vision de loin : bonne (reconnaît les visages de l'autre côté de la rue).
    \item \textbf{Examen clinique :} $AV_{loin} = 10/10$ (OD et OG). $AV_{près} = P8$ (insuffisant) à la distance habituelle de lecture ($33$ cm). Amélioration nette à $P2$ en écartant le texte à $50$ cm (distance de près « confortable » pour elle).
    \item \textbf{Document 2 (Données physiologiques) :} Pouvoir d'accommodation maximal mesuré à $2,50$ D (dioptries) (norme pour 50 ans : $\approx 2,50$ D). Distance du point proche (punctum proximum) : $40$ cm.
\end{itemize}

\begin{popfigure}
\begin{tikzpicture}[scale=0.8, every node/.style={font=\small}]
    % Presbytie : Objet près, foyer arrière rétine
    \draw[popdark, thick] (-1.5, 1) -- (-1.5, -1) node[midway, left] {Objet près (33 cm)};
    \draw[->, popblue, thick] (-1.5, 1) -- (-1, 0.5);
    \draw[->, popblue, thick] (-1.5, -1) -- (-1, -0.5);
    \draw[->, popblue, thick] (-1.5, 0) -- (-1, 0);
    \draw[poporange, very thick] (-1, -1) -- (-1, 1) node[midway, above, rotate=90] {Cristallin vieilli (peu puissant)};
    % Foyer ARRIÈRE rétine
    \draw[popblue, thick] (-1, 0.5) -- (2, 0);
    \draw[popblue, thick] (-1, -0.5) -- (2, 0);
    \draw[popblue, thick] (-1, 0) -- (2, 0);
    \fill[popred] (2, 0) circle (2pt) node[below] {$F'$};
    % Rétine
    \draw[poppink, very thick] (1.5, -1.5) -- (1.5, 1.5) node[midway, right] {Rétine};
    \draw[popred, dashed] (1.5, 0) -- (2, 0);
    \node[popred, align=center] at (2.8, 0) {Foyer \textbf{arrière} rétine};
    \node[align=center] at (-1, -2.5) {\textbf{Presbytie : Accommodation insuffisante}};
\end{tikzpicture}
\legende{Document 2 (complément) : Schéma optique de la presbytie. Pour un objet proche (33 cm), le cristallin vieilli (faible vergence max) ne converge pas assez : le foyer $F'$ se forme \textbf{en arrière de la rétine}. En écartant l'objet (50 cm), les rayons sont moins divergents, le foyer recule sur la rétine.}
\end{popfigure}

\vspace{0.3cm}
\noindent\textbf{Dossier n°3 : \textsc{M. Fouda}, 63 ans, retraité, ancien agent de la SONARA (Limbe), et \textsc{M. Njoya}, 61 ans, voisin de chambre.}
\begin{itemize}
    \item \textbf{Cas 3a (M. Fouda) :} Vision voilée progressive depuis 18 mois, « comme derrière une cascade » ou un verre dépoli. Éblouissement intense la nuit face aux phares des voitures (conduite difficile). Perception des couleurs ternes (le vert du drapeau camerounais lui paraît jaunâtre).
    \item \textbf{Examen à la lampe à fente (Biomicroscopie) :} Opacification diffuse du noyau du cristallin (cataracte nucléaire). Réflexe rouge au fond d'œil atténué.
    \item \textbf{Cas 3b (M. Njoya) :} Pas de plainte visuelle spontanée au début. Découvert lors d'un dépistage : réduction du champ visuel périphérique (vision tubulaire naissante). Tonométrie à l'air (sans contact) : Pression Intra-Oculaire (PIO) = $24$ mmHg (OD) et $26$ mmHg (OG) [Norme : $10$--$21$ mmHg]. Examen du fond d'œil : excavation du nerf optique (creusement de la papille) avec rapport excavation/disque $> 0,7$.
\end{itemize}

\begin{popfigure}
\begin{tikzpicture}[scale=0.9, every node/.style={font=\small}]
    % Cataracte M. Fouda
    \begin{scope}[xshift=-6cm]
        \draw[popdark, thick] (0,0) circle (1.8cm);
        \draw[popblue, thick, fill=popblue!10] (-1.8,0) arc (180:360:1.8cm and 0.5cm) -- cycle;
        \draw[poporange, thick, fill=poporange!60, pattern=north west lines, pattern color=popdark] (0,0) ellipse (0.6cm and 0.35cm);
        \draw[poppink, very thick] (1.0,1.3) arc (90:-90:0.35cm and 1.3cm);
        \node[align=center] at (0, -2.5){\textbf{M. Fouda : Cataracte nucléaire}\\Cristallin opaque / jauni};
        \draw[->, popblue] (-3.5, 0.8) -- (-1.8, 0.4);
        \draw[->, popblue] (-3.5, -0.8) -- (-1.8, -0.4);
        \draw[popblue, loosely dashed] (-1.8, 0.4) -- (0, 0.2) -- (1.5, 0.6);
        \draw[popblue, loosely dashed] (-1.8, -0.4) -- (0, -0.2) -- (1.5, -0.6);
        \node[popred, align=center] at (2.5, 0) {Diffusion / Voile\\Éblouissement / Couleurs ternes};
    \end{scope}
    % Glaucome M. Njoya
    \begin{scope}[xshift=6cm]
        \draw[popdark, thick, fill=poporange!20] (0,0) circle (1.5cm);
        \draw[poppurple, thick, fill=poppurple!30] (0.3, 0) circle (0.5cm);
        \draw[poppurple, thick, fill=white] (0.3, 0) circle (0.38cm);
        \node[align=center] at (0, -2.5){\textbf{M. Njoya : Glaucome (GCAO)}\\Excavation/Disque > 0,7\\PIO = 24-26 mmHg};
        \draw[->, popdark] (-2.5, 1.2) -- (-0.5, 0.6) node[left, align=center]{Champ visuel\\réduit (périphérie)};
        \draw[->, popdark] (-2.5, -1.2) -- (-0.5, -0.6) node[left] {Vision tubulaire};
    \end{scope}
\end{tikzpicture}
\legende{Documents 3a et 3b : Aspects cliniques schématiques. Gauche : Cataracte nucléaire (opacité centrale jaunâtre diffusant la lumière). Droite : Papille glaucomateuse (excavation majeure, amincissement de l'anneau neurorétinien) responsable de la constriction du champ visuel.}
\end{popfigure}

\courssub{Consignes de l'activité}
Travailler par groupes de 3--4 élèves. Pour chaque dossier, répondre aux questions ci-dessous en rédigeant des phrases complètes, justifiées par les connaissances du cours (sections 1 à 5). Une mise en commun orale suivra pour établir le tableau de synthèse final.

\vspace{0.2cm}
\noindent\textbf{Dossier 1 (Ékoué) :}
\begin{enumerate}
    \item À quel trouble de la vision correspond ce tableau clinique ? Justifier par les symptômes, l'acuité visuelle et le schéma optique (Document 1).
    \item Expliquer le mécanisme optique : pourquoi l'image se forme-t-elle en avant de la rétine ? Citer deux causes anatomiques possibles.
    \item Quelle correction optique proposez-vous ? Préciser la nature du verre (vergence, forme) et son rôle sur le trajet des rayons lumineux. Calculer la vergence approximative du verre correcteur si l'on considère que le point lointain de l'élève est à $33$ cm.
    \item Rédiger deux conseils d'hygiène visuelle adaptés à un collégien (écrans, éclairage, pauses, activités extérieures).
\end{enumerate}

\vspace{0.2cm}
\noindent\textbf{Dossier 2 (Mme Essomba) :}
\begin{enumerate}
    \item Identifier le trouble. Justifier par l'âge, les symptômes (distance de lecture), l'acuité visuelle de loin et de près, et la donnée sur le pouvoir d'accommodation (Document 2).
    \item Expliquer le mécanisme physiologique : rôle du cristallin et de l'accommodation. Pourquoi ce trouble apparaît-il inéluctablement avec l'âge ?
    \item Proposer la correction adaptée. Préciser la nature du verre, sa vergence (signe) et son effet sur l'image. Pourquoi parle-t-on d'« addition » pour la vision de près ? Calculer l'addition nécessaire pour lire à $33$ cm.
    \item Conseils de prévention : fatigue visuelle, éclairage du poste de travail (couture), contrôle régulier.
\end{enumerate}

\vspace{0.2cm}
\noindent\textbf{Dossier 3 (M. Fouda \& M. Njoya) :}
\begin{enumerate}
    \item Différencier les deux pathologies : cataracte (M. Fouda) et glaucome chronique à angle ouvert (M. Njoya). Pour chacune, citer le signe clinique majeur, la lésion anatomique et la conséquence fonctionnelle.
    \item Pour la cataracte : expliquer l'origine de la baisse d'acuité visuelle et de l'éblouissement (optique). Décrire le traitement de référence (chirurgie) et le principe de l'implant intraoculaire.
    \item Pour le glaucome : expliquer le lien entre PIO élevée, excavation de la papille et réduction du champ visuel. Pourquoi le dépistage est-il crucial ? Citer les principales classes de collyres hypotonisants.
    \item Conseil de prévention commun : importance du contrôle ophtalmologique annuel après 50 ans (dépistage glaucome, suivi cataracte, DMLA, rétinopathies).
\end{enumerate}

\courssub{Ressources à mobiliser (Rappels de cours)}
\begin{itemize}
    \item \textbf{Anatomie :} Cornée ($\approx 43$ D), Cristallin (variable $\approx 19$--$33$ D), Rétine (plan récepteur), Humor aqueux/vitré, Nerf optique, Muscle ciliaire, Zonule.
    \item \textbf{Optique :} Lois de la réfraction. Condition de netteté : Foyer image $F'$ sur la rétine. Vergence $V = 1/f$ (Dioptrie, $f$ en mètres). Verre divergent ($V<0$, concave) $\rightarrow$ écarte les rayons. Verge convergent ($V>0$, convexe) $\rightarrow$ rapproche les rayons.
    \item \textbf{Accommodation :} Contraction muscle ciliaire $\rightarrow$ relâchement zonule $\rightarrow$ bombement cristallin $\rightarrow$ $\uparrow$ Vergence. $P_{acc} = V_{max} - V_{min}$. $d_{pp} = 1/P_{acc}$.
    \item \textbf{Corrections :} Myopie : Verge divergent (recule image virtuelle au Point Lointain). Hypermétropie : Verge convergent (avance foyer). Presbytie : Verge convergent pour près (Addition $Add > 0$).
    \item \textbf{Pathologies :} Cataracte = Opacification cristallin (diffusion lumière) $\rightarrow$ Chirurgie (Phaco + Implant). Glaucome = Neuropathie optique (souvent PIO $\uparrow$) $\rightarrow$ Excavation papille $\rightarrow$ Perte champ visuel $\rightarrow$ Collyres hypotonisants $\pm$ Laser/Chirurgie.
    \item \textbf{Hygiène :} Distance lecture $35$--$40$ cm, Éclairage $\ge 300$ lux, Règle 20-20-20, Protection UV cat. 3, Alimentation antioxydants, Non-tabac, Contrôle annuel $> 50$ ans.
\end{itemize}

\courssub{Démarche et corrigé commenté}
\begin{exoresolu}[Corrigé commenté de l'activité d'intégration]
\textbf{DOSSIER 1 : Ékoué (13 ans) -- Myopie (probablement axile)}

\textbf{1. Identification et diagnostic.}
Ékoué présente une \textbf{myopie bilatérale}. Les arguments convergent :
\begin{itemize}
    \item \textbf{Symptômes :} Vision floue de loin (tableau $> 3$ m) conservée de près (cahier à $30$ cm). Céphalées de fin de journée (fatigue oculaire / sur-accommodation vaine).
    \item \textbf{Acuité visuelle :} $AV_{loin} = 3/10$ (baisse nette bilatérale) ; $AV_{près} = P2$ (normale). \emph{Caractère distinctif :} le myope voit net de près sans correction.
    \item \textbf{Schéma (Doc 1) :} Les rayons parallèles (objet à l'infini) convergent en un foyer $F'$ situé \textbf{en avant de la rétine}. La rétine reçoit un cercle de diffusion $\rightarrow$ image floue.
\end{itemize}

\textbf{2. Mécanisme optique et causes anatomiques.}
L'œil myope est \textbf{trop puissant pour sa longueur}, ou \textbf{trop long pour sa puissance}. Deux causes anatomiques principales (souvent associées) :
\begin{itemize}
    \item \textbf{Myopie axile (la plus fréquente à l'adolescence) :} Allongement excessif du globe oculaire (longueur axiale $> 24$ mm). La rétine est « trop loin » du système optique (cornée + cristallin).
    \item \textbf{Myopie d'indice / de courbure :} Augmentation de la puissance de la cornée (courbure excessive / kératocône) ou du cristallin (indice de réfraction élevé, ex. début de cataracte nucléaire, rare à 13 ans).
\end{itemize}
\emph{Conséquence optique :} Le système oculaire converge trop fort $\rightarrow$ foyer antérieur à la rétine.

\textbf{3. Correction optique.}
\begin{itemize}
    \item \textbf{Principe :} Placer un \textbf{verge divergent (concave, vergence $V < 0$)} devant l'œil. Ce verre crée une \textbf{image virtuelle} de l'objet lointain \textbf{au point lointain (PL)} de l'œil myope (le point le plus loin vu net sans effort d'accommodation). L'œil n'a alors plus qu'à regarder cet objet virtuel placé à son PL.
    \item \textbf{Calcul de la vergence (approximatif) :} Le point lointain d'Ékoué est estimé à $33$ cm ($0,33$ m) car il voit net son cahier à cette distance sans accommodation excessive. La vergence du verre correcteur $V_{verre}$ doit former l'image d'un objet à l'infini ($u = \infty$) au PL ($v' = -0,33$ m, signe $-$ car image virtuelle du même côté que l'objet pour un verre divergent).
    \[ V_{verre} = \frac{1}{v'} - \frac{1}{u} = \frac{1}{-0,33} - 0 \approx -3,03 \text{ D} \]
    On prescrira des verres \textbf{$\approx -3,00$ D} (ou $-3,25$ D selon la réfraction subjective sous cycloplégie pour relâcher l'accommodation résiduelle).
    \item \textbf{Effet sur le schéma :} Le verre divergent « écarte » les rayons incidents avant qu'ils n'entrent dans l'œil. Le système cornée-cristallin, moins sollicité, forme alors l'image \textbf{sur la rétine}.
\end{itemize}

\textbf{4. Conseils d'hygiène visuelle (Collégien, contexte camerounais).}
\begin{enumerate}
    \item \textbf{Règle du 20-20-20 :} Toutes les 20 minutes de travail de près (cahier, téléphone, tablette), regarder un objet à au moins 20 pieds ($\approx 6$ m) pendant 20 secondes. Relâche l'accommodation et la convergence, prévient la fatigue et la progression de la myopie.
    \item \textbf{Éclairage, posture et plein air :} Travailler sous un éclairage suffisant ($\ge 300$ lux, lampe de bureau orientée à gauche pour un droitier), distance de lecture $\approx 35$--$40$ cm (ne pas coller le nez au cahier), dos droit. Limiter le temps d'écran récréatif ($< 2$ h/jour hors cours). \textbf{Privilégier les activités en plein air} (cour de récréation, sport) : la lumière naturelle (dopamine rétinienne) est protectrice contre l'allongement axial (progression myopie).
\end{enumerate}

\vspace{0.5cm}
\noindent\textbf{DOSSIER 2 : Mme Essomba (52 ans) -- Presbytie}

\textbf{1. Identification et diagnostic.}
Mme Essomba présente une \textbf{presbytie} (débutante/modérée).
\begin{itemize}
    \item \textbf{Âge :} 52 ans (pic de survenue : 40--50 ans).
    \item \textbf{Symptômes :} Allongement de la distance de lecture (bras tendu $\approx 50$ cm), fatigue visuelle de près (asthénopie : picotements, larmoiement, céphalées), vision de loin intacte ($10/10$).
    \item \textbf{Acuité visuelle :} $AV_{loin} = 10/10$ (émétropie de loin). $AV_{près} = P8$ à $33$ cm (insuffisant) ; $P2$ à $50$ cm (son point proche actuel).
    \item \textbf{Donnée physiologique (Doc 2) :} Pouvoir d'accommodation résiduel $P_{acc} = 2,50$ D. Distance du point proche $d_{pp} = 1/2,50 = 0,40$ m $= 40$ cm. Cela correspond exactement à sa plainte : elle ne peut plus mettre au net plus près que $40$--$50$ cm.
\end{itemize}

\textbf{2. Mécanisme physiologique.}
L'accommodation repose sur la \textbf{déformation du cristallin} : contraction du muscle ciliaire $\rightarrow$ relâchement de la zonule $\rightarrow$ cristallin devient plus bombé (augmentation de sa vergence). Avec l'âge, le cristallin \textbf{sclérose (durcit)} : il perd son élasticité, sa capacité à se déformer diminue. Le pouvoir d'accommodation $P_{acc}$ chute inéluctablement ($\approx 14$ D à 10 ans $\rightarrow$ $2,5$ D à 50 ans $\rightarrow$ $0$ D vers 60--65 ans). Le point proche recule. L'œil devient incapable d'augmenter sa puissance pour former l'image d'un objet proche sur la rétine : l'image se forme \textbf{en arrière de la rétine} (mécanisme optique inverse de la myopie, mais dû à une puissance insuffisante \emph{variable}).

\textbf{3. Correction optique.}
\begin{itemize}
    \item \textbf{Verge convergent (convexe, $V > 0$)} pour la vision de près uniquement (lunettes de près) ou en partie basse de verres progressifs.
    \item \textbf{Rôle :} Le verre convergent « pré-converge » les rayons venant de l'objet proche (ex. à $33$ cm). L'œil presbyte, qui ne peut plus accommoder suffisamment, reçoit des rayons déjà convergents et forme l'image \textbf{sur la rétine}.
    \item \textbf{Calcul de l'addition (Add) :} Pour lire à $33$ cm ($0,33$ m), l'œil a besoin d'une vergence d'incidence de $1/0,33 \approx 3,00$ D. Son pouvoir d'accommodation max est $2,50$ D. L'addition nécessaire est $Add \approx 3,00 - 2,50 = +0,50$ D. On arrondit au quart de dioptrie supérieur pour le confort : \textbf{$+0,75$ D ou $+1,00$ D}.
    \item \textbf{Terminologie :} On parle d'« \textbf{addition} » car cette vergence positive s'\textbf{ajoute} à la correction de loin (ici nulle, $0,00$ D) pour la vision de près.
\end{itemize}

\textbf{4. Conseils (Couturière, Marché Central).}
\begin{enumerate}
    \item \textbf{Éclairage ciblé :} Utiliser une lampe articulée (type lampe d'architecte) dirigée sur la zone de couture (aiguille, tissus), puissance $\ge 500$ lux, température de couleur neutre ($4000$ K) pour bien discriminer les couleurs des fils. Éviter l'éblouissement direct.
    \item \textbf{Pauses actives :} Toutes les 45--60 min, lever les yeux, regarder loin (fenêtre), cligner délibérément (réhydratation cornée), étirer nuque/épaules. Prévenir le syndrome sec oculaire aggravé par la concentration (baisse fréquence clignement).
    \item \textbf{Suivi :} Contrôle ophtalmologique tous les 18--24 mois : évolution de l'addition (augmente $\approx +0,50$ D tous les 2--3 ans), dépistage glaucome (PIO, fond d'œil), cataracte, DMLA.
\end{enumerate}

\vspace{0.5cm}
\noindent\textbf{DOSSIER 3 : M. Fouda (Cataracte) \& M. Njoya (Glaucome)}

\textbf{1. Différenciation des deux pathologies.}

\begin{center}
\begin{tabularx}{\linewidth}{|l|X|X|}\hline
\rowcolor{popblueL} \textbf{Critère} & \textbf{M. Fouda : Cataracte (Sénile Nucléaire)} & \textbf{M. Njoya : Glaucome Chronique à Angle Ouvert (GCAO)} \\\hline
\textbf{Signe clinique majeur} & Baisse AV progressive, voilée, \textbf{éblouissement} (phares nuit), couleurs ternes (jaunissement noyau). & \textbf{Asymptomatique} au début. Découvert : \textbf{champ visuel réduit} (périphérique), PIO $\uparrow$. \\\hline
\textbf{Lésion anatomique} & \textbf{Opacification du cristallin} (perte transparence, protéines dénaturées, jaunissement noyau). & \textbf{Neuropathie optique} : destruction fibres nerveuses ganglionnaires $\rightarrow$ \textbf{excavation papille} (creusement disque optique). \\\hline
\textbf{Conséquence fonctionnelle} & Baisse \textbf{quantitative} (AV) et \textbf{qualitative} (contraste, couleurs) \textbf{réversible} par chirurgie. & Perte \textbf{irréversible} du champ visuel (périphérique $\rightarrow$ central = vision tubulaire $\rightarrow$ cécité). \\\hline
\textbf{Paramètre clé examen} & Biomicroscopie : opacité cristallin. Fond d'œil : réflexe rouge atténué. & \textbf{PIO} $> 21$ mmHg (facteur de risque majeur). Fond d'œil : excavation $> 0,7$. Champ visuel : scotomes. \\\hline
\end{tabularx}
\end{center}

\textbf{2. Cataracte (M. Fouda) : Mécanisme et Traitement.}
\begin{itemize}
    \item \textbf{Optique :} Le cristallin opaque \textbf{diffuse} la lumière (diffusion sur agrégats protéiques). Deux effets : (1) Moins de lumière atteint la rétine $\rightarrow$ baisse AV (image sombre). (2) Lumière parasite diffusée dans l'œil $\rightarrow$ \textbf{voile lumineux} $\rightarrow$ perte de contraste, éblouissement incapacitant (phares nuit). Jaunissement noyau $\rightarrow$ filtration bleu $\rightarrow$ couleurs ternes/jaunâtres.
    \item \textbf{Traitement de référence :} \textbf{Phacoémulsification} (ultrasons) + \textbf{Implant Intraoculaire (IOL)}. Sous anesthésie locale (collyre/injection), incision cornéenne minime ($\approx 2,2$ mm). Le cristallin opaque est fragmenté/aspiré. Le sac capsulaire (enveloppe) est conservé. Un implant souple (acrylique) de puissance calculée (biométrie : longueur axiale, kératométrie) est injecté dans le sac. Récupération visuelle rapide (J1--J7). Succès $> 95\%$.
\end{itemize}

\textbf{3. Glaucome (M. Njoya) : Physiopathologie, Dépistage, Traitement.}
\begin{itemize}
    \item \textbf{Physiopathologie :} PIO élevée ($24$--$26$ mmHg) $\rightarrow$ contrainte mécanique + ischémie sur la \textbf{lamelle cribreuse} (zone de passage du nerf optique dans la sclère) $\rightarrow$ compression/écrasement des axones des cellules ganglionnaires $\rightarrow$ mort cellulaire (apoptose) $\rightarrow$ \textbf{excavation de la papille} (creusement, rapport Excavation/Disque $> 0,7$ pathologique) $\rightarrow$ perte correspondante du \textbf{champ visuel} (fibres périphériques d'abord $\rightarrow$ scotome arcué $\rightarrow$ vision tubulaire).
    \item \textbf{Dépistage crucial :} Maladie \textbf{silencieuse} (pas de douleur, AV centrale conservée tardivement). Dépistage systématique après 40 ans (tonométrie, fond d'œil, champ visuel si suspect) permet de traiter \textbf{avant} perte irréversible. Au Cameroun, dépistage intégré dans les CSI (ex. Journée Mondiale de la Vue, programmes ONG/Ministère).
    \item \textbf{Traitement médical (1\textsuperscript{re} ligne) :} Collyres hypotonisants (visent $\downarrow$ PIO de $20$--$30\%$).
    \begin{itemize}
        \item \textbf{Analogues des prostaglandines} (Latanoprost, Travoprost, 1x/jour soir) : $\uparrow$ débit uvéo-scléral. \emph{Souvent 1\textsuperscript{re} intention.}
        \item \textbf{Bêta-bloquants} (Timolol, 2x/jour) : $\downarrow$ sécrétion humor aqueux. Contre-indication : asthme, bradycardie.
        \item \textbf{Inhibiteurs anhydrase carbonique} (Dorzolamide, Brinzolamide, 2--3x/jour) : $\downarrow$ sécrétion.
        \item \textbf{Alpha-2 agonistes} (Brimonidine, 2x/jour) : $\downarrow$ sécrétion + $\uparrow$ débit uvéo-scléral. Attention enfant $< 2$ ans.
    \end{itemize}
    Si échec médical $\rightarrow$ Laser (Trabéculoplastie SLT) $\rightarrow$ Chirurgie (Trabéculectomie, Implants drainage).
\end{itemize}

\textbf{4. Conseil de prévention commun (Senior 60+).}
\textbf{« Après 50 ans, une visite ophtalmologique complète tous les 12 à 18 mois, même sans plainte. »} Justification : Dépistage précoce du glaucome (seule façon d'éviter la cécité), suivi de l'évolution de la cataracte (décider le moment opératoire optimal), dépistage de la DMLA (dégénérescence maculaire liée à l'âge), contrôle de la presbytie (adaptation addition), dépistage rétinopathie diabétique/hypertensive (fréquentes au Cameroun). Hygiène de vie : arrêt tabac (facteur risque cataracte, DMLA, glaucome), équilibre glycémique/lipidique/tensionnel, protection UV (lunettes solaires cat. 3), alimentation riche légumes verts (lutéine/zéaxanthine), fruits colorés (vit C/E), poisson (oméga-3).
\end{exoresolu}

\courssub{Bilan de l'activité et Synthèse de la leçon}
Cette activité d'intégration a permis de mettre en évidence les points clés suivants, validant l'acquisition des compétences visées par le programme de 3ème :

\begin{enumerate}
    \item \textbf{Maîtrise du raisonnement clinique :} Les élèves ont appris à croiser des données hétérogènes (anamnèse, acuité visuelle, schémas optiques, mesures biométriques/pressionnelles) pour poser un diagnostic différentiel structuré. La distinction entre \textbf{trouble réfractif} (myopie, hypermétropie, presbytie : erreur de mise au point corrigible par verres) et \textbf{pathologie oculaire} (cataracte, glaucome : lésion tissulaire nécessitant traitement médical/chirurgical) est fondamentale.
    \item \textbf{Lien Anatomie-Fonction-Pathologie :} L'explication des mécanismes (foyer antérieur/postérieur à la rétine, sclérose cristallinienne, diffusion lumineuse, excavation papillaire) ancre la compréhension dans l'anatomie (cornée, cristallin, rétine, nerf optique, humor aqueux) et l'optique géométrique (vergence, foyer, accommodation).
    \item \textbf{Compétence « Prescription/Conseil » :} Le calcul approximatif des vergences (myopie $-3,00$ D, presbytie $Add +0,75$ D) et le choix justifié du type de verre (divergent/convergent, unifocal/progressif) ou du traitement (collyres, chirurgie implant) développent l'esprit critique et l'application concrète des formules ($V = 1/f$).
    \item \textbf{Dimension de Santé Publique (Contexte Cameroun) :} L'activité ancre les savoirs dans la réalité locale : consultation en CSI, métier de couturière (vision de près), conduite nocturne (éblouissement cataracte), dépistage glaucome (maladie silencieuse, cécité évitable), coût accessibilité des implants/lunettes, rôle des agents de santé communautaire. L'hygiène visuelle (écrans, éclairage, UV, tabac, nutrition) devient un message de prévention primaire/secondaire que l'élève peut diffuser dans sa famille.
    \item \textbf{Qualité de la communication scientifique :} La rédaction des réponses (diagnostic, mécanisme, conseil) oblige à utiliser un vocabulaire précis (\emph{accommodation, point lointain, addition, excavation, pression intra-oculaire, phacoémulsification}) et à structurer l'argumentaire (Observation $\rightarrow$ Interprétation $\rightarrow$ Conclusion), compétence transversale évaluée au BEPC.
\end{enumerate}

\noindent\textbf{Tableau de synthèse final (à construire collectivement au tableau) :}

\begin{center}
\begin{tabularx}{\linewidth}{|l|X|X|X|}\hline
\rowcolor{popblueL} \textbf{Trouble / Maladie} & \textbf{Mécanisme principal} & \textbf{Correction / Traitement} & \textbf{Prévention / Hygiène clé} \\\hline
\textbf{Myopie} & Foyer image \textbf{avant} rétine (œil trop long/puissant) & Verge \textbf{divergent} (concave, $V < 0$) & Activité plein air, pauses écrans (20-20-20), éclairage \\\hline
\textbf{Hypermétropie} & Foyer image \textbf{arrière} rétine (œil trop court/faible) & Verge \textbf{convergent} (convexe, $V > 0$) & Correction précoce enfant (éviter strabisme accommodatif) \\\hline
\textbf{Presbytie} & \textbf{Perte accommodation} (cristallin sclérosé, $P_{acc} \downarrow$) & Verge \textbf{convergent pour près} (Addition $> 0$) & Éclairage près, pauses, contrôle régulier après 40 ans \\\hline
\textbf{Cataracte} & \textbf{Opacification cristallin} (diffusion lumière) & \textbf{Chirurgie} (Phaco + Implant IOL) & Protection UV, arrêt tabac, contrôle glycémie, antioxydants \\\hline
\textbf{Glaucome} & \textbf{Neuropathie optique} (souvent PIO $\uparrow$ $\rightarrow$ excavation papille) & \textbf{Hypotonisants} (Collyres 1\textsuperscript{re} ligne : PG analogues) $\pm$ Laser/Chirurgie & \textbf{Dépistage annuel > 40 ans} (PIO, Fond d'œil, Champ visuel) \\\hline
\end{tabularx}
\end{center}

\begin{exercice}[Entraînement BEPC - QCM / Réponse courte]
\begin{enumerate}
    \item Un élève de 3ème voit flou le tableau mais lit bien son cahier. Son acuité visuelle de loin est de 4/10. Quel trouble ? \textbf{Myopie}. Quel verre ? \textbf{Divergent (concave)}.
    \item Une personne de 55 ans doit écarter son journal à 50 cm pour le lire. Sa vision de loin est 10/10. Quel trouble ? \textbf{Presbytie}. Cause ? \textbf{Sclérose du cristallin (perte élasticité)}. Correction ? \textbf{Verge convergent pour près (Addition)}.
    \item Un patient de 65 ans a une PIO de 28 mmHg, une excavation de la papille à 0,8 et une réduction du champ visuel périphérique. Pathologie ? \textbf{Glaucome chronique à angle ouvert}. Traitement de 1ère ligne ? \textbf{Collyres hypotonisants (ex: analogues prostaglandines)}.
    \item Un patient de 70 ans voit « comme derrière une cascade », est ébloui la nuit, les couleurs sont jaunâtres. Pathologie ? \textbf{Cataracte (nucléaire)}. Traitement ? \textbf{Chirurgie (Phacoémulsification + Implant)}.
    \item Calculer la vergence d'un verre correcteur pour un myope dont le point lointain est à 25 cm. $V = 1/(-0,25) = \textbf{-4,00 D}$.
\end{enumerate}
\end{exercice}

\begin{corrige}[Entraînement BEPC]
\begin{enumerate}
    \item Myopie. Verre divergent (concave, $V<0$). L'image se forme en avant de la rétine, le verre divergent la recule sur la rétine.
    \item Presbytie. Perte d'élasticité du cristallin (sclérose) $\rightarrow$ baisse pouvoir accommodation $\rightarrow$ point proche recule. Verre convergent pour la vision de près (addition positive).
    \item Glaucome chronique à angle ouvert. Neuropathie optique liée à l'hypertonie oculaire. Collyres hypotonisants (prostaglandines, bêta-bloquants, etc.) à vie.
    \item Cataracte. Opacification du cristallin diffusant la lumière. Chirurgie : extraction cristallin opaque + pose implant intraoculaire.
    \item $PL = 25 \text{ cm} = 0,25 \text{ m}$. Verre divergent : $V = -1/0,25 = -4,00 \text{ D}$.
\end{enumerate}
\end{corrige}

\begin{savaistu}[Contexte Cameroun \& Histoire des sciences]
\begin{itemize}
    \item \textbf{Ibn al-Haytham (Alhazen), XI\textsuperscript{e} siècle :} Premier à décrire correctement l'anatomie de l'œil et le mécanisme de la vision (la lumière entre dans l'œil, ne sort pas de l'œil). Base de l'optique géométrique moderne.
    \item \textbf{Kepler (1604) :} Modélise l'œil comme une chambre noire (caméra obscura) : image inversée sur la rétine.
    \item \textbf{Au Cameroun :} La cataracte est la 1ère cause de cécité évitable. Le glaucome est la 2ème cause de cécité \textbf{irréversible}. Le Programme National de Lutte contre la Cécité (PNLC) organise des campagnes de chirurgie de la cataracte (ex. « Opération Œil » avec ONG) et de dépistage du glaucome dans les CSI. Le port de lunettes de soleil (cat. 3) est encore insuffisant chez les travailleurs extérieurs (agriculteurs, conducteurs moto, vendeurs marché).
    \item \textbf{Myopie scolaire :} En forte augmentation à Yaoundé/Douala (écrans, travail de près intense, manque d'activités extérieures). Prévention : 2h/jour dehors recommandées par l'OMS.
\end{itemize}
\end{savaistu}

\newpage

\coursec{Méthodes et exercices résolus}

\courssub{Méthode 1 : Identifier les structures de l'œil sur un schéma}

\begin{definition}[L'œil, organe de la vision]
L'\cle{œil} est l'organe récepteur de la lumière. Il transforme les rayons lumineux en \textbf{messages nerveux} transmis au cerveau par le \cle{nerf optique}. Il est constitué de trois enveloppes superposées (sclérotique, choroïde, rétine) et de milieux transparents (cornée, humeur aqueuse, cristallin, humeur vitrée) qui dirigent la lumière vers la rétine.
\end{definition}

\begin{methode}[Légender un schéma de l'œil]
Pour identifier correctement les structures de l'œil en coupe :
\begin{enumerate}
\item Repérer d'abord les \cle{trois tuniques} de l'extérieur vers l'intérieur : la \cle{sclérotique} (blanc de l'œil, protectrice), la \cle{choroïde} (riche en vaisseaux sanguins, nourrit l'œil) et la \cle{rétine} (membrane sensible à la lumière).
\item Identifier la \cle{cornée} (partie transparente de la sclérotique, à l'avant), puis l'\cle{iris} (diaphragme coloré) percé de la \cle{pupille}.
\item Localiser le \cle{cristallin} (lentille transparente derrière l'iris), l'\cle{humeur aqueuse} (liquide devant le cristallin) et l'\cle{humeur vitrée} (gros volume gélatineux derrière le cristallin).
\item Repérer le \cle{nerf optique} qui quitte l'œil à l'arrière et la \cle{tache aveugle} (zone de la rétine sans cellules sensibles, au point d'émergence du nerf optique).
\item Vérifier que chaque légende est reliée à la structure par un trait propre, sans croisement.
\end{enumerate}
\textbf{Réflexe} : raisonner toujours dans le même ordre (dehors $\rightarrow$ dedans, avant $\rightarrow$ arrière) pour ne rien oublier.
\end{methode}

\begin{exoresolu}[Légendage d'une coupe d'œil]
Sur le schéma ci-dessous d'une coupe horizontale de l'œil, on a numéroté six structures : 1 (membrane externe blanche), 2 (milieu transparent bombé à l'avant), 3 (ouverture centrale), 4 (lentille biconvexe), 5 (membrane interne tapissant le fond), 6 (cordon nerveux à l'arrière). Nommer chaque structure et donner son rôle.
\begin{center}
\begin{tikzpicture}[scale=0.9]
\draw[thick, popblue] (0,0) circle (2.2);
\draw[thick, popblue!60] (0,0) circle (2.05);
\draw[thick, poporange] (0,0) circle (1.9);
\draw[thick, popink] (2.2,0) arc[start angle=-35, end angle=35, radius=2.9];
\draw[thick, popgreen] (1.35,-0.75) -- (1.35,-0.28);
\draw[thick, popgreen] (1.35,0.75) -- (1.35,0.28);
\draw[thick, poppurple] (0.95,0) ellipse (0.18 and 0.6);
\draw[very thick, popdark] (-2.2,0.25) -- (-3.1,0.25) (-2.2,-0.25) -- (-3.1,-0.25) (-3.1,0.25) -- (-3.1,-0.25);
\node[popink, font=\small] at (3.6,1.1) {1};
\draw[->, popink] (3.4,1.0) -- (1.9,1.15);
\node[popink, font=\small] at (4.4,0.2) {2};
\draw[->, popink] (4.2,0.2) -- (2.75,0.15);
\node[popink, font=\small] at (2.6,-1.6) {3};
\draw[->, popink] (2.4,-1.5) -- (1.4,-0.15);
\node[popink, font=\small] at (0.9,1.5) {4};
\draw[->, popink] (0.9,1.3) -- (0.95,0.65);
\node[popink, font=\small] at (-0.9,-2.7) {5};
\draw[->, popink] (-0.9,-2.5) -- (-0.9,-1.95);
\node[popink, font=\small] at (-3.6,-1.2) {6};
\draw[->, popink] (-3.4,-1.1) -- (-2.9,-0.15);
\end{tikzpicture}
\end{center}
\tcblower
\textbf{Solution détaillée.}

On procède de l'extérieur vers l'intérieur, et de l'avant vers l'arrière :
\begin{enumerate}
\item \textbf{Structure 1} : membrane externe blanche et résistante $\rightarrow$ c'est la \cle{sclérotique}. Rôle : protection mécanique de l'œil.
\item \textbf{Structure 2} : milieu transparent et bombé à l'avant $\rightarrow$ c'est la \cle{cornée}. Rôle : laisser passer la lumière et participer à sa convergence (c'est une lentille fixe).
\item \textbf{Structure 3} : ouverture centrale de l'iris $\rightarrow$ c'est la \cle{pupille}. Rôle : réguler la quantité de lumière qui pénètre dans l'œil (elle se rétrécit en pleine lumière, se dilate dans l'obscurité).
\item \textbf{Structure 4} : lentille biconvexe derrière l'iris $\rightarrow$ c'est le \cle{cristallin}. Rôle : faire converger les rayons lumineux sur la rétine ; sa courbure peut varier (accommodation).
\item \textbf{Structure 5} : membrane interne tapissant le fond de l'œil $\rightarrow$ c'est la \cle{rétine}. Rôle : recevoir la lumière grâce à ses cellules sensibles (bâtonnets et cônes) et transformer le message lumineux en message nerveux.
\item \textbf{Structure 6} : cordon nerveux à l'arrière $\rightarrow$ c'est le \cle{nerf optique}. Rôle : conduire les messages nerveux de la rétine vers le cerveau.
\end{enumerate}
\textbf{Conclusion} : la lumière traverse successivement la cornée, l'humeur aqueuse, la pupille, le cristallin, l'humeur vitrée, puis forme une image sur la rétine ; le nerf optique transmet l'information au cerveau.
\end{exoresolu}

\courssub{Méthode 2 : Expliquer la formation de l'image et l'accommodation}

\begin{definition}[Accommodation]
L'\cle{accommodation} est la modification de la courbure du \cle{cristallin}, sous l'action des \textbf{muscles ciliaires}, qui permet d'obtenir une image nette sur la rétine quelle que soit la distance de l'objet regardé.
\end{definition}

\begin{methode}[Décrire le trajet de la lumière et l'accommodation]
\begin{enumerate}
\item Citer dans l'ordre les \cle{milieux transparents} traversés par la lumière : cornée $\rightarrow$ humeur aqueuse $\rightarrow$ pupille $\rightarrow$ cristallin $\rightarrow$ humeur vitrée $\rightarrow$ rétine.
\item Préciser que l'image se forme sur la \cle{rétine} et qu'elle est \textbf{renversée} (plus petite et à l'envers) ; c'est le cerveau qui la redresse.
\item Pour l'\cle{accommodation} (vision nette à différentes distances) :
\begin{itemize}
\item objet \textbf{proche} : les muscles ciliaires se contractent, le cristallin \textbf{se bombe} (devient plus convergent) ;
\item objet \textbf{éloigné} : les muscles ciliaires se relâchent, le cristallin \textbf{s'aplatit} (devient moins convergent).
\end{itemize}
\item Conclure : l'accommodation est la modification de la courbure du cristallin qui permet une image nette sur la rétine quelle que soit la distance de l'objet.
\end{enumerate}
\textbf{À retenir} : image nette = image \emph{sur la rétine}. Si l'image se forme devant ou derrière la rétine, la vision est floue.
\end{methode}

\begin{exoresolu}[Observer un livre puis le paysage]
Un élève lit son cahier posé à \SI{30}{cm} de ses yeux, puis lève la tête pour regarder un arbre situé à \SI{50}{m}. Sa vision reste nette dans les deux cas.
\begin{enumerate}
\item Nommer le phénomène qui permet cette vision nette dans les deux situations.
\item Décrire l'état du cristallin et des muscles ciliaires dans chaque cas.
\end{enumerate}
\tcblower
\textbf{Solution détaillée.}
\begin{enumerate}
\item Le phénomène mis en jeu est l'\cle{accommodation} : c'est la modification de la courbure du cristallin qui permet d'obtenir une image nette sur la rétine, que l'objet soit proche ou éloigné.
\item \textbf{Lecture du cahier (objet proche, \SI{30}{cm})} : les rayons lumineux issus d'un point proche arrivent très divergents sur l'œil. Pour les faire converger sur la rétine, les \textbf{muscles ciliaires se contractent} et le \textbf{cristallin se bombe} : il devient plus épais et plus convergent.\\
\textbf{Observation de l'arbre (objet éloigné, \SI{50}{m})} : les rayons arrivent presque parallèles. Les \textbf{muscles ciliaires se relâchent} et le \textbf{cristallin s'aplatit} : il devient moins convergent, ce qui suffit à projeter l'image exactement sur la rétine.
\end{enumerate}
\textbf{Conclusion} : grâce à l'accommodation, le cristallin joue le rôle d'une lentille à courbure variable : bombé pour la vision de près, aplati pour la vision de loin, l'image se forme toujours sur la rétine.
\end{exoresolu}

\courssub{Méthode 3 : Diagnostiquer et corriger une anomalie de la vision}

\begin{methode}[Raisonner sur myopie, hypermétropie et presbytie]
\begin{enumerate}
\item Se poser la question clé : \textbf{où se forme l'image par rapport à la rétine ?}
\item Appliquer le tableau de diagnostic :
\begin{center}
\begin{tabularx}{\linewidth}{|l|X|X|X|}
\hline
\rowcolor{popblueL} \textbf{Anomalie} & \textbf{Cause} & \textbf{Image} & \textbf{Correction} \\
\hline
Myopie & Œil trop long ou cristallin trop convergent & En \textbf{avant} de la rétine & Verres \textbf{divergents} \\
\hline
Hypermétropie & Œil trop court ou cristallin pas assez convergent & En \textbf{arrière} de la rétine & Verres \textbf{convergents} \\
\hline
Presbytie & Cristallin qui perd son élasticité avec l'âge & Accommodation insuffisante de près & Verres \textbf{convergents} (lecture) \\
\hline
\end{tabularx}
\end{center}
\item Déduire la conséquence sur la vision : le myope voit flou \textbf{de loin}, l'hypermétrope voit flou \textbf{de près} (effort permanent d'accommodation), le presbyte ne voit plus net \textbf{de près} à partir d'environ 45 ans.
\item Conclure en reliant la correction optique au déplacement de l'image : les verres ramènent l'image \emph{exactement sur la rétine}.
\end{enumerate}
\end{methode}

\begin{exoresolu}[Le cas de M. Etoundi]
M. Etoundi, commerçant à Yaoundé, 52 ans, se plaint : « Je vois très bien les motos au loin, mais pour lire mes factures, je dois les tenir à bout de bras. » Sa fille de 16 ans, elle, voit parfaitement son cahier mais ne lit plus ce qui est écrit au tableau en classe.
\begin{enumerate}
\item Identifier l'anomalie de vision de chacun en justifiant.
\item Indiquer le type de verres correcteurs adapté à chacun et expliquer leur action.
\end{enumerate}
\tcblower
\textbf{Solution détaillée.}
\begin{enumerate}
\item \textbf{M. Etoundi (52 ans)} : il voit bien de loin mais mal de près, et ce trouble est apparu avec l'âge. Il s'agit de la \cle{presbytie} : avec le vieillissement, le cristallin perd son élasticité et ne peut plus se bomber suffisamment pour accommoder sur les objets proches. Tenir le texte à bout de bras éloigne l'objet et compense partiellement ce défaut.\\
\textbf{Sa fille (16 ans)} : elle voit bien de près mais flou de loin (le tableau). Il s'agit de la \cle{myopie} : son œil est trop long (ou son cristallin trop convergent), donc l'image d'un objet éloigné se forme \textbf{en avant de la rétine}.
\item \textbf{Pour M. Etoundi} : des verres \textbf{convergents} pour la lecture. Ils ajoutent la convergence que le cristallin vieillissant ne peut plus fournir et ramènent l'image des objets proches sur la rétine.\\
\textbf{Pour sa fille} : des verres \textbf{divergents}. Ils font diverger légèrement les rayons avant leur entrée dans l'œil, ce qui recule l'image jusqu'à la rétine : la vision de loin redevient nette.
\end{enumerate}
\textbf{Conclusion} : M. Etoundi est presbyte (verres convergents) et sa fille est myope (verres divergents) ; dans les deux cas, la correction a pour but de replacer l'image exactement sur la rétine.
\end{exoresolu}

\courssub{Méthode 4 : Distinguer une anomalie optique d'une maladie de l'œil}

\begin{methode}[Différencier cataracte et glaucome]
\begin{enumerate}
\item Vérifier d'abord si le trouble est une simple \textbf{anomalie optique} (myopie, hypermétropie, presbytie) : l'œil est sain, seule la mise au point est déficiente, et des verres corrigent totalement le trouble.
\item Si le trouble persiste malgré les verres, penser à une \textbf{maladie} :
\begin{itemize}
\item \cle{cataracte} : \textbf{opacification du cristallin} ; la lumière ne traverse plus bien, la vision devient voilée, comme à travers un brouillard ; fréquente chez la personne âgée ; traitement : \textbf{chirurgie} (remplacement du cristallin opaque par un cristallin artificiel).
\item \cle{glaucome} : \textbf{augmentation de la pression à l'intérieur de l'œil} qui détruit progressivement les fibres du nerf optique ; le champ visuel se rétrécit peu à peu ; peut conduire à la \textbf{cécité} si non traité ; traitement : collyres ou chirurgie pour faire baisser la pression.
\end{itemize}
\item Argumenter à partir du \textbf{siège de l'atteinte} (cristallin pour la cataracte, nerf optique pour le glaucome) et de la \textbf{conséquence} (lumière bloquée / message nerveux détruit).
\end{enumerate}
\end{methode}

\begin{exoresolu}[Deux patients au centre de santé]
Deux personnes âgées consultent au centre de santé de Bafoussam.\\
Patient A : « Tout est devenu flou et voilé depuis deux ans, même avec mes lunettes ; les lumières m'éblouissent la nuit. » À l'examen, son cristallin apparaît opaque.\\
Patient B : « Je vois encore net au centre, mais je ne vois plus ce qui se passe sur les côtés ; j'ai souvent mal à la tête. » Le médecin mesure une pression intraoculaire élevée.
\begin{enumerate}
\item Identifier la maladie de chaque patient en justifiant par deux arguments.
\item Expliquer pourquoi le patient B risque la cécité s'il n'est pas soigné.
\end{enumerate}
\tcblower
\textbf{Solution détaillée.}
\begin{enumerate}
\item \textbf{Patient A} : il souffre de la \cle{cataracte}. Argument 1 : son cristallin est devenu \textbf{opaque} (signe direct observé à l'examen). Argument 2 : sa vision est \textbf{voilée} et les verres ne corrigent rien, car le problème n'est pas un défaut de convergence mais l'impossibilité pour la lumière de traverser le cristallin. Le traitement est chirurgical : on remplace le cristallin opaque par un cristallin artificiel.\\
\textbf{Patient B} : il souffre d'un \cle{glaucome}. Argument 1 : la \textbf{pression à l'intérieur de son œil est élevée} (mesurée par le médecin). Argument 2 : son \textbf{champ visuel se rétrécit} (il ne voit plus sur les côtés), signe de la destruction progressive des fibres du nerf optique, en périphérie d'abord.
\item La pression intraoculaire élevée comprime et détruit peu à peu les fibres du \cle{nerf optique}. Or les cellules nerveuses détruites \textbf{ne se régénèrent pas} : les messages nerveux ne parviennent plus au cerveau. La destruction progresse de la périphérie vers le centre ; si rien n'est fait (collyres ou chirurgie pour abaisser la pression), la perte de vision devient totale et \textbf{irréversible} : c'est la cécité.
\end{enumerate}
\textbf{Conclusion} : le patient A présente une cataracte (cristallin opaque, curable par chirurgie) et le patient B un glaucome (pression élevée détruisant le nerf optique, urgence à traiter pour éviter une cécité définitive).
\end{exoresolu}

\courssub{Méthode 5 : Appliquer les règles d'hygiène de la vision à une situation concrète}

\begin{methode}[Conseiller une bonne hygiène visuelle]
\begin{enumerate}
\item Identifier dans la situation décrite les \textbf{comportements à risque} : lecture dans l'obscurité, écran trop proche ou trop longtemps, soleil regardé directement, frottement des yeux avec les mains sales, automédication.
\item Proposer les règles d'\cle{hygiène de la vision} correspondantes :
\begin{itemize}
\item lire et travailler avec un \textbf{bon éclairage}, sans éblouissement ;
\item respecter une \textbf{distance suffisante} (environ \SI{30}{cm} pour un livre, \SI{50}{cm} à \SI{70}{cm} pour un écran) et faire des \textbf{pauses régulières} ;
\item ne jamais regarder \textbf{directement le soleil} ou une source intense ; porter des lunettes de soleil adaptées ;
\item ne pas se frotter les yeux avec les mains sales ; se laver les mains ;
\item consulter régulièrement un \textbf{ophtalmologiste}, surtout en cas de baisse de vision, et éviter l'automédication (gouttes non prescrites).
\end{itemize}
\item Justifier chaque conseil par la structure de l'œil concernée (fatigue des muscles ciliaires, lésion de la rétine, infection de la conjonctive\ldots).
\end{enumerate}
\end{methode}

\begin{exoresolu}[Les habitudes de Sandrine]
Sandrine, élève de 3ème, prépare son BEPC. Chaque soir, elle révise couchée dans son lit, à la lueur de son téléphone portable tenu à \SI{15}{cm} des yeux, la lampe de la chambre éteinte, souvent jusqu'à 1 h du matin. Depuis quelques semaines, elle a mal à la tête et les yeux qui piquent ; elle se frotte souvent les yeux avec les mains.
\begin{enumerate}
\item Relever trois comportements à risque pour la vision de Sandrine.
\item Donner pour chacun un conseil d'hygiène visuelle justifié.
\end{enumerate}
\tcblower
\textbf{Solution détaillée.}
\begin{enumerate}
\item Comportements à risque relevés :
\begin{itemize}
\item elle révise \textbf{dans l'obscurité} (lampe éteinte), seul l'écran éclaire ;
\item elle tient l'écran \textbf{très près des yeux} (\SI{15}{cm}) et pendant \textbf{plusieurs heures sans pause} ;
\item elle \textbf{se frotte les yeux avec les mains} (probablement sales).
\end{itemize}
\item Conseils justifiés :
\begin{itemize}
\item \textbf{Travailler avec un bon éclairage} : dans l'obscurité, la pupille se dilate et l'effort d'accommodation est maximal ; un éclairage correct réduit la fatigue visuelle et les maux de tête.
\item \textbf{Respecter une distance de \SI{30}{cm} à \SI{50}{cm} et faire des pauses} (regarder au loin quelques minutes toutes les 20 à 30 minutes) : une distance trop courte oblige les muscles ciliaires à rester contractés en permanence pour bomber le cristallin, d'où la fatigue et les maux de tête ; la pause les relâche.
\item \textbf{Ne pas se frotter les yeux avec les mains sales et se laver les mains} : les mains transportent des microorganismes qui peuvent provoquer une infection de l'œil (conjonctivite), ce qui explique les picotements.
\end{itemize}
\end{enumerate}
\textbf{Conclusion} : en adoptant un bon éclairage, une distance de lecture correcte, des pauses régulières et une hygiène des mains, Sandrine protégera ses yeux et révisera dans de meilleures conditions pour son BEPC.
\end{exoresolu}

\courssub{Méthode 6 : Rédiger et justifier une conclusion en démarche scientifique}

\begin{methode}[Construire une réponse argumentée au BEPC]
\begin{enumerate}
\item \textbf{Observer / analyser} : dégager les faits du document ou de l'expérience (ce que l'on voit, les mesures, les symptômes) sans les interpréter tout de suite.
\item \textbf{Formuler une hypothèse} : proposer une explication possible (ex. « l'image se forme en avant de la rétine »).
\item \textbf{Confronter} : vérifier que l'hypothèse rend compte de \emph{tous} les faits observés ; éliminer les hypothèses incompatibles (ex. la presbytie est impossible chez un enfant de 10 ans).
\item \textbf{Conclure} : rédiger une phrase complète qui répond à la question posée, en reliant le fait observé au mécanisme biologique (structure de l'œil concernée $\rightarrow$ conséquence sur l'image $\rightarrow$ trouble de la vision).
\item \textbf{Vérifier la rédaction} : vocabulaire scientifique exact (rétine, cristallin, accommodation\ldots), orthographe soignée, unités si des valeurs sont citées.
\end{enumerate}
\textbf{Réflexe} : une conclusion sans justification (sans le « parce que ») ne rapporte que la moitié des points.
\end{methode}

\begin{exoresolu}[L'expérience de l'œil réduit]
Au laboratoire, on réalise un « œil réduit » : une lentille convergente (le cristallin) projette l'image d'une bougie sur un écran translucide (la rétine) placé à \SI{20}{cm} derrière la lentille. L'image de la bougie est nette sur l'écran. On éloigne alors l'écran de \SI{3}{cm} vers l'arrière sans rien changer d'autre : l'image devient floue.
\begin{enumerate}
\item À quelle anomalie de la vision cette deuxième situation correspond-elle ? Justifier.
\item Proposer, à l'aide du montage, un moyen de rendre l'image à nouveau nette, et en déduire le principe de correction de cette anomalie.
\end{enumerate}
\tcblower
\textbf{Solution détaillée.}
\begin{enumerate}
\item \textbf{Analyse des faits} : la lentille et la bougie n'ont pas bougé ; seul l'écran (la rétine) a été reculé. L'image, qui se forme toujours à la même place (à \SI{20}{cm} de la lentille), se trouve désormais \textbf{en avant de l'écran} : sur l'écran, elle est floue.\\
\textbf{Hypothèse} : cela correspond à un œil trop long, dont la rétine est trop éloignée du cristallin.\\
\textbf{Vérification} : dans la \cle{myopie}, l'image d'un objet éloigné se forme en avant de la rétine parce que l'œil est trop long (ou le cristallin trop convergent) : tous les faits sont expliqués. L'hypermétropie (image derrière la rétine) est incompatible avec l'observation.\\
\textbf{Conclusion} : cette situation modélise la \textbf{myopie}, car l'image se forme en avant de la rétine.
\item \textbf{Expérience} : on place devant la lentille une \textbf{lentille divergente} bien choisie : elle fait légèrement diverger les rayons lumineux, l'image recule et redevient nette sur l'écran reculé.\\
\textbf{Conclusion} : le principe de correction de la myopie est le port de \textbf{verres divergents}, qui ramènent l'image exactement sur la rétine.
\end{enumerate}
\textbf{Conclusion générale} : l'œil réduit montre qu'un œil trop long est myope (image en avant de la rétine) et que des verres divergents corrigent ce défaut en replaçant l'image sur la rétine.
\end{exoresolu}

\begin{aretenir}[L'essentiel de la leçon]
\begin{itemize}
\item La lumière traverse la \textbf{cornée}, l'\textbf{humeur aqueuse}, la \textbf{pupille}, le \textbf{cristallin}, l'\textbf{humeur vitrée}, puis forme une image renversée sur la \textbf{rétine} ; le \textbf{nerf optique} transmet le message nerveux au cerveau.
\item L'\textbf{accommodation} est la modification de la courbure du cristallin : il se \textbf{bombe} pour voir de près, s'\textbf{aplatit} pour voir de loin.
\item \textbf{Myopie} : image en avant de la rétine, vision floue de loin, verres divergents. \textbf{Hypermétropie} : image en arrière de la rétine, vision floue de près, verres convergents. \textbf{Presbytie} : perte d'élasticité du cristallin avec l'âge, verres convergents pour lire.
\item \textbf{Cataracte} : cristallin opaque, vision voilée, traitement chirurgical. \textbf{Glaucome} : pression intraoculaire élevée détruisant le nerf optique, risque de cécité irréversible.
\item \textbf{Hygiène de la vision} : bon éclairage, distance de lecture suffisante, pauses, ne pas regarder le soleil, mains propres, consultation régulière chez l'ophtalmologiste.
\end{itemize}
\end{aretenir}

\begin{attention}[Les 5 erreurs qui coûtent des points au BEPC]
\begin{enumerate}
\item \textbf{Confondre pupille et iris} : l'iris est le diaphragme \emph{coloré}, la pupille est l'\emph{ouverture} centrale qui laisse passer la lumière. Écrire « la lumière traverse l'iris » est faux.
\item \textbf{Inverser myopie et hypermétropie} : le \textbf{myope} voit flou de \emph{loin} (image en avant de la rétine, verres \emph{divergents}) ; l'\textbf{hypermétrope} voit flou de \emph{près} (image en arrière de la rétine, verres \emph{convergents}). Une inversion dans le tableau de diagnostic fait perdre tous les points de la question.
\item \textbf{Dire que l'image se forme sur le cristallin} : le cristallin est une \emph{lentille} qui dirige la lumière ; l'image se forme sur la \textbf{rétine}, et c'est le \textbf{cerveau} qui interprète le message (et redresse l'image renversée).
\item \textbf{Confondre cataracte et glaucome} : la cataracte est l'opacification du \emph{cristallin} (vision voilée, guérissable par chirurgie) ; le glaucome est une \emph{pression intraoculaire élevée} qui détruit le \emph{nerf optique} (cécité irréversible si non traitée). Ne pas écrire que le glaucome se soigne avec des lunettes.
\item \textbf{Conclure sans justifier} : affirmer « c'est une myopie » sans expliquer (œil trop long $\rightarrow$ image en avant de la rétine $\rightarrow$ vision floue de loin) ne vaut que la moitié des points. Toute conclusion doit contenir le mécanisme et le vocabulaire scientifique exact.
\end{enumerate}
\end{attention}

\newpage

\coursec{Exercices}

\courssub{Je vérifie mes connaissances}

\begin{exercice}[QCM : Anatomie et formation de l'image]
Parmi les affirmations suivantes, une seule est exacte. La recopier en justifiant brièvement votre choix.
\begin{enumerate}
\item La cornée est la lentille principale responsable de l'accommodation.
\item L'iris régule la quantité de lumière entrant dans l'œil en modifiant le diamètre de la pupille.
\item Le cristallin est un muscle qui se contracte pour faire la mise au point.
\item L'image se forme sur la choroïde, qui la transmet au nerf optique.
\end{enumerate}
\end{exercice}

\begin{exercice}[Vrai ou Faux justifié : Anomalies de la vision]
Dire si les affirmations suivantes sont vraies ou fausses. Justifier chaque réponse par une phrase scientifique précise.
\begin{enumerate}
\item Le myope voit net de loin mais flou de près.
\item L'hypermétropie est due à un œil trop court ou pas assez convergent ; l'image d'un objet lointain se forme en arrière de la rétine.
\item La presbytie est une anomalie de la réfraction corrigée par des verres divergents.
\item Le glaucome se caractérise par une opacification du cristallin entraînant une baisse de l'acuité visuelle.
\end{enumerate}
\end{exercice}

\begin{exercice}[Définitions et vocabulaire]
Donner la définition précise des termes suivants en utilisant le vocabulaire scientifique du cours :
\begin{enumerate}
\item Accommodation
\item Point proche (punctum proximum)
\item Cataracte
\item Pression intra-oculaire
\item Hygiène visuelle
\end{enumerate}
\end{exercice}

\begin{exercice}[Schéma à compléter : Le trajet de la lumière]
Recopier et compléter la chaîne suivante, qui décrit le trajet de la lumière et la formation du message nerveux, en plaçant dans l'ordre les structures de l'œil :
\begin{center}
Objet lumineux $\rightarrow$ \ldots\ldots $\rightarrow$ pupille $\rightarrow$ \ldots\ldots $\rightarrow$ corps vitré $\rightarrow$ \ldots\ldots $\rightarrow$ photorécepteurs $\rightarrow$ \ldots\ldots $\rightarrow$ cerveau
\end{center}
Préciser sur quelle structure se forme l'image réelle et inversée de l'objet.
\end{exercice}

\courssub{Je m'entraîne}

\begin{exercice}[Analyse de schéma : Myopie et correction]
On observe le schéma simplifié d'un œil myope regardant un objet lointain (à l'infini). Les rayons lumineux parallèles convergent en un point $F'$ situé \textbf{en avant de la rétine}.
\begin{enumerate}
\item Sur une copie du schéma, tracer le chemin des rayons après passage dans un verre correcteur placé devant l'œil, de manière à ce que l'image se forme sur la rétine.
\item Quel est le type de verre utilisé (convergent ou divergent) ? Justifier par l'effet de ce verre sur les rayons incidents.
\item Expliquer, avec vos mots, pourquoi ce verre permet au myope de voir net de loin.
\end{enumerate}
\end{exercice}

\begin{exercice}[Mécanisme de l'accommodation]
On compare deux situations pour un œil emmétrope (normal) : observation d'un objet lointain (A) et observation d'un objet proche à \SI{25}{cm} (B).
\begin{enumerate}
\item Recopier et compléter le tableau comparatif ci-dessous :
\begin{center}
\begin{tabularx}{\linewidth}{|l|X|X|}\hline
\textbf{Paramètre} & \textbf{Situation A : Objet lointain} & \textbf{Situation B : Objet proche (25 cm)} \\ \hline
État des muscles ciliaires & & \\ \hline
Tension du ligament suspenseur (zonule de Zinn) & & \\ \hline
Forme du cristallin (bombé / aplati) & & \\ \hline
Puissance de convergence du cristallin & & \\ \hline
\end{tabularx}
\end{center}
\item Expliquer pourquoi l'accommodation a une limite (point proche) chez le jeune adulte.
\end{enumerate}
\end{exercice}

\begin{exercice}[Diagnostic différentiel : Cataracte ou Glaucome]
Deux patients consultent à l'hôpital de district de Bafoussam.
\begin{itemize}
\item \textbf{Patient 1 (M. K., 68 ans)} : « Docteur, je vois comme à travers une chute d'eau ou un brouillard. Les couleurs sont ternes. Ça empire depuis deux ans. » L'examen à la lampe à fente montre un cristallin blanchâtre, opaque.
\item \textbf{Patient 2 (Mme N., 55 ans)} : « Je ne ressens pas de douleur, mais je me cogne souvent dans les meubles sur les côtés. Je ne vois plus bien les côtés. » La mesure de la pression intra-oculaire (PIO) donne \SI{28}{mmHg} (normale : \SI{10}{à 21}{mmHg}). Le fond d'œil montre une atteinte du nerf optique.
\end{itemize}
\begin{enumerate}
\item Pour chaque patient, poser le diagnostic (cataracte ou glaucome) et justifier par deux arguments cliniques tirés du texte.
\item Préciser le traitement curatif principal pour le patient 1.
\item Expliquer pourquoi le glaucome (patient 2) est une urgence thérapeutique même s'il est indolore.
\end{enumerate}
\end{exercice}

\begin{exercice}[Hygiène visuelle en milieu scolaire]
Vous êtes délégué de la classe de 3ème au Lycée de Nkolbisson. La salle informatique vient d'être équipée de 30 nouveaux ordinateurs. Le professeur principal vous demande de proposer \textbf{5 règles d'hygiène visuelle} à afficher près de chaque poste de travail pour prévenir la fatigue oculaire (asthénopie) et la myopie scolaire.
\begin{enumerate}
\item Rédiger ces 5 règles sous forme d'affiches courtes et percutantes.
\item Pour chaque règle, justifier scientifiquement son efficacité (ex : lien avec l'accommodation, la convergence, le clignement, l'éclairage, la distance de lecture).
\end{enumerate}
\end{exercice}

\courssub{Je me prépare au BEPC}

\begin{exercice}[Sujet type BEPC : Étude intégrée de la vision — 12 points]
\textbf{Contexte :} Au Cameroun, les erreurs de réfraction non corrigées sont la première cause de déficience visuelle chez l'enfant scolarisé. Un dépistage est organisé dans un collège de Yaoundé.

\textbf{Document 1 : Résultats de dépistage pour trois élèves}
\begin{center}
\begin{tabularx}{\linewidth}{|l|X|X|X|}\hline
\textbf{Élève} & \textbf{A (14 ans)} & \textbf{B (14 ans)} & \textbf{C (14 ans)} \\ \hline
Vision de loin (sans correction) & Floue & Nette & Nette \\ \hline
Vision de près (à 30 cm, sans correction) & Nette & Floue (fatigue, maux de tête) & Floue (bras trop courts) \\ \hline
Point lointain (punctum remotum) & \SI{40}{cm} & À l'infini & À l'infini \\ \hline
Point proche (punctum proximum) & \SI{10}{cm} & \SI{15}{cm} & \SI{60}{cm} \\ \hline
\end{tabularx}
\end{center}

\textbf{Document 2 : Schémas optiques simplifiés}
\begin{itemize}
\item \textbf{Schéma 1 :} Œil regardant l'infini. Rayons parallèles $\rightarrow$ convergence en $F'$ \textbf{avant} la rétine.
\item \textbf{Schéma 2 :} Œil regardant l'infini. Rayons parallèles $\rightarrow$ convergence en $F'$ \textbf{sur} la rétine. Pour un objet à 30 cm, convergence en $F''$ \textbf{après} la rétine.
\item \textbf{Schéma 3 :} Œil regardant l'infini. Rayons parallèles $\rightarrow$ convergence en $F'$ \textbf{sur} la rétine. Pour un objet à 30 cm, convergence en $F''$ \textbf{après} la rétine malgré l'accommodation, car le cristallin est peu déformable.
\end{itemize}

\textbf{Questions :}
\begin{enumerate}
\item \textbf{Diagnostic (3 pts) :} À l'aide du Document 1, identifier l'anomalie visuelle de chaque élève (A, B, C). Justifier chaque diagnostic par un élément du tableau.
\item \textbf{Optique (3 pts) :} Associer chaque élève (A, B, C) au schéma correspondant (1, 2 ou 3) du Document 2. Justifier l'association pour l'élève A en décrivant le trajet des rayons.
\item \textbf{Correction optique (3 pts) :}
\begin{enumerate}
\item Préciser la nature des verres correcteurs (convergents ou divergents) pour l'élève A et l'élève B. Justifier.
\item Expliquer, sans calcul, comment le verre de l'élève A modifie le trajet des rayons pour que l'image se forme sur la rétine.
\item Pourquoi l'élève C ne peut-il pas être corrigé par un verre unique pour la vision de loin et de près ? Proposer une solution pratique.
\end{enumerate}
\item \textbf{Prévention (3 pts) :} L'infirmière scolaire remarque que l'élève A passe 4 heures par jour sur son téléphone portable à \SI{20}{cm} des yeux, dans le noir. Rédiger un message de sensibilisation de 5 lignes maximum adressé aux élèves, expliquant \textbf{deux mécanismes physiologiques} par lesquels cette habitude aggrave la fatigue visuelle et favorise la myopie.
\end{enumerate}
\end{exercice}

\begin{exercice}[Sujet type BEPC : Glaucome et pression intra-oculaire — 10 points]
\textbf{Contexte :} M. Essomba, 52 ans, chauffeur de poids lourd sur l'axe Douala--Bertoua, consulte pour « vision en tunnel » et difficultés de conduite de nuit. Il n'a jamais consulté d'ophtalmologiste.

\textbf{Document 1 : Circulation de l'humeur aqueuse (description)}
L'humeur aqueuse est un liquide transparent sécrété en permanence à l'intérieur de l'œil (1). Elle circule dans la partie antérieure du globe, passe par la pupille (2), puis est normalement évacuée par un système de drainage (3) vers les vaisseaux sanguins (4).

\textbf{Document 2 : Résultats cliniques}
\begin{itemize}
\item Acuité visuelle : œil droit : 8/10 ; œil gauche : 6/10.
\item Champ visuel : rétrécissement concentrique important de l'œil gauche (vision tubulaire), début de rétrécissement à l'œil droit.
\item Pression intra-oculaire (PIO) : œil droit : \SI{24}{mmHg} ; œil gauche : \SI{32}{mmHg}. (Norme : \SI{10}{à 21}{mmHg}.)
\item Fond d'œil : atteinte du nerf optique, plus marquée à gauche.
\end{itemize}

\textbf{Questions :}
\begin{enumerate}
\item \textbf{Physiologie (2 pts) :} À l'aide du Document 1, expliquer comment s'établit l'équilibre de la pression intra-oculaire dans un œil sain. Que se passe-t-il si le drainage (3) est ralenti ?
\item \textbf{Diagnostic (2 pts) :} Poser le diagnostic de M. Essomba. Justifier par \textbf{trois} arguments cliniques issus du Document 2.
\item \textbf{Mécanisme lésionnel (2 pts) :} Expliquer le lien entre l'élévation prolongée de la PIO et l'atteinte du nerf optique observée au fond d'œil. Quelle est la conséquence sur le champ visuel ?
\item \textbf{Prise en charge et santé publique (4 pts) :}
\begin{enumerate}
\item Citer le principe du traitement du glaucome (but recherché et moyens : médicaments en collyre ou chirurgie).
\item Pourquoi le dépistage systématique de la PIO après 40 ans est-il recommandé au Cameroun, notamment chez les conducteurs professionnels ? Argumenter à partir du cas de M. Essomba.
\item Rédiger un message court (3 lignes) pour une campagne d'affichage dans les gares routières : « Glaucome : le voleur silencieux de la vue. »
\end{enumerate}
\end{enumerate}
\end{exercice}

\begin{exercice}[Sujet type BEPC : Synthèse — L'œil, de la lumière à l'image perçue — 8 points]
\textbf{Consigne :} Répondre aux questions suivantes en mobilisant vos connaissances sur l'anatomie, la physiologie et la pathologie de l'œil. La clarté du raisonnement et la rigueur du vocabulaire seront évaluées.

\begin{enumerate}
\item \textbf{Trajet de la lumière (2,5 pts) :} En partant d'un objet lumineux situé devant l'œil, décrire le trajet du rayon lumineux à travers les milieux transparents de l'œil jusqu'à sa transformation en message nerveux. Citer dans l'ordre : cornée, humeur aqueuse, pupille, cristallin, corps vitré, rétine, photorécepteurs, nerf optique. Préciser sur quelle structure se forme l'image réelle et inversée.
\item \textbf{Accommodation et presbytie (2,5 pts) :}
\begin{enumerate}
\item Rappeler le rôle des muscles ciliaires et du ligament suspenseur dans la variation de la courbure du cristallin lors de la vision de près.
\item Expliquer pourquoi, vers 45--50 ans, l'accommodation devient insuffisante pour la lecture à \SI{30}{cm} (presbytie), en lien avec les propriétés du cristallin.
\item Quel verre correcteur prescrit-on ? Pourquoi ce verre ne convient-il pas pour la vision de loin simultanée ?
\end{enumerate}
\item \textbf{Pathologies comparées (3 pts) :} Recopier et compléter le tableau récapitulatif suivant :
\begin{center}
\begin{tabularx}{\linewidth}{|l|X|X|X|}\hline
\textbf{Pathologie} & \textbf{Structure atteinte / Mécanisme principal} & \textbf{Symptôme fonctionnel majeur} & \textbf{Traitement principal} \\ \hline
Cataracte & & & \\ \hline
Glaucome chronique & & & \\ \hline
Myopie & & & \\ \hline
\end{tabularx}
\end{center}
\end{enumerate}
\end{exercice}



\newpage

\coursec{Corrigés des exercices}

\begin{corrige}[Exercice 1 — QCM : Anatomie et formation de l'image]
La réponse exacte est la \textbf{2}.
\begin{itemize}
\item \textbf{1. Fausse} : la cornée est bien un milieu réfringent (elle fait converger les rayons), mais sa courbure est fixe ; c'est le \cle{cristallin} qui, en se déformant, réalise l'accommodation.
\item \textbf{2. Vraie} : l'\cle{iris} est un diaphragme pigmenté ; ses muscles modifient le diamètre de la \cle{pupille} (réflexe pupillaire) : la pupille se rétrécit en lumière vive et se dilate dans l'obscurité.
\item \textbf{3. Fausse} : le cristallin n'est pas un muscle ; c'est une lentille transparente et élastique, déformée par l'action des \cle{muscles ciliaires} et du ligament suspenseur.
\item \textbf{4. Fausse} : l'image se forme sur la \cle{rétine} (et non la choroïde, qui est une couche nutritive et pigmentée) ; ce sont les cellules photoréceptrices de la rétine qui transforment la lumière en message nerveux transmis par le nerf optique.
\end{itemize}
\end{corrige}

\begin{attention}[Point de vigilance — Exercice 1]
Ne confonds jamais \emph{cristallin} (lentille) et \emph{muscles ciliaires} (muscles qui déforment la lentille). C'est la confusion la plus fréquente au BEPC.
\end{attention}

\begin{corrige}[Exercice 2 — Vrai ou Faux justifié : Anomalies de la vision]
\begin{enumerate}
\item \textbf{Faux} : c'est l'inverse. Le myope a un œil trop long (ou trop convergent) : l'image d'un objet lointain se forme \emph{en avant} de la rétine ; il voit donc flou de loin et net de près.
\item \textbf{Vrai} : dans l'hypermétropie, le globe oculaire est trop court ou le système optique pas assez convergent ; l'image d'un objet lointain tend à se former \emph{en arrière} de la rétine. Elle se corrige par des verres \emph{convergents}.
\item \textbf{Faux} : la presbytie est due au vieillissement du cristallin, qui perd son élasticité et ne peut plus se bomber suffisamment pour la vision de près. Elle se corrige par des verres \emph{convergents} (et non divergents) pour la lecture.
\item \textbf{Faux} : l'opacification du cristallin définit la \cle{cataracte}. Le \cle{glaucome} est dû à une élévation de la pression intra-oculaire qui détruit progressivement les fibres du nerf optique, entraînant un rétrécissement du champ visuel.
\end{enumerate}
\end{corrige}

\begin{attention}[Point de vigilance — Exercice 2]
Une justification « par cœur » sans mécanisme (position de l'image, structure atteinte) ne vaut pas tous les points : cite toujours \emph{où} se forme l'image ou \emph{quelle} structure est lésée.
\end{attention}

\begin{corrige}[Exercice 3 — Définitions et vocabulaire]
\begin{enumerate}
\item \textbf{Accommodation} : modification de la courbure du cristallin (il se bombe) sous l'action des muscles ciliaires, permettant la mise au point pour voir net un objet proche.
\item \textbf{Point proche (punctum proximum)} : point le plus rapproché de l'œil vu nettement avec une accommodation maximale (environ \SI{10}{cm} chez le jeune adulte emmétrope).
\item \textbf{Cataracte} : opacification progressive du cristallin, qui devient trouble et empêche la lumière d'atteindre la rétine ; la vision devient floue, comme à travers un brouillard.
\item \textbf{Pression intra-oculaire (PIO)} : pression régnant à l'intérieur du globe oculaire, maintenue par l'équilibre entre la sécrétion et l'évacuation de l'humeur aqueuse (normale : \SI{10}{mmHg} à \SI{21}{mmHg}).
\item \textbf{Hygiène visuelle} : ensemble des règles de comportement (éclairage suffisant, distance de lecture, pauses, dépistage) visant à préserver la vision et à prévenir la fatigue oculaire et les anomalies de la vision.
\end{enumerate}
\end{corrige}

\begin{attention}[Point de vigilance — Exercice 3]
Respecte le vocabulaire officiel : « \emph{point proche} » et non « distance de lecture » ; « \emph{opacification du cristallin} » et non « du cristallin qui devient blanc » ; « \emph{pression intra-oculaire} » et non « tension de l'œil » (terme courant mais imprécis).
\end{attention}

\begin{corrige}[Exercice 4 — Schéma à compléter : Le trajet de la lumière]
\begin{center}
Objet lumineux $\rightarrow$ \textbf{cornée} $\rightarrow$ \textbf{humeur aqueuse} $\rightarrow$ pupille $\rightarrow$ \textbf{cristallin} $\rightarrow$ corps vitré $\rightarrow$ \textbf{rétine} $\rightarrow$ photorécepteurs $\rightarrow$ \textbf{nerf optique} $\rightarrow$ cerveau
\end{center}
L'image réelle et inversée se forme sur la \cle{rétine}. Les photorécepteurs (cônes et bâtonnets) transforment l'énergie lumineuse en message nerveux, acheminé par le nerf optique jusqu'au cerveau, qui redresse et interprète l'image.
\end{corrige}

\begin{attention}[Point de vigilance — Exercice 4]
L'ordre est impératif : la lumière rencontre la cornée \emph{avant} l'humeur aqueuse, puis la pupille, et le cristallin \emph{avant} le corps vitré. Omettre l'humeur aqueuse ou inverser deux structures rend la chaîne fausse.
\end{attention}

\begin{corrige}[Exercice 5 — Analyse de schéma : Myopie et correction]
\begin{enumerate}
\item Le verre correcteur écarte légèrement les rayons parallèles avant leur entrée dans l'œil : ils convergent alors plus loin, exactement sur la rétine. (Sur le schéma, les rayons ressortent du verre légèrement divergents, puis l'œil les fait converger en $F'$ situé \emph{sur} la rétine.)
\item Il s'agit d'un verre \textbf{divergent} : il fait diverger les rayons incidents, ce qui compense la convergence excessive de l'œil myope (œil trop long ou trop convergent).
\item L'œil myope converge trop : l'image se forme en avant de la rétine. Le verre divergent « retarde » la convergence des rayons, de sorte que le point de convergence $F'$ recule et vient se placer exactement sur la rétine : la vision de loin devient nette.
\end{enumerate}
\end{corrige}

\begin{attention}[Point de vigilance — Exercice 5]
Moyen mnémotechnique : \textbf{M}yopie = verres \textbf{M}oins (divergents) ; \textbf{H}ypermétropie et presbytie = verres \textbf{P}lus (convergents). Ne jamais écrire « le verre allonge l'œil » : il agit sur les \emph{rayons}, pas sur le globe oculaire.
\end{attention}

\begin{corrige}[Exercice 6 — Mécanisme de l'accommodation]
\begin{enumerate}
\item Tableau complété :
\begin{center}
\begin{tabularx}{\linewidth}{|l|X|X|}\hline
\rowcolor{popblueL}
\textbf{Paramètre} & \textbf{Situation A : Objet lointain} & \textbf{Situation B : Objet proche (\SI{25}{cm})} \\ \hline
État des muscles ciliaires & Relâchés & Contractés \\ \hline
Tension du ligament suspenseur (zonule de Zinn) & Forte (ligament tendu) & Faible (ligament relâché) \\ \hline
Forme du cristallin & Aplati & Bombé \\ \hline
Puissance de convergence du cristallin & Faible & Forte \\ \hline
\end{tabularx}
\end{center}
\item L'accommodation repose sur l'élasticité du cristallin et sur l'amplitude de contraction des muscles ciliaires. Le cristallin ne peut pas se bomber indéfiniment : lorsque sa courbure maximale est atteinte, l'image d'un objet plus proche se forme en arrière de la rétine et la vision devient floue. Cette distance minimale de vision nette est le \cle{point proche} (environ \SI{10}{cm} chez le jeune adulte).
\end{enumerate}
\end{corrige}

\begin{attention}[Point de vigilance — Exercice 6]
Le raisonnement en chaîne est exigé : muscles ciliaires contractés $\Rightarrow$ ligament suspenseur relâché $\Rightarrow$ cristallin bombé $\Rightarrow$ convergence augmentée. Oublier un maillon fait perdre la logique du mécanisme.
\end{attention}

\begin{corrige}[Exercice 7 — Diagnostic différentiel : Cataracte ou Glaucome]
\begin{enumerate}
\item \textbf{Patient 1 : cataracte.} Arguments : (a) vision floue « comme à travers un brouillard », couleurs ternes, d'évolution progressive ; (b) l'examen montre un cristallin blanchâtre et opaque, c'est-à-dire opacifié. \textbf{Patient 2 : glaucome.} Arguments : (a) rétrécissement du champ visuel périphérique (« je ne vois plus bien les côtés ») ; (b) PIO élevée à \SI{28}{mmHg} (norme : \SI{10}{mmHg} à \SI{21}{mmHg}) avec atteinte du nerf optique au fond d'œil.
\item Le traitement curatif de la cataracte est \textbf{chirurgical} : ablation du cristallin opacifié et son remplacement par un implant (cristallin artificiel).
\item Le glaucome est indolore et évolue silencieusement : la pression intra-oculaire élevée détruit progressivement et \emph{irrémédiablement} les fibres du nerf optique. Sans traitement précoce (collyres hypotenseurs, parfois chirurgie), le rétrécissement du champ visuel aboutit à la cécité définitive. Il faut donc traiter avant l'apparition de symptômes graves.
\end{enumerate}
\end{corrige}

\begin{attention}[Point de vigilance — Exercice 7]
Les arguments doivent être \emph{tirés du texte} (citation de la valeur de la PIO, du symptôme rapporté) : un argument général non appuyé sur le document est insuffisant.
\end{attention}

\begin{corrige}[Exercice 8 — Hygiène visuelle en milieu scolaire]
Exemple de réponse attendue (toute formulation équivalente est acceptée) :
\begin{enumerate}
\item \textbf{« Écran à distance de bras : au moins \SI{50}{cm} ! »} — Une distance suffisante limite l'effort d'accommodation (bombement du cristallin) et de convergence des yeux, qui fatiguent les muscles ciliaires.
\item \textbf{« Toutes les 20 minutes, regarde au loin 20 secondes ! »} — Regarder au loin relâche les muscles ciliaires et le cristallin s'aplatit : cela repose le système d'accommodation.
\item \textbf{« Travaille dans une salle bien éclairée, sans reflets sur l'écran ! »} — Un bon éclairage évite les contrastes excessifs et les efforts permanents d'adaptation de la pupille et de la rétine.
\item \textbf{« Pense à cligner des yeux ! »} — Devant un écran, le clignement diminue ; le film lacrymal ne se renouvelle plus, l'œil devient sec et irrité. Cligner humidifie la cornée.
\item \textbf{« Fais contrôler ta vue régulièrement ! »} — Un dépistage précoce permet de corriger à temps une myopie ou une hypermétropie naissante et d'éviter maux de tête et fatigue visuelle.
\end{enumerate}
\end{corrige}

\begin{attention}[Point de vigilance — Exercice 8]
Chaque règle doit être accompagnée de sa \emph{justification scientifique} (accommodation, convergence, clignement, éclairage, dépistage) : une liste de conseils sans explication ne répond pas à la question.
\end{attention}

\begin{corrige}[Exercice 9 — Sujet type BEPC : Étude intégrée de la vision (12 points)]
\begin{enumerate}
\item \textbf{Élève A : myopie} — vision de loin floue, vision de près nette, point lointain très proche (\SI{40}{cm} au lieu de l'infini). \textbf{Élève B : hypermétropie} — vision de loin nette mais vision de près floue avec fatigue et maux de tête (effort permanent d'accommodation). \textbf{Élève C : presbytie} — vision de loin nette (point lointain à l'infini) mais point proche très reculé (\SI{60}{cm}) : le cristallin ne s'accommode plus assez pour la vision de près. \hfill \emph{[3 pts : 1 pt par diagnostic justifié]}
\item \textbf{A $\rightarrow$ Schéma 1} ; \textbf{B $\rightarrow$ Schéma 2} ; \textbf{C $\rightarrow$ Schéma 3}. Pour l'élève A (myope) : les rayons parallèles venant d'un objet lointain convergent trop tôt, en un point $F'$ situé \emph{en avant} de la rétine, car l'œil est trop long (ou trop convergent) ; l'image sur la rétine est donc floue. \hfill \emph{[3 pts : 1 pt par association + justification du trajet pour A]}
\item \begin{enumerate}
\item Élève A : verres \textbf{divergents} (pour compenser la convergence excessive). Élève B : verres \textbf{convergents} (pour aider l'image à se former sur la rétine et soulager l'accommodation).
\item Le verre divergent écarte les rayons parallèles avant leur entrée dans l'œil ; la convergence se fait alors plus loin, et le point image $F'$ recule exactement jusqu'à la rétine : la vision de loin devient nette.
\item L'élève C (presbyte) voit net de loin sans correction, mais a besoin d'une aide à la convergence pour lire. Un verre convergent unique le gênerait pour la vision de loin. Solution : des verres \textbf{uniquement pour la lecture} (ou des verres à double foyer : partie supérieure sans correction pour le loin, partie inférieure convergente pour le près).
\end{enumerate}
\hfill \emph{[3 pts : 1 pt par sous-question]}
\item Exemple de message : « Élèves, protégez vos yeux ! Regarder un écran à \SI{20}{cm} pendant des heures impose à votre cristallin un effort d'accommodation intense et prolongé : les muscles ciliaires se fatiguent et l'œil peut s'allonger, ce qui aggrave la myopie. Dans le noir, le fort contraste entre l'écran lumineux et l'obscurité fatigue la rétine et la pupille. Règle d'or : écran à plus de \SI{30}{cm}, lumière allumée, et une pause « regard au loin » toutes les 20 minutes ! » \hfill \emph{[3 pts : 2 mécanismes physiologiques corrects (1 pt chacun) + clarté et respect de la contrainte de longueur (1 pt)]}
\end{enumerate}
\end{corrige}

\begin{attention}[Points de vigilance — Barème Exercice 9]
\begin{itemize}
\item Chaque diagnostic doit être appuyé par \emph{un élément chiffré du tableau} (point lointain à \SI{40}{cm}, point proche à \SI{60}{cm}...).
\item Pour la question 3b, le raisonnement doit suivre l'ordre : verre divergent $\rightarrow$ rayons écartés $\rightarrow$ convergence retardée $\rightarrow$ image sur la rétine.
\item Pour la question 4, les \emph{deux mécanismes physiologiques} sont obligatoires : effort d'accommodation prolongé (muscles ciliaires, cristallin) ET fatigue liée au contraste écran/obscurité (pupille, rétine).
\end{itemize}
\end{attention}

\begin{corrige}[Exercice 10 — Sujet type BEPC : Glaucome et pression intra-oculaire (10 points)]
\begin{enumerate}
\item Dans l'œil sain, la PIO est stable car il y a \textbf{équilibre} entre la sécrétion permanente d'humeur aqueuse (1) et son évacuation par le système de drainage (3 puis 4). Si le drainage est ralenti ou bloqué, l'humeur aqueuse s'accumule : la pression intra-oculaire \textbf{augmente}. \hfill \emph{[2 pts : équilibre sécrétion/évacuation (1 pt) + conséquence du ralentissement (1 pt)]}
\item \textbf{Diagnostic : glaucome chronique.} Arguments : (a) PIO élevée dans les deux yeux (\SI{24}{mmHg} et \SI{32}{mmHg}, norme \SI{10}{mmHg} à \SI{21}{mmHg}) ; (b) rétrécissement du champ visuel (vision tubulaire) ; (c) atteinte du nerf optique visible au fond d'œil. \hfill \emph{[2 pts : diagnostic (0,5 pt) + 3 arguments cliniques (0,5 pt chacun)]}
\item La pression élevée comprime de façon prolongée les fibres du nerf optique à l'arrière de l'œil ; privées d'une vascularisation normale, ces fibres meurent progressivement et de façon irréversible. La perte des fibres périphériques se traduit par un \textbf{rétrécissement concentrique du champ visuel} (vision en tunnel), pouvant aboutir à la cécité. \hfill \emph{[2 pts : mécanisme de compression/destruction (1 pt) + conséquence sur le champ visuel (1 pt)]}
\item \begin{enumerate}
\item Le traitement vise à \textbf{faire baisser la pression intra-oculaire} pour stopper la destruction du nerf optique : collyres hypotenseurs (qui diminuent la sécrétion d'humeur aqueuse ou facilitent son évacuation) et, si nécessaire, chirurgie ou laser pour améliorer le drainage. \hfill \emph{[1,5 pt]}
\item Le glaucome est fréquent, indolore et silencieux : le malade ne consulte que lorsque le champ visuel est déjà très réduit, comme M. Essomba. Un dépistage régulier de la PIO après 40 ans permet un diagnostic précoce et un traitement avant les lésions irréversibles. Chez les conducteurs professionnels, une vision tubulaire non dépistée est un danger majeur pour la sécurité routière (mauvaise perception des obstacles latéraux, conduite de nuit difficile). \hfill \emph{[1,5 pt]}
\item Exemple : « GLAUCOME : le voleur silencieux de la vue. Il ne fait pas mal, mais il aveugle. Après 40 ans, faites mesurer votre tension oculaire chaque année : votre vue n'a pas de prix ! » \hfill \emph{[1 pt]}
\end{enumerate}
\end{enumerate}
\end{corrige}

\begin{attention}[Points de vigilance — Barème Exercice 10]
\begin{itemize}
\item Les trois arguments du diagnostic doivent être \emph{distincts} et issus du Document 2 (PIO, champ visuel, fond d'œil) : répéter le même argument sous deux formes ne rapporte qu'un point.
\item Bien distinguer le \emph{but} du traitement (baisser la PIO) des \emph{moyens} (collyres, chirurgie) : les deux sont attendus.
\item Le message de campagne doit mentionner le caractère \emph{silencieux} (indolore) de la maladie et l'appel au \emph{dépistage}.
\end{itemize}
\end{attention}

\begin{corrige}[Exercice 11 — Sujet type BEPC : Synthèse — L'œil, de la lumière à l'image perçue (8 points)]
\begin{enumerate}
\item La lumière issue de l'objet traverse successivement la \textbf{cornée}, l'\textbf{humeur aqueuse}, la \textbf{pupille} (orifice de l'iris qui dose la lumière entrante), le \textbf{cristallin}, puis le \textbf{corps vitré}. Les rayons convergent et forment une image \textbf{réelle et inversée sur la rétine}. Les \textbf{photorécepteurs} (cônes et bâtonnets) de la rétine transforment l'énergie lumineuse en message nerveux, acheminé par le \textbf{nerf optique} jusqu'au cerveau, qui redresse et interprète l'image. \hfill \emph{[2,5 pts : ordre exact des milieux (1,5 pt) + image réelle et inversée sur la rétine (0,5 pt) + rôle des photorécepteurs (0,5 pt)]}
\item \begin{enumerate}
\item Pour la vision de près, les \textbf{muscles ciliaires se contractent}, ce qui \textbf{relâche le ligament suspenseur} ; libéré de la traction, le cristallin élastique \textbf{se bombe} : sa courbure et sa puissance de convergence augmentent, et l'image de l'objet proche se forme sur la rétine. \hfill \emph{[1 pt]}
\item Avec l'âge, le cristallin \textbf{perd progressivement son élasticité} (il durcit) : même lorsque les muscles ciliaires se contractent, il ne peut plus se bomber suffisamment. Le point proche recule au-delà de \SI{30}{cm} : c'est la presbytie. \hfill \emph{[1 pt]}
\item On prescrit des verres \textbf{convergents} pour la lecture, qui apportent la convergence que le cristallin ne fournit plus. Ce verre ne convient pas pour la vision de loin, car il rendrait floue l'image des objets éloignés (convergence excessive) ; le presbyte emmétrope n'en a besoin que pour la vision de près (verres de lecture ou verres à double foyer). \hfill \emph{[0,5 pt]}
\end{enumerate}
\item Tableau complété :
\begin{center}
\begin{tabularx}{\linewidth}{|l|X|X|X|}\hline
\rowcolor{popblueL}
\textbf{Pathologie} & \textbf{Structure atteinte / Mécanisme principal} & \textbf{Symptôme fonctionnel majeur} & \textbf{Traitement principal} \\ \hline
Cataracte & Cristallin opacifié (perte de transparence) & Vision floue, voile, couleurs ternes & Chirurgie : remplacement du cristallin par un implant \\ \hline
Glaucome chronique & Nerf optique détruit par une pression intra-oculaire trop élevée (mauvais drainage de l'humeur aqueuse) & Rétrécissement progressif du champ visuel (vision en tunnel), maladie indolore & Collyres hypotenseurs, parfois chirurgie ou laser \\ \hline
Myopie & Œil trop long ou trop convergent : image formée en avant de la rétine & Vision floue de loin, nette de près & Verres divergents (lunettes ou lentilles) \\ \hline
\end{tabularx}
\end{center}
\hfill \emph{[3 pts : 1 pt par ligne complète et exacte]}
\end{enumerate}
\end{corrige}

\begin{attention}[Points de vigilance — Barème Exercice 11]
\begin{itemize}
\item Dans le trajet de la lumière, l'\emph{ordre} des structures est noté : une structure oubliée ou mal placée fait perdre des points.
\item Pour la presbytie, le mot-clé attendu est la \emph{perte d'élasticité du cristallin} (et non « les muscles ciliaires ne fonctionnent plus », ce qui est faux).
\item Dans le tableau, chaque case doit être renseignée avec le vocabulaire du cours ; pour le glaucome, citer à la fois la PIO élevée et le nerf optique.
\end{itemize}
\end{attention}

\newpage

\coursec{Fiche bilan}

\begin{popfigure}
\begin{tikzpicture}[mindmap, grow cyclic, every node/.style={concept, text=white, font=\small\bfseries},
root concept/.append style={concept color=popblue, font=\large\bfseries},
level 1/.append style={level distance=3.6cm, sibling angle=60},
level 1 concept/.append style={font=\small\bfseries}]
\node[root concept]{L'{\oe}il et la vision}
  child[concept color=poporange]{ node{Anatomie de l'{\oe}il : cornée, iris, cristallin, rétine, nerf optique} }
  child[concept color=popgreen]{ node{Formation de l'image : image renversée sur la rétine} }
  child[concept color=poppurple]{ node{Accommodation : cristallin qui se bombe} }
  child[concept color=popgold]{ node{Anomalies : myopie, hypermétropie, presbytie} }
  child[concept color=poppink]{ node{Maladies : cataracte, glaucome} }
  child[concept color=popdark]{ node{Hygiène de la vision} };
\end{tikzpicture}
\legende{Carte mentale de la leçon : l'{\oe}il et la vision.}
\end{popfigure}

\begin{bilan}
\textbf{1. Définitions clés}
\begin{itemize}
\item \cle{Vision} : fonction qui permet de percevoir les objets grâce à la lumière qu'ils renvoient ; l'organe de la vision est l'\cle{{\oe}il}.
\item \cle{Accommodation} : modification de la courbure du cristallin (il se bombe) qui permet de voir net un objet proche.
\item \cle{Myopie} : anomalie où l'image se forme \textbf{en avant} de la rétine ({\oe}il trop long) ; on voit mal de loin ; correction par une lentille \textbf{divergente}.
\item \cle{Hypermétropie} : anomalie où l'image se forme \textbf{en arrière} de la rétine ({\oe}il trop court) ; on voit mal de près ; correction par une lentille \textbf{convergente}.
\item \cle{Presbytie} : baisse de l'accommodation avec l'âge (cristallin moins souple) ; correction par une lentille convergente.
\item \cle{Cataracte} : opacification du cristallin qui rend la vision trouble ; traitement chirurgical.
\item \cle{Glaucome} : maladie due à une pression trop élevée dans l'{\oe}il, qui détruit le nerf optique et peut conduire à la cécité.
\end{itemize}

\textbf{2. Anatomie de l'{\oe}il à connaître par c{\oe}ur}
\begin{center}
\begin{tabularx}{\linewidth}{|l|X|}\hline
\rowcolor{popblueL} \textbf{Partie} & \textbf{Rôle} \\ \hline
Cornée & Membrane transparente qui laisse entrer la lumière \\ \hline
Iris et pupille & L'iris colore l'{\oe}il ; la pupille règle la quantité de lumière qui entre \\ \hline
Cristallin & Lentille convergente qui dirige les rayons lumineux vers la rétine \\ \hline
Rétine & Membrane sensible à la lumière (cellules : cônes et bâtonnets) où se forme l'image \\ \hline
Nerf optique & Conduit le message nerveux de la rétine vers le cerveau \\ \hline
\end{tabularx}
\end{center}

\textbf{3. Méthodes express}
\begin{itemize}
\item Trajet de la lumière : objet $\rightarrow$ cornée $\rightarrow$ pupille $\rightarrow$ cristallin $\rightarrow$ rétine (image \textbf{renversée}) $\rightarrow$ nerf optique $\rightarrow$ cerveau (image redressée).
\item Pour diagnostiquer une anomalie : repérer où se forme l'image par rapport à la rétine (en avant = myopie ; en arrière = hypermétropie).
\item Pour choisir la correction : myopie $\rightarrow$ verre divergent ; hypermétropie et presbytie $\rightarrow$ verre convergent.
\item Hygiène : bon éclairage, distance de lecture, repos des yeux, dépistage régulier chez l'ophtalmologue.
\end{itemize}
\end{bilan}

\begin{aretenir}[Checklist avant le BEPC]
\begin{itemize}
\item $\square$ Je sais légender un schéma de l'{\oe}il (cornée, iris, pupille, cristallin, rétine, nerf optique).
\item $\square$ Je connais le rôle de chaque partie de l'{\oe}il.
\item $\square$ Je sais décrire le trajet de la lumière de l'objet jusqu'au cerveau.
\item $\square$ Je sais que l'image se forme renversée sur la rétine.
\item $\square$ Je sais définir l'accommodation et le rôle du cristallin.
\item $\square$ Je distingue myopie, hypermétropie et presbytie (cause, symptôme, correction).
\item $\square$ Je sais que la myopie se corrige par une lentille divergente.
\item $\square$ Je sais que l'hypermétropie et la presbytie se corrigent par une lentille convergente.
\item $\square$ Je connais la cause et les conséquences de la cataracte et du glaucome.
\item $\square$ Je sais citer au moins quatre règles d'hygiène de la vision.
\item $\square$ Je sais rédiger une conclusion justifiée à partir d'un schéma ou d'un document.
\end{itemize}
\end{aretenir}

\findecours

\end{document}
